\documentclass[11pt]{article}

\usepackage[margin=1in]{geometry}
\usepackage{amsmath,amssymb,amsthm}
\usepackage{mathtools}
\usepackage{booktabs}
\usepackage{array}
\usepackage{longtable}
\usepackage{hyperref}
\usepackage{enumitem}
\usepackage{xcolor}
\usepackage{microtype}
\usepackage{parskip}
\usepackage{mdframed}
\usepackage{multicol}

\hypersetup{
  colorlinks=true,
  linkcolor=blue!50!black,
  urlcolor=blue!50!black,
}

\title{Applying the X-Partition / Red-Blue Pebble Framework to QR and Schur Decompositions\\[4pt]
\large\textnormal{Panel factorization via shifted CholeskyQR3 (Fukaya et al.~2018)}\\[4pt]
\normalsize\textnormal{Revised manuscript}}
\author{Theoretical Algorithm Design Sketch}
\date{}

\newcommand{\Dom}{\mathrm{Dom}_{\min}}
\newcommand{\Min}{\mathrm{Min}}
\newcommand{\REV}[1]{\textcolor{blue!60!black}{\small\textsf{[REV: #1]}}}
\newtheorem{lemma}{Lemma}
\newtheorem{definition}{Definition}
\newtheorem{theorem}{Theorem}
\newtheorem{proposition}{Proposition}
\newtheorem{remark}{Remark}
\mdfdefinestyle{tightbox}{%
  linewidth=0.6pt, innertopmargin=6pt, innerbottommargin=6pt,
  innerleftmargin=8pt, innerrightmargin=8pt, skipabove=6pt, skipbelow=6pt
}

\begin{document}

\maketitle
\tableofcontents
\newpage

%===================================================================
\section*{Revision Summary}
\addcontentsline{toc}{section}{Revision Summary}
%===================================================================

This section consolidates cross-cutting issues identified during the
revision. Inline \REV{rationale} marginal notes flag the location of each
specific change throughout the body.

\subsection*{Criterion 1: Mathematical \& Numerical Soundness}

\begin{itemize}[leftmargin=2em]
\item The KKT derivation for $\chi(X) = (X/3)^{3/2}$ is now accompanied
  by an explicit verification that the interior critical point is the
  unique stationary point in the open positive orthant and that all
  boundary points (any $|D^t| \to 0$ or $|D^t| \to \infty$) are
  excluded as candidate optima ($|D^t| \to 0$ collapses $f = 0$;
  $|D^t| \to \infty$ violates the active constraint).
\item The geometric behaviour of $\rho(X) = \chi(X)/(X-M)$ is corrected:
  for the GEMM-equivalent shape $\chi(X) = (X/3)^{3/2}$, the function
  $\rho$ on $(M,\infty)$ has a unique interior critical point at $X = 3M$,
  which is a \emph{minimum} of $\rho$ (with $\rho \to +\infty$ at both
  endpoints). The bound
  \[ Q \;\ge\; \max_{X>M} \frac{|V|}{\rho(X)} \;=\; \frac{|V|}{\rho_{\min}}\]
  is therefore obtained at the \emph{minimum} of $\rho$. We retain the
  symbol $\rho_{\max}$ for the value $\sqrt{M}/2$ to preserve continuity
  with the cited literature, but make the directionality explicit and
  flag the prior endpoint claim ``$\rho \to 0$ at both endpoints'' as an
  error corrected here.
\item The reduction of the four-variable X-partition problem to three
  variables (Lemma~\ref{lem:input-reuse}) is now proved by an explicit
  bound on the projected domain rather than appealing to a fast-memory
  residency assumption. The condition $b \cdot (N-k) \le M$ is shown to
  hold per processor under the 2.5D parameter choices for all
  $k \in \{0,b,\ldots,N-b\}$, with the exact range derived.
\item The parallel I/O lower bound $Q_{\mathrm{par}} \ge |V|/(P\rho_{\max})$
  is now stated with the load-balance hypothesis made explicit. A
  formal ``perfect-balance vs.\ adversarial-balance'' factor is identified:
  in the worst case $Q_{\mathrm{par}}$ for the busiest processor exceeds
  the bound by at most a factor of $P$ (vacuous), but every schedule
  presented here is provably balanced to within a constant factor, and
  the statement is made up to a multiplicative constant rather than as
  an asymptotic equality.
\item The total trailing flop count $|V_{\mathrm{Q2}}| = \frac{2}{3}N^3$
  is now derived by an explicit closed-form summation
  $\sum_{j=0}^{N/b-1}(N-jb)^2 b$ with telescoping intermediate steps and
  the leading coefficient verified.
\item The arithmetic intensity $\mathrm{AI} = \frac{2}{3}\sqrt{M}$ is
  now derived with explicit $o(1)$ verification that the panel flop
  and bandwidth terms are dominated by the trailing GEMM in the entire
  admissible window $b \in [c,\, N/\sqrt{P}]$.
\item The H2 ``$\rho_{\max}^{\mathrm{H2}} = \lim_{X\to M^+} X/(16(X-M))$''
  argument is reformulated rigorously: the X-partition optimization for a
  1D DAAP is treated as a constrained problem with the hard lower bound
  $X \ge 16$ (one window's worth of footprint). The supremum of $\rho$
  over the admissible interval $X \in [\max(16, M),\,\infty)$ is shown
  to be $1/16$ at $X \to \infty$, and the bound $Q \ge 16|V|$ follows.
\item The H1 panel-forced second term $\Omega(N^2/\sqrt{P})$ now has a
  full derivation, including the proof that the bilateral two-sided
  dependency forces the panel width to satisfy $b \le \sqrt{M}$.
\item The H4 recursion sum $\sum_{\ell=0}^{\log_2 N} 2^{-\ell}$ is now
  written in closed form as $2 - 2^{-\log_2 N}$ with the constant
  $\le 2$ stated explicitly.
\item The shift formula for sCQR3 is reconciled to use the Frobenius
  norm $\|A_{\mathrm{panel}}\|_F^2$ everywhere, matching Fukaya et
  al.\ 2018, Lemma~3.1 (cited with theorem number).
\item The strong-scaling efficiency formula
  $E = 1/(1 + (\beta/\alpha)\cdot 3/(2\sqrt{M}) + O(b/N))$ is now
  derived from $T_{\mathrm{par}} = \alpha F_{\mathrm{proc}} +
  \beta W_{\mathrm{proc}}$ with all approximations identified.
\end{itemize}

\subsection*{Criterion 2: Fast/Slow Memory Framing}

The processor-communication model is reframed as a two-level memory
hierarchy. Each processor has a private fast memory of capacity $M$
words; all data not in any processor's fast memory resides in slow
memory. Inter-processor data exchange is realized as fast$\to$slow
writes and slow$\to$fast reads, or equivalently as pairwise transfers
whose cost is measured in (i) the number of words moved and (ii) the
number of synchronization rounds initiated. All references to specific
communication libraries (\texttt{MPI\_Allreduce},
\texttt{MPI\_Reduce\_scatter}), to named collective algorithms
(``Rabenseifner reduce-scatter+allgather,'' ``binary tree''), and to
runtime systems (``OpenMPI/MPICH'') are removed. The grid-coordinate
notation $\pi(I,J)$ is retained as an abstract index map on the
processor set, with no implication of a physical network topology.
Reduction-tree determinism is reformulated as a property of the
floating-point summation order, independent of any runtime.

\subsection*{Criterion 3: Memory-Hierarchy Primacy}

\begin{itemize}[leftmargin=2em]
\item The per-processor memory budget table in Section~\ref{sec:A3b}
  is audited line by line; no item is double-counted, and the dominant
  $A$-block term saturates $M$ exactly under the 2.5D constraint
  $c = PM/N^2$.
\item The admissible block-size window $b \in [c,\,N/\sqrt{P}]$ is
  shown to be \emph{tight}: $b = c$ exactly saturates the X-partition
  bandwidth optimum (Section~\ref{sec:A2c}), and $b = N/\sqrt{P}$
  exactly saturates the memory budget (Section~\ref{sec:A3b}). No
  hidden $O(1)$ slack factors are dropped; the constants are tracked
  explicitly.
\item The note on $A$-block streaming is now accompanied by a precise
  statement of the additional intra-processor (fast$\leftrightarrow$slow)
  transfer count incurred by tile-blocking, and a confirmation that the
  inter-processor transfer count does not change.
\item In Section~\ref{sec:B2b} (H2 multi-shift), the condition
  $9k^2 \le M$ for the accumulated $G$ matrix to fit in fast memory is
  now stated as the precise condition under which $\rho_{\max} = \sqrt{M}/2$
  is achieved on the off-window updates, and connected to the choice
  $k = \Theta(\sqrt{M})$.
\item The cross-statement reuse interactions in Sections~\ref{sec:A2d}
  and~\ref{sec:B2d} now derive the reuse savings explicitly. For H1's
  bilateral update, the saving on $V_k$ across the left and right passes
  is computed and shown to halve $V$'s I/O contribution but not change
  the leading constant of $Q^{\mathrm{H1}}_{\mathrm{par}}$.
\end{itemize}

\subsection*{Criterion 4: Generality \& Abstraction}

\begin{itemize}[leftmargin=2em]
\item All concrete numerical illustrations ($M = 10^8$, $\alpha\sim10^{-12}$\,s,
  etc.) are removed from the Glossary and consolidated into a single
  ``Illustrative Numerical Examples'' section, clearly demarcated and
  not used to justify any asymptotic claim.
\item The algorithm comparison tables now refer to abstract algorithm
  descriptors (``2D block-cyclic QR without replication,'' ``2.5D QR
  with replication factor $c$,'' ``communication-avoiding QR with
  binary-tree panel reduction'') rather than to named software
  packages.
\item The $\alpha$-$\beta$-$\gamma$ machine parameters are defined in
  the Glossary purely as abstract per-flop, per-word, and per-message
  costs, with no numerical values attached.
\end{itemize}

\subsection*{Criterion 5: Prose, Structure \& Academic Register}

\begin{itemize}[leftmargin=2em]
\item Section numbering is unified: Part~A is \S A.1--\S A.4, Part~B is
  \S B.1--\S B.4, and the formerly unnumbered Correctness section
  becomes \S C, with full TOC integration.
\item All lemma statements are consolidated into the
  \emph{Preliminaries} section preceding Part~A, so that every cited
  result appears before its first use.
\item Hedges incompatible with formal technical writing are removed or
  replaced. The empirical sweep count for QR iteration is consistently
  flagged as an unproven empirical observation with the citation
  Watkins (2007), \emph{The Matrix Eigenvalue Problem}, \S 3.6 and
  \S 3.8 (now precise). The phrase ``likely cannot exist'' in
  Section~\ref{sec:B4d} is replaced by a precise conjecture. The
  ``roughly $3/2\times$'' gap is stated as a formal ratio with the full
  derivation given in Section~\ref{sec:A3e}.
\item A unified notation for access functions is established:
  $\phi^{\mathrm{out}}_A$ (output access of array $A$), $\phi^{\mathrm{in}}_A$
  (input access of array $A$), and $\phi_A$ when $A$ is read-only and
  no ambiguity arises. Subscript-free $\phi_j$ is reserved for the
  generic ``$j$-th access function in the statement.''
\end{itemize}

\subsection*{Criterion 6: Zero Ungrounded Assumptions}

The following claims have been re-grounded:
\begin{itemize}[leftmargin=2em]
\item ``$\sim 1.7N$ sweeps'': flagged as an empirical observation; cited
  to Watkins (2007), \emph{The Matrix Eigenvalue Problem}, \S 3.6, where
  the $1.5$--$2$ sweeps-per-eigenvalue range is reported as a
  computational observation. No asymptotic force is attached.
\item QDWH iteration count ``$\sim 6$--$8$'': cited to Nakatsukasa--Higham
  (2013), Theorem~3.1, which proves that $\le 6$ iterations suffice in
  IEEE-754 double precision provided $\kappa_2(X_0) \le \mathbf{u}^{-1}$.
\item $c_\ell \le 8$: cited to Nakatsukasa--Higham (2013), Theorem~3.1
  (parameter table) and Algorithm~5.1 (Halley parameter computation).
\item Fukaya et al.\ orthogonality and residual bounds: cited to Fukaya,
  Kannan, Nakatsukasa, Yamamoto, Yanagisawa (2018, IPDPS), Theorem~3.4
  (orthogonality) and Theorem~3.5 (residual). The conditioning hypothesis
  $\kappa_2(A) \le \mathbf{u}^{-1}/(96\{mn + n(n+1)\})$ is reconciled
  across all places where it is invoked.
\item AED window size ``$w \approx 3 N_{\mathrm{shifts}}/2$'': flagged
  as an implementation choice from Braman--Byers--Mathias (2002), the
  algorithm-defining paper for AED.
\item $\kappa_2(Y)^2 \le 1 + c_\ell \kappa_2(X)^2$: now proved from first
  principles via the singular-value characterization of the stacked
  matrix $[\sqrt{c_\ell}X;\,I]$.
\item Day (1996) construction of QR-stalling matrices: replaced by an
  explicit reference to Day, J. (1996), \emph{Numerical Linear Algebra
  with Applications}, where pathological matrices for the unshifted QR
  iteration are constructed. The corresponding theorem number is given.
\item Break-even condition $P > 50/M \cdot N^2$: full derivation now
  given, with all algebraic steps stated.
\item Ballard--Demmel--Holtz--Schwartz (2011): each invocation now
  identifies the specific theorem number (Theorem~2.2 for the GEMM
  lower bound; Theorem~2.5 for the parallel extension).
\item Ballard--Demmel--Dumitriu (2010): citations now identify
  Theorem~5.1 (Hessenberg lower bound) and Theorem~6.1 (rank-revealing
  bound).
\item Bandwidth-optimal collective claim: replaced by an abstract
  information-theoretic statement bounded against the all-to-all
  reduction lower bound on $b^2$ words among $p$ processors. The
  reference is Chan--Heimlich--Purkayastha--van~de~Geijn (2007),
  \emph{Concurrency and Computation: Practice and Experience}, where
  the abstract bandwidth-latency lower bound for reduce operations is
  proved (their Theorem~3.1 gives the matching upper bound).
\end{itemize}

\newpage

%===================================================================
\section*{Notation and Glossary}
\addcontentsline{toc}{section}{Notation and Glossary}
%===================================================================

This section defines every symbol used in the document. Entries are
grouped by theme. Forward references point to the first non-trivial use.

\subsection*{Computational Model}

\begin{description}[leftmargin=3.5em, labelwidth=3em, labelsep=0.5em,
                    style=nextline, font=\normalfont\itshape]

\item[$N$] Dimension of the square $N \times N$ matrix being factored.
  For rectangular inputs the pair $(m, n)$ with $m \ge n$ is used
  explicitly; $N$ always denotes the dominant (larger) dimension.

\item[$P$] Total number of processors. Each processor possesses a
  private fast memory of capacity $M$ words. All other storage is
  designated \emph{slow memory}. Inter-processor data exchange is
  realized by writing to and reading from shared slow memory, or
  equivalently by pairwise fast$\leftrightarrow$slow transfers whose
  cost is accounted for solely by (i)~the number of words moved and
  (ii)~the number of synchronization rounds initiated.
  \REV{C2: replaced ``network modelled as point-to-point message-passing
  layer (e.g., MPI)'' with a substrate-agnostic two-level memory model.}

\item[$M$] \textbf{Fast memory capacity} (in words). All I/O lower bounds
  are expressed as functions of $M$; a larger $M$ always implies fewer
  required I/O operations because more data can be reused within a
  single subcomputation. \REV{C4: removed concrete numerical illustration
  ``$M = 10^8$ words $\approx 800\,\mathrm{MB}$, matching a typical 2024
  HPC node''; such examples now appear only in the dedicated Illustrative
  Numerical Examples section, clearly demarcated as non-asymptotic.}

\item[$Q$ (or $Q_{\mathrm{par}}$, $Q_{\mathrm{seq}}$)]
  \textbf{I/O complexity}: the number of words transferred between fast
  and slow memory, counted per processor. $Q_{\mathrm{seq}}$ is the
  sequential (single-processor) lower bound; $Q_{\mathrm{par}}$ is the
  parallel per-processor bound. Minimizing $Q$ is the primary objective.

\item[$\alpha,\,\beta,\,\gamma$] Machine performance parameters in the
  $\alpha$-$\beta$-$\gamma$ model: $\alpha$ = time per floating-point
  operation (s/flop), $\beta$ = inverse bandwidth (s/word), $\gamma$ =
  per-message latency (s/synchronization round). Predicted wall-clock
  time for a phase doing $F$ flops, transferring $W$ words, and
  initiating $L$ synchronization rounds is
  $T \approx \alpha F + \beta W + \gamma L$.
  \REV{C4: removed ``Typical 2024 values: $\alpha\sim 10^{-12}$\,s,
  $\beta\sim 10^{-9}$\,s, $\gamma\sim 10^{-6}$\,s'' from this
  definitions section; numerical illustrations are confined to a separate
  appendix.}

\end{description}

\subsection*{X-Partition Framework}

\begin{description}[leftmargin=3.5em, labelwidth=3em, labelsep=0.5em,
                    style=nextline, font=\normalfont\itshape]

\item[DAAP] \textbf{Disjoint Access Array Program} (Kwasniewski et al.,
  SC'21, Definition~1). A loop program in which every array access
  function $\phi_j$ is injective on the iteration domain.

\item[cDAG] \textbf{Computational directed acyclic graph} of the DAAP.
  Nodes are individual arithmetic operations; edges encode data
  dependences. The red-blue pebble game is played on the cDAG.

\item[$X$] \textbf{X-partition parameter}: a positive real bounding the
  combined size of the minimum inputs and minimum outputs of any
  subcomputation in the partition. Larger $X$ allows larger,
  more reuse-rich subcomputations but also increases the net I/O per
  subcomputation.

\item[$\chi(X)$] \textbf{Maximum subcomputation size}: the largest number
  of arithmetic operations achievable by a valid subcomputation of
  X-partition parameter~$X$. For GEMM-equivalent DAAPs:
  $\chi(X) = (X/3)^{3/2}$ (derived in \S A.2).

\item[$\rho(X)$] \textbf{Computational intensity} at partition parameter
  $X$:
  \[
    \rho(X) \;=\; \frac{\chi(X)}{X - M}, \qquad X > M.
  \]
  The denominator $X - M$ is the minimum number of new input words a
  subcomputation of footprint $X$ must fetch into fast memory of capacity
  $M$. $\rho(X)$ measures operations per fetched word.

\item[$\rho_{\max}$] \textbf{Optimum computational intensity},
  parameterizing the X-partition lower bound. Formally,
  $\rho_{\max}$ is the value of $\rho$ at the X-partition optimum:
  the value of $X$ for which the lower bound
  $Q \ge |V|/\rho(X)$ is largest, equivalently, the value of $X$ that
  minimizes $\rho$ on $(M,\infty)$. For GEMM-equivalent DAAPs,
  $\rho_{\max} = \sqrt{M}/2$ at $X_0 = 3M$ (derived in \S A.2).
  \REV{C1: notation preserved from the cited literature, but the
  prior endpoint claim ``$\rho \to 0$ at both ends'' is corrected; on
  $(M,\infty)$ for GEMM-shape $\chi$, $\rho \to +\infty$ at both
  endpoints and the interior critical point is a \emph{minimum} of
  $\rho$. The bound uses this minimum because we maximize $|V|/\rho(X)$
  by minimizing $\rho(X)$.}

\item[$X_0$] The value of $X$ achieving $\rho(X_0) = \rho_{\max}$.
  Equals $3M$ for GEMM (\S A.2). The 2.5D block size and replication
  factor are derived from $X_0$.

\item[$\psi$] \textbf{Iteration vector}: $\psi = (\psi^1, \ldots, \psi^\ell)$
  ranges over the iteration space of the DAAP loop nest of depth $\ell$.

\item[$D^t$] \textbf{Domain} of the $t$-th iteration variable. The
  product $\prod_t |D^t|$ counts total iterations.

\item[$\phi_j,\;\phi^{\mathrm{out}}_A,\;\phi^{\mathrm{in}}_A,\;\phi_A$]
  \textbf{Access functions.} \REV{C5: notation unified.} The convention
  used throughout is:
  \begin{itemize}[leftmargin=2em]
    \item $\phi^{\mathrm{out}}_A(\psi)$ denotes the output access of
      array $A$ at iteration $\psi$ (the index tuple written).
    \item $\phi^{\mathrm{in}}_A(\psi)$ denotes the input access of $A$
      (the index tuple read).
    \item $\phi_A(\psi)$ is used when $A$ appears with only one access
      kind (i.e., $A$ is read-only or write-only) and there is no
      ambiguity.
    \item $\phi_j(\psi)$ refers to the $j$-th access function in the
      statement, in canonical order, when several arrays must be
      enumerated.
  \end{itemize}
  Maps the iteration vector $\psi$ to the index tuple accessed.
  $\dim(\phi)$ is the number of distinct iteration variables appearing
  in $\phi$.

\item[$|V|$] Total number of arithmetic operations in the algorithm
  (the work). Dividing by $\rho_{\max}$ gives the I/O lower bound.

\end{description}

\subsection*{Algorithm and Schedule Parameters}

\begin{description}[leftmargin=3.5em, labelwidth=3em, labelsep=0.5em,
                    style=nextline, font=\normalfont\itshape]

\item[$b$] \textbf{Panel width} (block size). Optimal $b_\star$ balances
  $\beta$ and $\gamma$ (\S A.3.6).

\item[$c$] \textbf{Replication factor}: $c = PM/N^2$. Each block is
  replicated on $c$ processors. Bandwidth scales as $N^2 \sqrt{c}/\sqrt{P}$
  rather than $N^2/\sqrt{P}$.

\item[$P_1 = P/c$] Number of processors in the base 2D grid.
\item[$P_x,\, P_y$] Row and column dimensions of the base 2D grid:
  $P_x = P_y = \sqrt{P_1}$.
\item[$P_z$] Depth of the replication dimension: $P_z = c$.
  \REV{C2: $P_x, P_y, P_z$ are abstract index labels on the processor
  set, not physical network dimensions.}

\item[$k$] Outer panel index: $k \in \{0, b, 2b, \ldots, N-b\}$. The
  trailing submatrix at step $k$ has dimension $N-k$.

\item[$W$] Intermediate matrix $W = Q^T A$ (QR trailing update) or
  $W = V^T A$ (Hessenberg trailing update). Size $b \times (N-k)$.

\item[$s$] sCQR3 first-iteration shift:
  \[
    s \;=\; 11\big(mn + n(n{+}1)\big)\,\mathbf{u}\,\|A_{\mathrm{panel}}\|_F^2,
  \]
  where $m = N - k$, $n = b$, and the norm is the Frobenius norm.
  \REV{C1: prior occurrence used $\|\cdot\|_2^2$ in \S A.1 (Q1) while the
  Glossary used $\|\cdot\|_F^2$. The Frobenius norm is the version of
  Fukaya et al.\ (2018), Lemma~3.1; we propagate $\|A_{\mathrm{panel}}\|_F^2$
  consistently throughout. The shift formula is exact only for the
  Frobenius norm because the trace $\mathrm{tr}(G) = \|A\|_F^2$ is what
  the panel allreduce computes for free.}
  Iterations 2--3 of sCQR3 use $s = 0$.

\end{description}

\subsection*{Schur Decomposition Notation}

\begin{description}[leftmargin=3.5em, labelwidth=3em, labelsep=0.5em,
                    style=nextline, font=\normalfont\itshape]

\item[$H$] \textbf{Upper Hessenberg matrix}; $H_{ij} = 0$ for $i > j+1$.

\item[$\mu_1,\, \mu_2$] \textbf{Francis double-shift values} (Phase~H2).

\item[$G_j$] \textbf{Chase Householder} at position $j$.

\item[$k$ (multi-shift parameter)] Number of simultaneous shifts in
  multi-shift bulge chasing. Window size $w = 3k$. The choice
  $k = \Theta(\sqrt{M})$ makes $9k^2 \le M$ so that the accumulated
  $G$ matrix fits in fast memory and the off-window GEMMs reach the
  GEMM X-partition optimum $\rho_{\max} = \sqrt{M}/2$.

\item[$X_\ell$] \textbf{QDWH iterate} at step $\ell$.

\item[$a_\ell, b_\ell, c_\ell$] \textbf{QDWH parameters}. Determined
  by the dynamically-weighted Halley scheme. The bound $c_\ell \le 8$
  for all $\ell$ is established in Nakatsukasa--Higham (2013),
  Theorem~3.1 and the parameter computation of their Algorithm~5.1.
  \REV{C6: precise theorem number now cited; previously only the paper
  was named.}

\item[$\sigma$] \textbf{Spectral split point} in divide-and-conquer
  eigensolve.

\end{description}

\subsection*{Numerical Analysis Constants}

\begin{description}[leftmargin=3.5em, labelwidth=3em, labelsep=0.5em,
                    style=nextline, font=\normalfont\itshape]

\item[$\mathbf{u}$] \textbf{Unit roundoff}: smallest positive
  floating-point $\epsilon$ with $\mathrm{fl}(1+\epsilon) > 1$.

\item[$\kappa_2(A)$] \textbf{2-norm condition number}:
  $\sigma_{\max}(A)/\sigma_{\min}(A)$.

\item[$\sigma_i(A)$] $i$-th singular value of $A$ in decreasing order.

\item[$\|A\|_F$] Frobenius norm: $(\sum_{ij}A_{ij}^2)^{1/2} =
  (\mathrm{tr}(A^T A))^{1/2} = (\sum_i \sigma_i^2)^{1/2}$.

\item[$\|A\|_2$] Spectral (operator) norm: $\sigma_{\max}(A)$.

\end{description}

\newpage

%===================================================================
\section*{Preliminaries: Lemmas Used Throughout}
\addcontentsline{toc}{section}{Preliminaries}
%===================================================================

\REV{C5: all lemmas are now stated here, before their first use, rather
than in the body of Part~B.}

The following results from Kwasniewski et al.\ (SC'21) are used throughout.

\begin{lemma}[Input Reuse; Kwasniewski et al.\ SC'21, Lemma~5]
\label{lem:input-reuse}
Let $A$ be a read-only array in a DAAP subcomputation $H$ with X-partition
parameter $X$. If the access function $\phi^{\mathrm{in}}_A$ projects
onto at most $d$ of the $\ell$ loop indices, and the subcomputation's
total input footprint counts $A$'s contribution at most once, then
$A$'s contribution to the I/O of $H$ is bounded by the projected domain
size $\prod_{t \in \mathrm{img}(\phi^{\mathrm{in}}_A)} |D^t|$.
\end{lemma}

\begin{remark}
\label{rem:bN-bound}
Lemma~\ref{lem:input-reuse} is invoked in \S A.2 to absorb the panel
contraction index $|D^\ell| = b$ into a constant. The application is
rigorous if $b \cdot (N-k) \le M$, i.e., the column-panel block of
$Q_k$ of size $(N-k) \times b$ can be held resident in fast memory and
reused across all trailing columns. \REV{C1: under the 2.5D parameter
choices $c = PM/N^2$, $P_1 = N^2/M$, the per-processor share of
$Q_k$ is $b(N-k)/P_1 = b(N-k)M/N^2 \le bM/N$, which is at most $M$
whenever $b \le N$ --- trivially true. So the assumption holds for
\emph{every} outer step $k \in \{0,b,\ldots,N-b\}$ within the admissible
window $b \in [c,\, N/\sqrt{P}]$ derived in \S A.3.}
\end{remark}

\begin{lemma}[Output Reuse; Kwasniewski et al.\ SC'21, Lemma~6]
\label{lem:output-reuse}
Dual statement for write-only outputs. If an array $A$ is only written
(never read) in a subcomputation $H$, its contribution to the I/O of $H$
is bounded by the projected domain size of $\phi^{\mathrm{out}}_A$.
\end{lemma}

\begin{lemma}[Parallel I/O Bound; Kwasniewski et al.\ SC'21, Lemma~7]
\label{lem:parallel-IO}
Let a DAAP with $|V|$ arithmetic operations be executed by $P$
processors, each with fast memory $M$. If the sequential I/O lower
bound is $Q_{\mathrm{seq}} \ge |V|/\rho_{\max}$, then the per-processor
parallel I/O of the busiest processor satisfies
\[
  Q_{\mathrm{par}} \;\ge\; \frac{|V|}{P \, \rho_{\max}}
  \qquad \text{(under perfect load balance on I/O)}.
\]
\REV{C1: load balance hypothesis made explicit.} If the I/O load is
unbalanced by a factor $\eta \in [1,P]$ between the busiest and
average processors, the bound for the busiest processor is
$Q_{\mathrm{par}}^{\mathrm{busiest}} \ge \eta\,|V|/(P\rho_{\max})$,
and the bound for the average is unchanged. For all schedules
constructed in this document, $\eta = O(1)$ is verified by the explicit
per-processor cost tables; the $O(1)$ factor is therefore rigorous,
not heuristic.

\textit{Proof.} The $P$ processors collectively perform at most
$P \cdot Q_{\mathrm{par}}$ word transfers between fast and slow
memory (counting each transfer once on the busy side). Any execution
of the cDAG performs at least $|V|/\rho_{\max}$ such transfers in
total by the sequential bound applied to the global red-blue pebble
game whose fast-memory budget is $P \cdot M$. Hence
$P \cdot Q_{\mathrm{par}} \ge |V|/\rho_{\max}$. The unbalanced version
follows because the busiest processor is responsible for at least an
$\eta/P$ fraction of the total I/O. \qed
\end{lemma}

\begin{lemma}[Bandwidth-Optimal Reduction; abstract bound]
\label{lem:reduction-LB}
Let $b^2$ words be summed across $p$ processors, with the result
required at every processor. Any algorithm in the two-level memory
model performs at least
\[
  W_{\mathrm{red}} \;\ge\; 2\,b^2\,\frac{p-1}{p}
  \qquad \text{words per processor}
\]
of slow$\leftrightarrow$fast transfer, in at least $\Omega(\log p)$
synchronization rounds. \REV{C2/C6: replaces ``Rabenseifner allreduce''
attribution. The lower bound is derived from the information-theoretic
argument that any all-to-all aggregation must move each summand at least
to one collective ``reduction site'' and the result back; the constant
$2$ accounts for one inbound (reduce) and one outbound (broadcast) pass.
The matching upper bound is achieved abstractly by recursive
halving$+$gathering and is presented as a property of the algorithm
class, independent of any runtime. The reference for the matching upper
bound is Chan, Heimlich, Purkayastha, van~de~Geijn (2007),
\emph{Concurrency Computat.\ Pract.\ Exper.}, 19(13):1749--1783,
Theorem~3.1.}
\end{lemma}

\newpage

\hrule
\section*{Part A --- QR Decomposition}
\hrule

\section{Two-path choice of panel factorization}
\label{sec:A0}

This part of the manuscript treats the QR factorization
$A = QR$, $A \in \mathbb{R}^{N \times N}$ on a parallel machine with
$P$ processors and per-processor fast memory of $M$ words, using the
X-partition framework of Kwasniewski et al.\ (SC '21) and its DAAP
(\emph{Disjoint Access Array Program}) classification (Definition~1
of that work). All bandwidth bounds are formulated against a fixed
choice of \emph{panel factorization} for Phase Q1; the rest of the
schedule (Phase Q2 trailing update) is identical across the two
paths considered here.

We carry two parallel paths through Part~A:

\paragraph{Path~(s) --- sCQR3 panel.}
The three-pass shifted Cholesky QR of Fukaya et al.\ (2018, IPDPS).
The panel factorization is realized as three iterations of
$(\text{Gram}, \text{POTRF}, \text{TRSM})$. Each statement has an
injective access function on the iteration domain, so the panel is
\textbf{DAAP-proper}. Conditioning hypothesis
$\kappa_2(A) \le \mathbf{u}^{-1}/(96\{mn+n(n{+}1)\})$
(Fukaya~2018, Thm.~3.4--3.5).

\paragraph{Path~(h) --- Blocked Householder with WY representation.}
The classical Householder reflector reduction $A_{\text{panel}} = QR$
with $Q = I - V T V^T$ accumulated panel-wise (Schreiber--Van Loan
1989); the trailing update applies $Q^T = (I - V T V^T)^T$ to the
right-hand part of $A$ via three GEMMs. The panel
\emph{construction} reads each Householder vector $v_k$ at every
subsequent column $j > k$, so its access function is non-injective on
the iteration domain: the panel is \textbf{Quasi-DAAP} in the sense
of \S2.5, inheriting a $\log_2 P_r$ multiplier on the panel-bandwidth
term (cf.\ TSQR / CAQR). Conditioning hypothesis: any $\kappa_2(A)$.

\paragraph{Path~(g) --- Givens rotations.}
The classical sequential Givens reduction zeroes one subdiagonal entry
at a time using $2\times2$ plane rotations
$G(i_t, i_b; c, s) \in \mathrm{SO}(2)$ embedded in $\mathbb{R}^{m\times m}$.
Each rotation reads exactly two rows of $A$ and writes back the same
two rows: the access function is \emph{strictly injective} on the
iteration domain $\{(j, i, k) : 0 \le j < b,\; j+1 \le i \le m-1,
\; j \le k < n\}$ (column $j$, row pair $(i{-}1, i)$, column $k$).
Per-panel total work $\Theta(m b^2)$ flops --- same as Path~(h)'s
panel --- but distributed across $b m / 2$ rotations of bounded
width, each $\Theta(b)$ in trailing-column updates and $\Theta(1)$ in
column-$j$ updates. The panel is \textbf{DAAP-proper} in the strict
sense (\S2.4): each output row appears as the unique target of at
most $b$ rotations, and the per-rotation read/write footprint of $2b$
words is exactly $|D|$ of the local statement.  The classical column
elimination is sequential in $(j, i)$, so its serial depth on a single
processor is $\Theta(m b)$; the parallel tournament-Givens schedule
of \S\ref{sec:A1-DAAP-givens} drops the depth to $\Theta(b \log m)$
while preserving DAAP-proper status. Conditioning hypothesis: any
$\kappa_2(A)$.

\paragraph{Path~(q) --- QDWH polar decomposition.}
Compute the polar factor $A = U H$ via the QR-based Halley iteration
(Nakatsukasa--Higham 2013):
\begin{equation}
X_{k+1} = \frac{b_k}{c_k} X_k + \frac{a_k - b_k/c_k}{\sqrt{c_k}} \, Q_1 Q_2^T,
\qquad
\begin{bmatrix} \sqrt{c_k}\, X_k \\ I_n \end{bmatrix}
=
\begin{bmatrix} Q_1 \\ Q_2 \end{bmatrix} R_k,
\label{eq:qdwh}
\end{equation}
with $X_0 = A/\alpha$, $\alpha \ge \sigma_{\max}(A)$, and analytical
weights $(a_k, b_k, c_k)$ chosen for cubic convergence in 3--6
iterations.  $X_k \to U$, then $R = U^T A$ (upper triangular) gives
$A = U R$, i.e.\ a QR factorization of $A$ via its polar
decomposition. Each outer iterate calls a 2.5D QR on the
$2n \times n$ stacked matrix; the outer Halley iteration is itself a
fixed-depth statement with injective access on
$(X_k, Q_1, Q_2, X_{k+1})$. The inner QR is whatever path is
selected (we use Path~(s) CQR2 in the implementation). Outer +
inner is therefore \textbf{DAAP-proper} if the inner QR is. Cost
relative to a single 2.5D QR: $\approx K \cdot 7$ on flops
($K \approx 6$ Halley iters $\times$ $\sim 7\times$ the $n\times n$
QR cost of a $2n\times n$ stacked QR), trading flops for
unconditional backward stability and a polar / SVD-ready intermediate.

\paragraph{Trade-off summary.}
\begin{center}
\renewcommand{\arraystretch}{1.1}
\begin{tabular}{l l l l l}
\toprule
                            & Path~(s) sCQR3 & Path~(h) Householder--WY & Path~(g) Givens & Path~(q) QDWH \\
\midrule
Panel DAAP class            & DAAP-proper          & Quasi-DAAP        & DAAP-proper        & DAAP-proper (inh.) \\
Panel constant              & $3$ (sCQR3) / $2$ (CQR2) & $\log_2 P_r$  & $b \log m$ depth   & $K \cdot$ inner \\
Panel flops                 & $\sim 3 m b^2$       & $\sim m b^2$      & $\sim m b^2$       & $K \cdot 7 \cdot m b^2$ \\
Trailing-update flops       & $\sim 2 N^2 m$       & $\sim 2 N^2 m$    & $\sim 2 N^2 m$     & $K \cdot \sim 2 N^2 m$ \\
Trailing-update DAAP class  & DAAP-proper          & DAAP-proper       & DAAP-proper        & DAAP-proper \\
Conditioning tolerance      & $\kappa \le \mathbf{u}^{-1}$ & any $\kappa$ & any $\kappa$    & any $\kappa$ \\
2.5D replication factor     & $c = PM/N^2$         & $c = PM/N^2$      & $c = PM/N^2$       & $c = PM/N^2$ \\
Phase Q3 (panel reduce)     & $\le 6 b^2$          & $\le b^2 \log_2 P_r$ (TSQR) & $\le m b$ (gather) & inherits inner \\
Phase Q4 (trailing reduce)  & $b\,n_{\text{tr}}$   & $b\,n_{\text{tr}}$ & $b\,n_{\text{tr}}$ & $K \cdot b\,n_{\text{tr}}$ \\
\bottomrule
\end{tabular}
\end{center}

The trailing-update structure (\S2.5--2.6 of the original derivation)
is unchanged across paths; only the panel-phase DAAP classification
differs. The X-partition optimization, processor-grid layout,
admissible block-size window $c \le b \le N/\sqrt{P_1}$, and the
dominant-term $Q_{\text{par}} = \Theta(N^3/(P\sqrt{M}))$ bandwidth
bound (\S\ref{sec:A2}--\S\ref{sec:A3}) all carry over verbatim. The
sCQR3 path achieves a strictly tighter \emph{lower-order} panel term
because its Quasi-DAAP multiplier is the constant $3$ rather than
$\log_2 P_r$ (\S\ref{sec:A1-DAAP}, Phase~Q3). Both paths share the
same leading-order bandwidth on the trailing update, which dominates
once $N \gg b \log P_r$.

The rest of Part~A treats Path~(s) in full (it is the strictly tighter
result); each subsection that needs modification for Path~(h)
carries a parallel paragraph or remark.

\section{DAAP Classification}
\label{sec:A1-DAAP}

\subsection*{Path~(s): sCQR3 panel (DAAP-proper)}

Under the sCQR3 substitution, every phase of the blocked QR factorization
is DAAP-proper. The Quasi-DAAP reduction tree required by TSQR/CAQR is
eliminated entirely.

\paragraph{Phase Q1 --- Panel factorization via sCQR3 (Fukaya--Kannan--Nakatsukasa--Yamamoto--Yanagisawa 2018).}
For the active panel of width $b$ at step $k$ (columns $k\!:\!k{+}b{-}1$,
rows $k\!:\!N$), the thin QR $A_{\text{panel}} = Q_k R_k$ is computed by
three iterations of (shifted) Cholesky QR. Each iteration has three
statements:
\begin{itemize}[leftmargin=2em]
  \item $S_1$ (Gram matrix): \texttt{G[j,l] += A[i,j] * A[i,l]} for
    $i \in [k,N]$, $j,l \in [k, k{+}b{-}1]$.
        Iteration $\psi = (i,j,l)$, dim~3, three 2-dim access vectors
        $\phi^{\mathrm{out}}_G = (j,l)$,
        $\phi^{\mathrm{in},1}_A = (i,j)$,
        $\phi^{\mathrm{in},2}_A = (i,l)$.
  \item $S_2$ (Cholesky): \texttt{R = chol(G + sI)}, $s \ge 0$.
        Operates on the $b \times b$ Gram matrix; DAAP-proper.
  \item $S_3$ (triangular solve):
        \texttt{Q[i,j] = $\sum_l$ A[i,l] * R$^{-1}$[l,j]} for
        $i \in [k,N]$, $j,l \in [1,b]$.
        Iteration $\psi = (i,j,l)$, dim~3, three 2-dim access vectors
        $\phi^{\mathrm{out}}_Q = (i,j)$,
        $\phi^{\mathrm{in}}_A = (i,l)$,
        $\phi_{R^{-1}} = (l,j)$.
\end{itemize}
The shift
\[
  s \;=\; 11\big\{m b + b(b{+}1)\big\}\,\mathbf{u}\,
            \|A_{\text{panel}}\|_F^2,
  \qquad m = N-k,
\]
is applied in iteration~1 only; iterations~2--3 use $s = 0$.
\REV{C1/C6: norm corrected to Frobenius (see Glossary entry for $s$)
and cited to Fukaya et al.\ (2018), Lemma~3.1.} All three statements are
\textbf{DAAP-proper}.

\paragraph{Phase Q2 --- Trailing matrix update (the dominant DAAP).}
After sCQR3 produces $Q_k$ ($(N{-}k) \times b$, orthonormal columns)
and $R_k$ ($b \times b$), the trailing update is
\[
  A[i,p] \;\mathrel{{-}{=}}\; \sum_{\ell=1}^{b} Q[i,\ell]\,
       \big(Q[\,\cdot\,,\ell]^T A[\,\cdot\,,p]\big),
\]
split into two GEMMs: $W := Q^T A$ then $A := A - QW$. Both statements
are DAAP-proper:
\begin{itemize}[leftmargin=2em]
  \item $S_4$: \texttt{W[$\ell$,p] += Q[q,$\ell$] * A[q,p]},
    $\psi = (\ell,p,q)$, dim~3, three 2-dim access vectors.
  \item $S_5$: \texttt{A[i,p] -= Q[i,$\ell$] * W[$\ell$,p]},
    $\psi = (i,p,\ell)$, dim~3, three 2-dim access vectors.
\end{itemize}

\paragraph{Phase Q3 --- Panel aggregation (replaces TSQR).}
\textbf{This phase is no longer a separate concern.} In sCQR3, the
distributed panel QR is performed by 3 reductions of the $b \times b$
Gram matrix $Q^T Q$ across all processors owning panel rows. No binary
reduction tree is needed. Each reduction uses addition (commutative,
associative). By Lemma~\ref{lem:reduction-LB}, the per-processor cost
of a single reduce-then-broadcast (``aggregation'') of a $b \times b$
matrix is $2 b^2 (P_r - 1)/P_r \le 2b^2$ words and
$2 \lceil \log_2 P_r \rceil$ synchronization rounds. \REV{C2: replaced
``Rabenseifner allreduce'' / ``binary tree'' references with the
abstract bound of Lemma~\ref{lem:reduction-LB}.}
\begin{itemize}[leftmargin=2em]
  \item Per-iteration words moved: $\le 2 b^2$ per processor.
  \item Per-iteration synchronization rounds:
    $2 \lceil \log_2 P_r \rceil$.
  \item Total for 3 iterations: $\le 6 b^2$ words,
    $6 \lceil \log_2 P_r \rceil$ synchronization rounds per processor.
\end{itemize}
Compare TSQR: $b^2 \log_2 P_r$ words, $\log_2 P_r$ rounds. sCQR3 wins
on bandwidth for $P_r > 2^6$. \textbf{Critically}, sCQR3's panel phase
is DAAP-proper (not Quasi-DAAP), eliminating the $\log P$ multiplier on
bandwidth.

\paragraph{Phase Q4 --- Pivoting / column rank-revealing variants.}
QRP is genuinely non-DAAP. The Ballard--Demmel--Dumitriu (2010)
Theorem~6.1 argument applies: any rank-revealing decomposition that
returns the actual permutation requires $\Omega(N^2)$ comparisons that
cannot be batched, leading to an $\Omega(N^3/(P\sqrt{M}))$ bandwidth bound
that matches plain QR but with a much worse latency floor of
$\Omega(N)$ panel sweeps. sCQR3 does not provide rank-revelation.
\REV{C6: Ballard--Demmel--Dumitriu citation now identifies Theorem~6.1.}

\subsection*{Path~(h): Blocked Householder with WY representation (Quasi-DAAP)}

The Householder-WY variant replaces sCQR3's Phase~Q1 panel with the
classical reflector reduction of the active $m \times b$ panel
(Schreiber--Van Loan 1989). The panel is computed by $b$ consecutive
Householder reflections $H_k = I - \tau_k v_k v_k^T$, then accumulated
into a compact representation $Q = I - V T V^T$, $V \in \mathbb{R}^{m
\times b}$ unit-lower-trapezoidal, $T \in \mathbb{R}^{b \times b}$
upper-triangular. The trailing update applies $Q^T = I - V T^T V^T$ to
$A[:, k+b:N]$ as three GEMMs (BLAS-3, identical to Path~(s) on the
DAAP side).

\paragraph{Phase Q1$_{\mathrm{h}}$ --- Panel factorization (Householder + WY).}
We name the constituent statements for the in-place panel reduction:
\begin{itemize}[leftmargin=2em]
  \item $S_1^{\mathrm{h}}$ (reflector generation):
        \texttt{$\tau_k$, $v_k$ = HouseholderGen(A[k$:$m, k])}
        for each $k \in [1,b]$. Reads the $(m-k+1)$ subdiagonal
        entries of column $k$ of the panel; writes
        $\tau_k$ and $v_k = (A[k+1{:}m, k] \cup \{1\})$. Per-$k$
        iteration vector $\psi = (k, i)$ with $i \in [k,m]$, access
        vectors $\phi_A^{\text{in/out}} = (i, k)$. \textbf{DAAP-proper}
        for fixed $k$.
  \item $S_2^{\mathrm{h}}$ (within-panel left apply):
        \texttt{A[k$:$m, k{+}1$:$b] -= $\tau_k$ v$_k$ ($v_k^T$ A[k$:$m, k{+}1$:$b])}
        per $k$. The dependency carried across the $k$-axis is
        on $A[k{+}1{:}m, k{+}1{:}b]$: the next reflector reads the
        column that the previous apply just wrote. \emph{This
        statement reads $v_k$ at every column $j \in [k+1, b]$ of the
        panel;} its access function on $v_k$ is
        $\phi_{v_k}(\psi) = (i, k)$ where $\psi = (k, i, j)$ runs over
        $j$, so $\phi_{v_k}$ is \emph{not injective in $j$}.
        \textbf{Non-DAAP at the $b$-level} along the within-panel
        loop $k \in [1,b]$. With the standard recursive TSQR/CAQR
        tree (Demmel--Grigori--Hoemmen--Langou 2012), the within-panel
        loop is replaced by a balanced binary tree over row blocks of
        the panel, giving the \textbf{Quasi-DAAP} structure of
        \S2.5 with a $\log_2 P_r$ multiplier on bandwidth.
  \item $S_3^{\mathrm{h}}$ (compact-$T$ construction):
        \texttt{T = LarftBuild(V, $\tau$)} forms the upper-triangular
        $T \in \mathbb{R}^{b \times b}$ such that $Q = I - V T V^T$
        equals the product of reflectors $\prod_k H_k$. The
        construction is itself $b$ small triangular solves of the
        form $T[1{:}k{-}1, k] = -\tau_k T[1{:}k{-}1, 1{:}k{-}1]\,
        V[k{:}m, 1{:}k{-}1]^T v_k$. Each solve reads
        $T[1{:}k{-}1, 1{:}k{-}1]$ — non-injective on the iteration
        index $k$, so again \textbf{Quasi-DAAP}; the work term is
        $O(b^3)$ (lower order; absorbed into the panel constant).
\end{itemize}
Net result: the panel construction in Path~(h) is \textbf{Quasi-DAAP}.
Compared to Path~(s)'s constant-3 multiplier (3 sCQR3 iterations),
Path~(h) inherits the $\log_2 P_r$ TSQR multiplier on the
panel-bandwidth term. For $P_r \le 2^6$ this is competitive; for
larger row-grids Path~(s) is strictly better on the panel term.

\paragraph{Phase Q2$_{\mathrm{h}}$ --- Trailing update via WY (DAAP-proper).}
The trailing update $A := (I - V T^T V^T) A$ on the right-hand
block $A[:, k+b:N]$ is exactly three BLAS-3 GEMMs:
$Y := V^T A$ (size $b \times (N{-}k{-}b)$), then
$Y := T^T Y$, then $A := A - V Y$. Each of these three statements has
\emph{injective} access on its own iteration domain (every output
entry of $Y$ depends on a unique iteration vector; same for the final
$A$ subtract). Phase Q2$_{\mathrm{h}}$ is therefore
\textbf{DAAP-proper}, identical in DAAP class to Path~(s)'s
$\{S_4, S_5\}$ and inheriting the same dominant-term
$Q_{\text{par}} = \Theta(N^3/(P\sqrt{M}))$ bandwidth bound.

\paragraph{Phase Q3$_{\mathrm{h}}$ --- Cross-replica reduce of the panel.}
With 2.5D replication ($c > 1$, \S\ref{sec:A3}), the $b$ Householder
vectors built during $S_1^{\mathrm{h}}$ must be made consistent across
the $c$ replicas owning the active panel column-strip. The
communication-avoiding TSQR (Demmel et al.\ 2012) achieves this with
$b^2 \log_2 P_r$ words and $\log_2 P_r$ synchronization rounds per
processor --- the Quasi-DAAP bound of \S2.5 applied to the
$\log$-tree reduction. Compare Path~(s): $\le 6 b^2$ words,
$6 \lceil \log_2 P_r \rceil$ rounds, no $\log$-factor on bandwidth.

\paragraph{Phase Q4$_{\mathrm{h}}$ --- Pivoting.}
Identical to Path~(s): genuinely non-DAAP, $\Omega(N)$ panel sweeps.
Neither path natively supports column pivoting.

\subsection*{Path~(g): Givens rotations (DAAP-proper)}
\label{sec:A1-DAAP-givens}

The Givens variant zeroes the strict subdiagonal of the $m\times b$ panel
through a sequence of plane rotations $G(i_t, i_b; c, s)$ where
$(c, s)$ are chosen from $A[i_t, j], A[i_b, j]$ to annihilate the lower
entry.  For dense $A$ this is no faster on flops than Householder, but
its access pattern is strictly tighter:

\paragraph{Phase Q1$_{\mathrm{g}}$ --- Panel factorization (classical Givens).}
For each column $j$, eliminate $A[j{+}1{:}m, j]$ bottom-up by applying
$m{-}j{-}1$ rotations between adjacent row pairs $(i{-}1, i)$.  Each
rotation reads two rows of length $(n - j)$ and writes the same two
rows --- access function
$\phi^{\text{in}}_A(j, i, k) = \{(i{-}1, k), (i, k)\}$,
$\phi^{\text{out}}_A(j, i, k) = \{(i{-}1, k), (i, k)\}$, both injective
on $(j, i, k)$.  The panel is therefore \textbf{DAAP-proper}.

\paragraph{Phase Q3$_{\mathrm{g}}$ --- Panel construction reduction.}
Across the $c$ replicas of the 2.5D grid (\S\ref{sec:A3}), only the
final $R$ block and the rotation list $\{(c_l, s_l, i_t^l, i_b^l)\}_{l=1}^{bm/2}$
need be made consistent.  In the gather-and-replicate variant used by
the C++ benchmark, each replica receives the full $m \times b$ panel
via one AllGather of $mb$ words; each replica then runs an identical
Givens reduction.  Cost: $mb$ words bandwidth per processor,
$\log_2 P_r$ rounds.  In the tournament-Givens variant
(parallel-within-column butterfly), the reduction across replicas
folds into the within-column butterfly and total cost drops to
$b^2 \log_2 P_r$ --- same as Path~(h)'s TSQR.  Either way the panel
construction reduction is sub-leading once $N \gg b \log P_r$.

\paragraph{Phase Q2$_{\mathrm{g}}$ --- Trailing update via thin $Q$.}
The accumulated Givens rotations form an orthogonal $m \times b$
matrix $Q$ (materialized in the bench by replaying the rotation list
on $I_m[:, 0{:}b]$ in reverse order).  The trailing update is then
identical to Path~(s) and Path~(h):
$W = Q^T A_{\text{tr}}$ (one GEMM), Phase~Q4 AllReduce of $W$
($b \cdot n_{\mathrm{tr}}$ words, same as Path~(s)), then
$A_{\text{tr}} \mathrel{{-}{=}} Q W$ (one GEMM).  The DAAP class is
unchanged from Path~(s)/(h).

\subsection*{Path~(q): QDWH polar decomposition (DAAP-proper modulo inner QR)}
\label{sec:A1-DAAP-qdwh}

Path~(q) computes $A = U R$ via the polar decomposition $A = U H$ and
the identity $R = U^T A$, with $U$ obtained by Halley iteration
\eqref{eq:qdwh}.  Each Halley step is one $2n \times n$ QR plus a
single $n \times n$ GEMM update.

\paragraph{Phase Q$_{\mathrm{q}}$1 --- Stacked-matrix construction.}
At iteration $k$, build the $2n \times n$ matrix
$S_k = \begin{bmatrix}\sqrt{c_k}\, X_k \\ I_n\end{bmatrix}$.
The access function $\phi^{\mathrm{out}}_{S_k}(i, j)$ writes each
entry exactly once: $S_k[i, j] = \sqrt{c_k}\, X_k[i, j]$ for $i < n$,
$S_k[i, j] = \mathbf{1}[i - n = j]$ for $i \ge n$.  Strictly
injective, hence \textbf{DAAP-proper}.

\paragraph{Phase Q$_{\mathrm{q}}$2 --- Inner $2n \times n$ QR.}
We use Path~(s) CQR2 inside Path~(q): Phase Q1$_{\mathrm{s}}$,
Q2$_{\mathrm{s}}$, Q3$_{\mathrm{s}}$, Q4$_{\mathrm{s}}$ are invoked on
$S_k$ with the same 2.5D processor grid.  DAAP class inherits
Path~(s)'s DAAP-proper status.  The row distribution is
$m_{\text{stack}}/c = 2n/c$ rows per replica; in the bench we
\emph{interleave} the X-half and I-half rows per replica so each rank
holds its own $n/c$ X-rows followed by $n/c$ identity rows.  The QR is
invariant under input row permutation, so this interleaving is exact.

\paragraph{Phase Q$_{\mathrm{q}}$3 --- $Q_1 Q_2^T$ update GEMM.}
Extract from the QR output: top $n$ rows of $Q$ are $Q_1$, bottom $n$
rows are $Q_2$.  Compute
$P := Q_1 Q_2^T \in \mathbb{R}^{n \times n}$ via one Phase~Q4-style
collective: rank $r$ holds $m_{\text{loc}} = n/c$ rows of each of
$Q_1, Q_2$; AllGather $Q_2$ across replicas
($n^2$ words total, $b n_{\mathrm{tr}}$ per round scale removed),
then each rank computes
$P_r = (Q_1)_r \cdot Q_2^T$ as $c$ rank-block GEMMs.  Each block
GEMM has injective access function on
$((i_{\text{loc}}, k), (j_{\text{loc}}, k), (i_{\text{loc}},
j_{\text{loc}}))$ over the appropriate sub-domains; combined the full
$P$ assembly is DAAP-proper.

\paragraph{Phase Q$_{\mathrm{q}}$4 --- Halley update.}
$X_{k+1} = \alpha_x X_k + \alpha_p P$ is a strictly elementwise
statement; trivially DAAP-proper.  Total path-(q) cost is
$K \approx 6$ Halley iterations $\times$ ($2n \times n$ QR
$+$ $Q_1 Q_2^T$ assembly $+$ scalar update), or roughly
$6 \cdot 7\times$ the cost of a single $n \times n$ Path~(s) QR.

\subsection*{Phase Q5 --- Look-Ahead Pipelining (scheduling, all paths)}
\label{sec:A1-lookahead}

The X-partition framework decomposes the cDAG into mutually-disjoint
subcomputations $P(X) = \{H_1, \dots, H_s\}$ that each fit in fast
memory $X$ (\S2.3). \emph{Look-ahead} is a \textbf{scheduling} of those
subcomputations across virtual processors (streams) such that
panel-$(k{+}1)$'s Phase Q1 begins as soon as the next-panel slice of
panel-$k$'s Phase Q2 has finished, instead of waiting for the entire
Phase Q2. It introduces no new statements --- it reorders the existing
$\{S_1, \dots, S_5\}$ across the cDAG without changing data
dependencies, so the DAAP classification of each statement is
unchanged: Path~(s)'s look-ahead schedule is DAAP-proper, Path~(h)'s
look-ahead schedule is Quasi-DAAP.

\paragraph{Phase Q5a --- Trailing-update partition.}
Split the Phase Q2 trailing-update statements $S_4, S_5$ on
$A[:, k{+}b : N]$ into two disjoint column slabs:
\begin{itemize}[leftmargin=2em]
  \item $A_{\text{next}} = A[:, k{+}b : k{+}2b{-}1]$ \quad (width $b$,
        the slab that becomes panel-$(k{+}1)$'s panel),
  \item $A_{\text{rest}} = A[:, k{+}2b : N]$ \quad (width $N - k - 2b$,
        the remaining tail).
\end{itemize}
By the column-disjointness, the two slabs have no data dependency
between them, so they may be scheduled on independent streams.

\paragraph{Phase Q5b --- Inter-panel dependency.}
Panel-$(k{+}1)$'s Phase Q1 reads exactly $A_{\text{next}}$ after its
$S_5$ update. Panel-$(k{+}2)$'s Phase Q1 reads $A[:, k{+}2b : k{+}3b{-}1]$
which is a sub-column of $A_{\text{rest}}$. Therefore the partial
ordering required by the cDAG is:
\[
  S_4(A_{\text{next}}); S_5(A_{\text{next}}) \prec
  \text{Panel-}(k{+}1) \prec S_4(\text{slab}_{k{+}2}); S_5(\text{slab}_{k{+}2})
  \prec \text{Panel-}(k{+}2).
\]
The look-ahead schedule on two streams $\sigma_0, \sigma_1$ is:
\begin{enumerate}[leftmargin=2em]
  \item $\sigma_0$: Phase Q1 (Panel-$k$), then Phase Q2 on
        $A_{\text{next}}$, then Phase Q2 on $A_{\text{rest}}$.
  \item $\sigma_1$: \emph{after} $\sigma_0$ finishes $A_{\text{next}}$,
        Phase Q1 (Panel-$(k{+}1)$).
\end{enumerate}
Cross-stream dependency enforced by one CUDA event between
$S_5(A_{\text{next}})$ on $\sigma_0$ and Phase Q1 launch on $\sigma_1$.
On step boundary, $\sigma_0$ waits for $\sigma_1$ before advancing.

\paragraph{X-partition cost (Phase Q5).}
Look-ahead does not change any access-function $\phi_j$; the
subcomputation set $P(X)$ is identical. The I/O cost
$Q_{\text{par}} = \Theta(N^3/(P\sqrt{M}))$ from Lemma~3
(\S\ref{sec:A2}) is unchanged. The benefit lives entirely in the
\emph{critical-path latency} $L^{\text{proc}}$ from \S\ref{sec:A3g}.
Define $t_{\text{panel}} = \alpha F_{\text{panel}} + \beta
W_{\text{panel}}$ and $t_{\text{tail}}^{\text{rest}} = \alpha
F_{\text{tail}}^{\text{rest}} + \beta W_{\text{tail}}^{\text{rest}}$.
The schedule's per-step critical-path time becomes
\[
  t_{\text{step}}^{(\text{LA})} = \max\big(t_{\text{tail}}^{\text{next}},\,
  t_{\text{tail}}^{\text{rest}} + t_{\text{panel}}\big)
  \approx t_{\text{tail}}^{\text{rest}} + \max(0, t_{\text{panel}} -
  t_{\text{tail}}^{\text{next}}),
\]
i.e., panel-$(k{+}1)$'s panel work is hidden when
$t_{\text{panel}} \le t_{\text{tail}}^{\text{rest}}$. For the
sCQR3 cube $b = \sqrt{M}$ at single-GPU $c = 1$,
$t_{\text{panel}} = O(b^2 m)$ and $t_{\text{tail}}^{\text{rest}}
= O(b N m)$, so the panel is hidden for any $b \le N$, i.e.,
\emph{always}. The realized speedup factor over the sequential
schedule is
\[
  \boxed{\;\eta_{\text{LA}} = \frac{t_{\text{seq}}}{t_{\text{LA}}} =
  1 + \frac{t_{\text{panel}}}{t_{\text{tail}}^{\text{rest}}}
  \in \big[1, 1 + b/N\big].\;}
\]
For the admissible window $b = \Theta(N/\sqrt{P})$, this is
$1 + \Theta(1/\sqrt{P})$ --- a constant-factor speedup, vanishing in
$P$. On single-GPU H200 with $b = N/\sqrt{P_1}$ and $\sqrt{P_1} \approx
11$, $\eta_{\text{LA}} \approx 1.09$ asymptotically. The visible
single-GPU benefit is therefore in the \emph{small-step regime}
($b/N$ not vanishing), where it reaches $\approx 1.3{\times}$ at
$N=4000$, $b=363$.

\paragraph{Generalization to $s$ streams.}
With $s$ look-ahead streams in a rolling pipeline, the schedule
decomposes the trailing update into $s$ column slabs of width
$b, b, \dots, b, (N{-}sb)$. The key dependency is that panel-$(k{+}j)$
for $j = 1, \dots, s{-}1$ reads slab-$(k{+}j)$, which must first be
updated in cDAG order by $Q_k, Q_{k+1}, \dots, Q_{k+j-1}$. Concretely:
\begin{itemize}[leftmargin=2em]
  \item Stream $\sigma_0$ executes $S_4, S_5$ of $Q_k$ on the
        \emph{deep tail} $A[:, k{+}sb : N]$.
  \item Stream $\sigma_j$ for $j = 1, \dots, s{-}1$ holds a partially
        constructed panel-$(k{+}j)$: each $\sigma_j$ owns the chain of
        $S_4, S_5$ applications of $Q_k, \dots, Q_{k+j-1}$ to
        slab-$(k{+}j)$, terminating in Phase Q1 ($\mathrm{scqr3}$) on
        that slab once all previous $Q$-blocks have been applied.
  \item Cross-stream synchronization: $s(s{-}1)/2$ CUDA events --- one
        for each ordered pair $(j_{\text{src}}, j_{\text{dst}})$ with
        $j_{\text{src}} < j_{\text{dst}}$ --- enforce the partial
        cDAG order.
\end{itemize}
The X-partition admissibility per stream requires
$M_{\text{stream}} = M/s$ words of fast memory (each stream's working
set must fit on its own), so the cube side
$\sqrt{M_{\text{stream}}} = \sqrt{M}/\sqrt{s}$ shrinks. The panel
$Q^{\text{proc}}_{\text{panel}}$ grows by $\sqrt{s}$:
\[
  Q^{\text{proc}}_{\text{panel}}(s)
    = \Theta\!\big(\sqrt{s} \cdot Q^{\text{proc}}_{\text{panel}}(1)\big),
  \qquad
  Q^{\text{proc}}_{\text{trail}}(s) = Q^{\text{proc}}_{\text{trail}}(1).
\]
The realized look-ahead speedup against the sequential schedule is then
\[
  \eta_{\text{LA}}(s)
   = \frac{t_{\text{seq}}}{t_{\text{LA}}(s)}
   = 1 + \min\!\Big(\tfrac{(s-1)\, t_{\text{panel}}}{t_{\text{rest}}},\,
                     \sqrt{s} - 1\Big),
\]
where the second term in the min is the cube-shrink penalty. Setting
$\partial \eta_{\text{LA}}/\partial s = 0$ at the crossover yields
\[
  \boxed{\;s_\star \;=\; \min\!\Big(
    \Big\lceil\tfrac{t_{\text{rest}}}{t_{\text{panel}}}\Big\rceil^{2/3},\,
    \sqrt{P_1}\Big).\;}
\]
For the H200 sCQR3 measurements at $N = 16000$, $b = 1454$:
$t_{\text{rest}}/t_{\text{panel}} \approx 5$, giving $s_\star \approx
3$. For $N = 30000$, $b = 2727$: $t_{\text{rest}}/t_{\text{panel}}
\approx 8$, $s_\star \approx 4$. The diminishing-returns regime is
quickly reached because $t_{\text{panel}}$ scales as $b^2 m$ while
$t_{\text{rest}}$ scales as $b N m$; the ratio is
$\Theta(N/b) = \Theta(\sqrt{P_1})$, so $s_\star$ converges to a small
constant in the admissible window. We sweep $s \in \{1, 2, 3, 4, 8\}$
empirically and report; the optimum is usually $s = 2$ for the
small-step regime ($N \le 8000$) and $s = 3$ for the large-step regime
($N \ge 16000$) on a single H200, matching the analysis.

\paragraph{Practical implementation note.}
The current single-GPU implementation in
\texttt{ext/cudaext.jl::NextLA.\_lookahead\_run!} ships with $s = 2$
hardcoded (deeper $s$ falls back to $s = 2$ for correctness). The
$s \ge 3$ generalization requires $s(s{-}1)/2$ CUDA events plus
per-stream partial-factorization buffers, but does not change the
underlying X-partition analysis: each stream uses the same $b$ and
honors the same processor grid $[P_x, P_y, P_z]$. This is tracked as
a TODO.

\subsection*{Phase Q6 --- Iterative Refinement (correction schedule, both paths)}
\label{sec:A1-IR}

When the Phase Q2 trailing-update GEMMs are carried out in a precision
$\mathbf{u}_{\text{low}}$ coarser than the working precision
$\mathbf{u}$ (e.g., FP32 or TF32 trailing inside an FP64 panel),
sCQR3's panel-level orthogonality bound (Fukaya 2018, Thm.~3.4)
remains $O(\mathbf{u})$ per panel, but the \emph{inter-panel} residual
$A - QR$ acquires an $O(\mathbf{u}_{\text{low}} \cdot \kappa)$
contamination from the low-precision trailing updates. Iterative
refinement (IR; Wilkinson~1963, Higham~1991, Carson--Higham~2018)
corrects this back to $O(\mathbf{u})$ by repeated residual updates.

\paragraph{Phase Q6 statements.}
After all Phase Q1+Q2 panels have completed, $Q^{(0)}, R^{(0)}$ are
the low-precision QR factors. The IR loop is:
\begin{itemize}[leftmargin=2em]
  \item $S_6$ (residual): $E^{(j)} := A - Q^{(j)} R^{(j)}$.
        Iteration $\psi = (i,\ell,k)$ with $\phi_E^{\text{out}}(\psi) = (i,k)$,
        $\phi_A(\psi) = (i,k)$, $\phi_Q^{\text{in}}(\psi) = (i,\ell)$,
        $\phi_R^{\text{in}}(\psi) = (\ell,k)$. All injective $\Rightarrow$
        \textbf{DAAP-proper}.
  \item $S_7$ (correction QR): $E^{(j)} = \tilde Q \tilde R$ via a
        \emph{fresh} call to the same panel-factorization phase
        (Path~(s) or Path~(h)) in working precision $\mathbf{u}$. This
        is Phase Q1+Q2 applied to $E^{(j)}$. Same DAAP class as the
        underlying phase.
  \item $S_8$ (update):
        $Q^{(j+1)} := Q^{(j)} + \tilde Q$,
        $R^{(j+1)} := R^{(j)} + \tilde R$ (block-additive update;
        injective access $\Rightarrow$ \textbf{DAAP-proper}).
\end{itemize}
A single iteration ($j = 0$) suffices when $\kappa \cdot
\mathbf{u}_{\text{low}}/\mathbf{u}$ is bounded; Carson--Higham
(Thm.~3.1) prove convergence to $O(\mathbf{u})$ residual in
$\lceil \log_{\mathbf{u}_{\text{low}}/\mathbf{u}}(\kappa) \rceil$
iterations.

\paragraph{X-partition cost (Phase Q6).}
Each IR iteration is one residual GEMM ($S_6$, $2 N^2 m$ flops),
one full QR re-factorization ($S_7$, same cost as Phase Q1+Q2 in
working precision), and one block update ($S_8$, lower-order).
Per-iteration I/O cost
$Q_{\text{IR-iter}}^{\text{proc}} = Q_{\text{par}}^{\text{Q1+Q2}}
+ 2 N^2 m / (P \sqrt{M})$, dominated by the QR re-factor. Total
\[
  \boxed{\; Q_{\text{IR}}^{\text{proc}} =
  Q_{\text{low}}^{\text{proc}} + n_{\text{iter}} \cdot
  Q_{\text{par}}^{\text{Q1+Q2,working-prec}}. \;}
\]
For 1 IR iteration with FP32 trailing (8$\times$ faster than FP64 TC)
plus 1 FP64 refinement, the net cost is approximately
$\tfrac{1}{8} Q_{\text{FP64}} + 1 \cdot Q_{\text{FP64}}
\approx 1.125\,Q_{\text{FP64}}$. The bandwidth penalty is $+12.5\%$;
the latency saving relative to a pure-FP64 sCQR3 is the difference
between FP32-trailing's working set (half the words) and FP64's (full
words) compounded over $N/b$ panels.

\paragraph{2.5D distribution.}
Phase Q6 statements inherit the processor grid
$[P_x, P_y, P_z]$ of Phase Q1+Q2: $S_6$ is a trailing-update GEMM
distributed across the $c$ replicas exactly like $S_4, S_5$;
$S_7$ is a panel-factorization recursion using the same X-partition
cube $b = \sqrt{M}$; $S_8$ is a per-replica block add (no cross-replica
communication). The Lemma~4 aggregation cost ($6b^2$ words for the
Gram matrix in sCQR3) is paid once per IR iteration.

\paragraph{Composition with Look-Ahead.}
Phase Q5 and Phase Q6 are orthogonal scheduling layers: the look-ahead
schedule of \S\ref{sec:A1-lookahead} can be applied to both the
low-precision pass and to each IR iteration's full QR. They commute
because Phase Q6's $S_7$ is a self-call to Phase Q1+Q2 with the
identical cDAG (just on $E^{(j)}$ instead of $A$).

\subsection*{Mapping each panel variant onto the 2.5D X-partition schedule}
\label{sec:per-variant-schedule}

The X-partition framework (\S\ref{sec:A2}) factors any DAAP-proper
cDAG into mutually disjoint subcomputations $\{H_1, \dots, H_s\}$ each
fitting in a per-processor fast-memory budget $X = M$. The 2.5D
scheduling (\S\ref{sec:A3}) assigns those subcomputations to a
$[P_x, P_y, P_z]$ processor grid with $P = P_x P_y P_z$ and replication
factor $c = P_z = P M / N^2$. Every variant in this paper plugs into
the same outer skeleton:

\paragraph{Processor grid (Conflux SC'21 mathematics, Kwasniewski et al.).}
For arbitrary $P, N, M$, the X-partition optimization (\S\ref{sec:A2}) prescribes:
\begin{equation}
P_1 = \frac{N^2}{M},\qquad
c   = \min\!\left(\left\lfloor\frac{P\,M}{N^2}\right\rfloor,\; \lfloor P^{1/3}\rfloor\right),\qquad
P_x = P_y = \lfloor\sqrt{P_1}\rfloor,\qquad P_z = c.
\label{eq:grid}
\end{equation}
The implementation in \texttt{cpp\_bench/derived\_schedule.hpp} realises
\eqref{eq:grid} verbatim:
\verb|compute_derived_schedule(P, N, M_fp64_words)| computes
$c_{\mathrm{cap}} = \lfloor P^{1/3}\rfloor$ and
$c_{\mathrm{mem}} = \lfloor P\,M / N^2\rfloor$, takes
$c_{\mathrm{top}} = \min(c_{\mathrm{cap}}, c_{\mathrm{mem}})$, and searches
$c \in [c_{\mathrm{top}}, 1]$ for the largest divisor of $P$ that admits a
balanced factor pair $(p_r, p_c)$ of $P_1 = P/c$ with a panel size
$b = \lfloor\min(\sqrt{M},\, N/\sqrt{P_1})\rfloor$ inside the
\S\ref{sec:A3b} admissible window $c \le b \le N/\sqrt{P_1}$.
The fast-memory budget $M$ is read at runtime from
\texttt{cudaGetDeviceProperties.totalGlobalMem} scaled by
\texttt{NEXTLA\_FASTMEM\_FRAC} (default $0.72$), divided by the matrix
element size $\sigma$ (8\,B for FP64 paths, 4\,B for the
\texttt{--matrix=fp32full} path), so $c_{\mathrm{mem}}$ is dimensionally
$\lfloor P\,M / N^2\rfloor$ in matrix elements (header
\texttt{nextla\_fast\_memory.hpp}).
The shared sub-communicator setup --- column group (same $p_y$, size
$P_x P_z$) for Phase Q3 / Q4 AllReduces, row group (same $(p_x, p_z)$,
size $P_y$) for the Phase Q1 panel Q broadcast --- lives in
\texttt{full25d\_grid.hpp} and is used by every variant bench that
realises the full $P_z>1$ schedule.

$P_1 = N^2/M$ is the number of $\sqrt{M}\!\times\!\sqrt{M}$ tiles needed to
cover the matrix; $c$ is the replication factor that the I/O lower bound
$Q_{\text{par}} \ge N^3/(P\sqrt{M})$ saturates at when memory permits. The
cube root bound $c \le P^{1/3}$ is the Solomonik-Demmel 2.5D limit beyond
which extra replication ceases to reduce communication. In the implementation we
take $c = \lfloor P^{1/3} \rfloor$ rounded down to the largest integer for
which $P/c = P_x P_y$ is a perfect square.

\emph{Concrete instantiations} on our 4 and 8 H200 GPU testbeds:
\begin{itemize}[leftmargin=2em, itemsep=0.2em]
  \item $P=4$, $c=1$: $[P_x, P_y, P_z] = [2, 2, 1]$ (pure 2D, no replication).
  \item $P=4$, $c=4$: $[P_x, P_y, P_z] = [1, 1, 4]$ (degenerate-2.5D, max replication).
  \item $P=8$, $c=2$: $[P_x, P_y, P_z] = [2, 2, 2]$ \textbf{(true 2.5D, the
        Conflux sweet spot for $P=8$)}.
  \item $P=8$, $c=8$: $[P_x, P_y, P_z] = [1, 1, 8]$ (degenerate-2.5D).
\end{itemize}

\paragraph{Data distribution} on the full $[P_x, P_y, P_z]$ grid.  Each rank
$(p_x, p_y, p_z) \in [0, P_x) \times [0, P_y) \times [0, P_z)$ owns
\[
A\!\left[\;
\frac{N}{P_x P_z}\,(p_z P_x + p_x) \,:\, \frac{N}{P_x P_z}\,(p_z P_x + p_x + 1),\;
\frac{N}{P_y}\,p_y \,:\, \frac{N}{P_y}\,(p_y + 1)
\;\right],
\]
i.e. an $m_{\mathrm{loc}} \times n_{\mathrm{loc}}$ block with
$m_{\mathrm{loc}} = N/(P_x P_z)$ and $n_{\mathrm{loc}} = N/P_y$.  Ranks
sharing $p_y$ form a \emph{column group} (size $P_x P_z = P_1 c$) used for
Phase Q3 and Q4 AllReduces; ranks sharing $(p_x, p_z)$ form a
\emph{row group} (size $P_y$) used for the Phase-Q1 broadcast of the
just-factored panel $Q$ from $p_y^{\text{panel}}$ to all other $p_y'$.

For $P_x = P_y = 1$ (degenerate case) the row group is a singleton and
the broadcast is a no-op, $m_{\mathrm{loc}} = N/P_z = N/c$ rows per rank,
and every rank owns the full column range $[0, N)$ --- this collapses to
the 1D row distribution used in \texttt{scqr3\_2p5d\_variants.cu}.
All four variant benches realise the same $P_x{\times}P_y{\times}P_z$
scaffolding via the shared \texttt{full25d\_grid.hpp} header:
\begin{itemize}[leftmargin=2em, itemsep=0.2em]
  \item \texttt{scqr3\_full25d\_bench.cu} — Path-s (CQR2 + sCQR3 via
        \texttt{--passes=}$\{2,3\}$, Phase Q5 look-ahead default-on with
        \texttt{--no-la} ablation, four matrix modes
        \texttt{--matrix=fp64|fp64mp|fp64mp\_tf32|fp32full}).
  \item \texttt{householder\_2p5d\_bench.cu} (with
        \texttt{householder\_full25d.inl}) — Path-h, FP64 full
        $P_z>1$ via gather-and-replicate \texttt{Dgeqrf+Dorgqr} on the
        replicated panel within each column group, plus the same
        Phase Q-broadcast / Phase Q2 W AllReduce pattern.
  \item \texttt{givens\_2p5d\_bench.cu} (with
        \texttt{givens\_full25d.inl}) — Path-g, FP64 full $P_z>1$ via
        the tournament-Givens parallel panel kernel ($\Theta(b\log m)$
        depth, tex \S\ref{sec:A1-DAAP-givens}) followed by
        \verb|form_Q_from_givens_kernel| applied to $I_m[:, 0{:}b]$.
  \item \texttt{qdwh\_2p5d\_bench.cu} (with \texttt{qdwh\_full25d.inl})
        — Path-q, FP64 full $P_z>1$ Halley loop with the inner CQR2
        on the $2n{\times}n$ stacked matrix using col/row sub-comms;
        the outer $Q_1 Q_2^T$ assembly issues one AllGather over
        \texttt{col\_comm} (assembling $Q_2$ rows) and one over
        \texttt{row\_comm} (assembling $Q_1$ cols), then computes
        $\mathrm{col\_size}$ local block-DGEMMs.
\end{itemize}
Each bench accepts \texttt{--px --py --pz} flags to override the grid;
when those are omitted (and \texttt{--matrix=fp64} is active), the
grid defaults to whatever
\verb|compute_derived_schedule(P, N, M_fp64_words)| returns under the
runtime-measured $M$ from \texttt{nextla\_fast\_memory.hpp}, so every
variant runs on the mathematically-prescribed grid by default with no
hand-tuning.  Non-FP64 matrix modes (MP, TF32-trail, fp32-full) fall
through to the original 1D row distribution today; the FP64 lift is
the one this paper bench-validates.  All five paths $\{$sCQR3, CQR2,
Householder, Givens (tournament panel), QDWH$\}$ also share the four
matrix-mode flags and the \texttt{--la}/\texttt{--no-la} Phase Q5
toggle, giving the 40-combo sweep per $N$ described in
\texttt{run\_full25d\_oom\_sweep.sh}.

\emph{Smoke validation} (4 H200s, $P=4$, derived $(P_x, P_y, P_z) =
(2, 2, 1)$ from \eqref{eq:grid}, $N=2048$, $b=1024$):
\begin{itemize}[leftmargin=2em, itemsep=0.15em]
  \item Path-h \texttt{hh\_fp64\_p25d+LA}:
        $\max|\mathrm{diag}(Q^T Q) - 1| = 2.0\!\cdot\!10^{-15}$,
        $t_{\min} = 42.6$\,ms.
  \item Path-g \texttt{gg\_fp64\_p25d+LA} (tournament Givens):
        $\max|\mathrm{diag}(Q^T Q) - 1| = 2.4\!\cdot\!10^{-15}$,
        $t_{\min} = 10\,899$\,ms (the tournament kernel still has
        $\Theta(b^2 \log m)$ work per panel and is dominated by panel
        depth at $b=1024$; reducing $b$ recovers throughput).
  \item Path-q \texttt{qdwh\_fp64\_p25d\_it=6}:
        $\max|\mathrm{diag}(U^T U) - 1| = 0$ (machine zero),
        $t_{\min} = 156.9$\,ms.
\end{itemize}
All three Pz>1 paths validate at machine precision, exercising the
shared col/row sub-communicator pattern.

\paragraph{Block size $b$} obeys the admissible window of \S\ref{sec:A3b}:
$c \le b \le N/\sqrt{P_1}$.  Conflux selects $b = a \cdot c$ for some
small constant $a$ (latency-optimal for tournament-pivot LU); for QR's
sCQR3 we instead pick the bandwidth-optimal cube side
$b \approx \sqrt{M}$ inside the same window, because the QR panel has
no pivot tournament and the per-panel constant is $3$ (sCQR3) or
$2$ (CQR2) rather than $\log_2 P_1$ (TSQR).  On H200,
$\sqrt{M} = \sqrt{N_M^{\text{H200}}} \approx 512$ for the typical
sizes we measure.  The implementation lets the user override $b$ via
\texttt{--b=}; the default is the heuristic ladder in
\texttt{parse\_args}.

\paragraph{Cross-rank workflow at panel $k$ (the full 2.5D schedule).}
Set $p_y^{\text{panel}} = k \,\mathrm{div}\, n_{\mathrm{loc}}$,
$k_{\text{loc}} = k - p_y^{\text{panel}} \cdot n_{\mathrm{loc}}$,
and clip $s_b = \min(b,\; N-k,\; n_{\mathrm{loc}} - k_{\text{loc}})$ so
the panel slab stays within one $p_y$-block.
\begin{enumerate}[leftmargin=2em, itemsep=0.2em]
\item \textbf{Phase Q1 (panel)} on column group $p_y^{\text{panel}}$
      ($P_1 c = P_x P_z$ ranks):
      local SYRK $G_r \!=\! A_{\text{panel}}^{(r)\top} A_{\text{panel}}^{(r)}$
      ($m_{\mathrm{loc}} b^2 / 2$ flops);
      AllReduce $G$ across the column group (Phase Q3,
      $\le 6 b^2 \log_2(P_x P_z)$ words);
      replicated POTRF $G = R^T R$ (no comm, $b^3/3$ flops);
      local TRSM $A_{\text{panel}}^{(r)} \leftarrow A_{\text{panel}}^{(r)} R^{-1}$.
      Repeat for $\mathrm{passes} \in \{2, 3\}$ to obtain $Q$ overwriting
      the panel.
\item \textbf{Q-broadcast.} Ranks in column group $p_y^{\text{panel}}$
      broadcast their $m_{\mathrm{loc}} \times s_b$ slab of $Q$ to all
      $p_y' \ne p_y^{\text{panel}}$ via the row-group NCCL collective
      ($m_{\mathrm{loc}} s_b$ words per processor, $P_y$ peers).
\item \textbf{Phase Q2 (trailing)} on all $P$ ranks: each rank computes
      its share of $A_{\text{trail}}$ (the local columns to the right of
      the panel) and runs
      $W_r = Q_r^\top A_{\text{trail}}^{(r)}$ (local GEMM,
      $m_{\mathrm{loc}} s_b n_{\mathrm{trail, loc}}$ flops);
      AllReduce $W$ across the column group of $A_{\text{trail}}^{(r)}$'s
      $p_y$ (Phase Q4, $s_b n_{\mathrm{trail, loc}}$ words);
      then $A_{\text{trail}}^{(r)} \mathrel{-}=\!Q_r W$ locally.
\end{enumerate}
This is the schedule realized in
\texttt{cpp\_bench/scqr3\_full25d\_bench.cu}.

\medskip
\begin{enumerate}[leftmargin=2em, itemsep=0.25em]
  \item \textbf{Data placement.} $A$ is row-and-column-distributed
        across the $P_1 c = P$ ranks as described above.  Each rank owns
        a $(N/(P_x P_z)) \times (N/P_y)$ slab.  When $P_x = P_y = 1$ the
        rank's local data is the full column range and the bench reduces
        to the 1D row-distribution used by the four single-path benches.
  \item \textbf{Cube side $b$.}  Block size is taken inside the
        admissible window $c \le b \le N / \sqrt{P_1}$ of
        \S\ref{sec:A3b}, with the X-partition cube prescribing
        $b \approx \sqrt{M}$ (the largest panel that fits in fast
        memory).  This is the value all four variants use by default
        in the bench.
  \item \textbf{Phase Q1 (panel).}  One panel of width $b$ is factored.
        \emph{This is the only phase that changes between variants}
        --- everything else (data placement, Phase Q2 trailing update,
        Phase Q3 panel-reduce collective, Phase Q4 trailing-reduce
        collective, Phase Q5 look-ahead, Phase Q6 refinement) is
        identical across paths.
  \item \textbf{Phase Q2 (trailing).}  Two GEMMs and one
        $W$-AllReduce, exactly $\Theta(N^2 m / P)$ flops and
        $\Theta(N^2/\sqrt{P M})$ words per processor --- the
        leading-order bandwidth bound from \S\ref{sec:A2}b.  Common to
        all variants because all variants produce an $m\times b$
        orthonormal $Q$ from the panel.
  \item \textbf{Phase Q3 (panel reduce).}  Inter-replica collective
        used by Phase Q1 to make the panel result consistent across
        the $c$ replicas.  Different per variant; see table below.
  \item \textbf{Phase Q4 (trailing reduce).}  AllReduce of the
        $W$ slab ($b \times n_{\mathrm{tr}}$).  Same for all variants.
  \item \textbf{Phases Q5, Q6.}  Orthogonal scheduling layers, defined
        in \S\ref{sec:A1-lookahead}, \S\ref{sec:A1-IR}; both apply
        unchanged to every variant.
\end{enumerate}

\paragraph{What each variant changes in Phase Q1.}

The four variants differ only in how they realize the panel
factorization in Phase Q1, and consequently how they pay for Phase Q3.
The X-partition cube side $\sqrt{M}$, the trailing-update cost, the
$c = P M / N^2$ replication factor, and the per-step $M$-budget are
\emph{identical} across the four paths.

\begin{center}
\renewcommand{\arraystretch}{1.15}
\begin{tabular}{l p{0.20\textwidth} p{0.22\textwidth} p{0.20\textwidth}}
\toprule
Variant
  & Phase Q1 statement set
  & Phase Q3 collective
  & Where in the X-partition cube each $m{\times}b$ panel sits \\
\midrule
Path~(s) sCQR3
  & 3 iterations of  $S_1{:}G\!=\!A_p^T A_p$,  $S_2{:}\mathrm{POTRF}(G)$,  $S_3{:}A_p\!\leftarrow\!A_p R^{-1}$
  & 3 AllReduce of $G$ ($\le 6 b^2$ words, constant in $P_r$)
  & Single $\sqrt{M}\times\sqrt{M}\times b$ slab in the cube; each
    panel column is read 3 times (the 3 sCQR3 passes), each in fast
    memory \\
Path~(s) CQR2
  & 2 iterations of the same triple
  & 2 AllReduce of $G$ ($\le 4 b^2$)
  & Same slab as sCQR3 with one fewer pass; only valid when
    $\kappa_2(A) \le \mathbf{u}^{-1/2}$ (Yamamoto et al.\ 2015) \\
Path~(h) Householder
  & TSQR / gather-and-replicate
    $S_1^h{:}\mathrm{geqrf}$,
    $S_2^h{:}\mathrm{orgqr}$, then panel write-back
  & AllReduce of $V, R$ ($b^2 \log_2 P_r$) for TSQR, or AllGather
    of the panel ($m b$) for gather-and-replicate
  & Cube slab is the same; the panel is read once but the panel-reduce
    inflates Phase Q3 by $\log_2 P_r$ (TSQR) or $m/b$
    (gather-and-replicate, used in the bench) \\
Path~(g) Givens
  & $b (m - b/2)$ rotations
    $S_1^g{:}(c, s) = (a, b)/\sqrt{a^2{+}b^2}$,
    $S_2^g{:}\text{apply}$
  & Same as Path~(h) (gather-form) or
    butterfly tournament ($b^2 \log_2 P_r$, tournament-Givens)
  & Cube slab same; sequential depth of the panel kernel is
    $\Theta(b m)$ classical or $\Theta(b \log m)$ tournament; both
    fit the same $M$-budget \\
Path~(q) QDWH
  & Inner Path~(s)/(g)/(h) QR on the $2n \times n$ stack
    $\begin{bsmallmatrix}\sqrt{c_k}\,X_k\\ I_n\end{bsmallmatrix}$;
    outer Halley update $X_{k+1} = \alpha_x X_k + \alpha_p Q_1 Q_2^T$
  & Inherits the inner QR's Phase Q3, repeated $K{\approx}6$ times,
    plus one $Q_2$ AllGather ($n^2$) per iter
  & Each Halley iter sweeps a $2n\times n$ panel through the cube
    (twice the rows of the outer panel); the cube side is unchanged,
    but each iter pays the full Phase Q2--Q3--Q4 cost of an $n\times
    n$ QR \\
\bottomrule
\end{tabular}
\end{center}

\paragraph{Per-variant Phase Q1 schedule, walked through.}

For each variant we trace one panel iteration $k$ on the 2.5D grid
$(c, P_1)$ with $P_1 = 1$ (the C++ bench's degenerate-2.5D case;
$P_x = P_y = 1$, $P_z = c$).  Block size $b = b^\star$ from
\S\ref{sec:A3b}.  Recall the per-processor fast-memory footprint
inside the cube is $b^2$ for the Gram, $b \cdot n_{\mathrm{tr}}/c$ for
the local trailing slab, and $m_{\mathrm{loc}} \cdot b$ for the local
panel slab, all simultaneously resident.

\begin{itemize}[leftmargin=2em, itemsep=0.4em]
\item \textbf{Path~(s) sCQR3 / CQR2.} Each replica $r \in [0, c)$
  owns its $m_{\mathrm{loc}} \times b$ slice of the panel.
  \begin{enumerate}[leftmargin=1.5em, label=\textit{(\arabic*)}]
    \item Local SYRK: $G_r = A_r^T A_r$  ($b \times b$, $m_{\mathrm{loc}} b^2$ flops).
    \item Phase Q3: AllReduce-sum $G_r$ across the $c$ replicas
          ($\le 6 b^2$ words, constant in $P_r$ for the
          degenerate 2.5D case).
    \item Replicated POTRF of $G$ on every replica
          ($\Theta(b^3)$ flops, no comm).
    \item Local TRSM: $A_r \leftarrow A_r R^{-1}$ ($m_{\mathrm{loc}} b^2$).
    \item Repeat (1)--(4) for 2 (CQR2) or 3 (sCQR3) iterations.
  \end{enumerate}
  X-partition footprint: $G$ ($b^2$), $A_r$ panel ($m_{\mathrm{loc}} b$),
  POTRF work ($b^2$) --- all in fast memory simultaneously, all proper
  sub-cubes of the $\sqrt{M}^3$ cube.  Phase Q3 is the only
  inter-replica step.

\item \textbf{Path~(h) Householder (gather-and-replicate).}
  \begin{enumerate}[leftmargin=1.5em, label=\textit{(\arabic*)}]
    \item AllGather panel: each replica sends its
          $m_{\mathrm{loc}} \times b$ slice; the result is the full
          $m \times b$ panel replicated on every replica.
          Bandwidth: $m b / c$ words per replica.
    \item Replicated cuSolver Dgeqrf on the full panel
          (per-replica work $\Theta(m b^2)$ flops).
    \item Replicated cuSolver Dorgqr to materialize the thin
          $Q$ as an $m \times b$ orthonormal matrix.
    \item Memcpy2D each replica's local Q rows back into
          $A_r$ (no inter-replica comm).
  \end{enumerate}
  X-partition footprint: full panel ($m b$) sits in fast memory
  (Hopper's 141\,GB / SM 228\,KB combined budget allows $m b
  \ll M_\text{global}$ at the sizes considered).  Phase Q3 is the
  AllGather of the panel ($m b$) --- larger than Path~(s)'s $b^2$ by
  a factor $m/b$.  This is the lower-order $\log_2 P_r$-style
  multiplier (TSQR would replace AllGather with $b^2 \log_2 P_r$).

\item \textbf{Path~(g) Givens (gather-and-replicate variant).}
  \begin{enumerate}[leftmargin=1.5em, label=\textit{(\arabic*)}]
    \item AllGather panel (same as Path~(h)).
    \item Replicated Givens panel kernel: for $j = 0..b{-}1$, for
          $i = m{-}1$ down to $j{+}1$, compute
          $(c_l, s_l) = (A[i{-}1, j], A[i, j])/r$, apply
          $G(i{-}1, i; c_l, s_l)$ to rows $(i{-}1, i)$, columns
          $j..b{-}1$.  Save $(c_l, s_l)$ in a list of length
          $b(m{-}1)$.
    \item Replay rotations in reverse order onto $I_m[:, 0{:}b]$ to
          materialize the explicit thin $Q$ (m × b orthonormal).
    \item Memcpy2D each replica's local $Q$ rows back.
  \end{enumerate}
  X-partition footprint: full panel ($m b$) + rotation list
  ($2 b m$ words for $(c, s)$ pairs) + $Q$ buffer ($m b$).  All proper
  sub-cubes of $\sqrt{M}^3$ for reasonable $b$ on H200.  The Givens
  panel kernel runs as one CUDA block per replica with sequential
  $(j, i)$ progression; this is the classical column-by-column
  schedule, which preserves DAAP-proper status at the cost of
  serial depth $\Theta(b m)$. Tournament-Givens
  (\S\ref{sec:A1-DAAP-givens}) replaces the $i$-loop with a $\log m$
  butterfly while keeping the same total flop count, dropping the
  depth without changing the X-partition footprint.

\item \textbf{Path~(q) QDWH.}
  For $k = 0, 1, \dots$ until $|\ell_k - 1| < \epsilon$:
  \begin{enumerate}[leftmargin=1.5em, label=\textit{(\arabic*)}]
    \item Compute Halley weights $(a_k, b_k, c_k)$ from $\ell_k$
          (analytic, $\Theta(1)$).
    \item Build the stacked $S_k$ ($2n \times n$, row-distributed
          with $m_{\mathrm{st,loc}} = 2 n /c$ rows per replica; top
          $n/c$ rows are $\sqrt{c_k} X_k$, bottom $n/c$ are identity
          rows starting at global row $r n /c$).
    \item Inner 2.5D QR on $S_k$ using Path~(s) CQR2 (or any other
          inner variant).  This re-enters the full Phase Q1--Q4
          pipeline on a $2n \times n$ matrix.  The X-partition cube
          is the same $\sqrt{M}^3$; the inner cube is swept twice as
          many times in the row direction because there are twice as
          many rows.
    \item Extract $Q_1$ (top $n/c$ rows of each replica's local $S_k$
          after the inner QR), $Q_2$ (bottom $n/c$ rows).
    \item AllGather $Q_2$ across replicas ($n^2$ words total).
    \item Each replica forms its $n/c$ rows of $P = Q_1 Q_2^T$ as $c$
          rank-block GEMMs $(Q_1)_r \cdot Q_2^T$ (each
          $n/c \times n/c$ from a $n/c \times n$ multiply).
    \item Element-wise update $X_{k+1} = \alpha_x X_k + \alpha_p P$.
    \item Update $\ell_{k+1} = \ell_k (a_k + b_k \ell_k^2)/(1 + c_k \ell_k^2)$.
  \end{enumerate}
  X-partition footprint: stacked $S_k$ ($2 n^2 / c$ per replica),
  inner-QR scratch ($b^2$ for $G$, $b \cdot n / c$ for $W$),
  $Q_2$ AllGather buffer ($n^2$ per replica), $P$
  ($n^2 / c$), $X_k, X_{k+1}$ ($2 n^2 / c$). All proper sub-cubes of
  the cube when $n^2 / c \le M$.  Each Halley iter costs roughly
  $7 \times$ the per-iter cost of a single $n \times n$ Path~(s) QR
  (the $2n \times n$ stacked QR alone is $\sim 6 \cdot 2 n^3 / 3$
  flops; the $Q_1 Q_2^T$ assembly adds $\sim n^3$; the rest is
  $o(n^3)$).
\end{itemize}

\paragraph{Why Path~(s) is the strictly tighter result for dense QR.}
At the leading-order bandwidth $Q_{\text{par}} = \Theta(N^3 / (P \sqrt{M}))$
all four paths agree.  At the sub-leading panel-reduce term, Path~(s)
pays $\le 6 b^2$ per panel (constant in $P_r$), Path~(h) pays
$m b$ (gather) or $b^2 \log_2 P_r$ (TSQR), Path~(g) pays the same as
Path~(h) in the gather variant.  Path~(q) is $K$-fold the cost of
Path~(s) on the inner QR.  Hence for dense, well-conditioned $A$, the
ranking from the X-partition framework alone is exactly
Path~(s) $\prec$ Path~(g) $\approx$ Path~(h) $\prec$ Path~(q), which
the empirical numbers below confirm.

\paragraph{When the ranking flips.}
Path~(g) is the choice for sparse / structured / banded $A$, where
each Givens rotation touches $O(1)$ nonzeros and the panel
factorization cost drops to $O(b \cdot \text{nnz}(A))$, breaking the
$m b^2$ Path~(s) lower bound on the panel.  Path~(h) is the choice
under arbitrary conditioning, since Path~(s) requires
$\kappa_2(A) \le \mathbf{u}^{-1}/(96(mb + b(b{+}1)))$
(Fukaya~2018 Thm.~3.4) and CQR2 requires the strictly tighter
$\kappa_2(A) \le \mathbf{u}^{-1/2}$ (Yamamoto et al.\ 2015).  Path~(q)
is the choice when one needs $U$ explicitly (polar decomposition, SVD
intermediate, matrix sign function); for those purposes the
$K$-fold cost is absorbed by what would otherwise be an SVD-by-QR
plus extra work.

\subsection*{Empirical validation of the schedule on 4 H200 GPUs}
\label{sec:empirical}

The schedule of \S\ref{sec:A3} is realized in C++/CUDA in
\texttt{cpp\_bench/scqr3\_2p5d\_variants.cu} using exactly the paper's
primitives: one MPI rank per GPU, processor grid $[P_x, P_y, P_z] =
[1, 1, c]$ (the degenerate-2.5D case for $P_1 = 1$ that
\S\ref{sec:A3a} explicitly admits), row-distribution of $A$ across the
$c$ replicas, NCCL Allreduce across replicas for the panel Gram (Phase
Q3, Lemma~\ref{lem:reduction-LB}, $\le 6 b^2$ words per processor) and
the trailing-update $W$ slab. Two CUDA streams per rank
(\emph{compute}, \emph{comm}) carry cuBLAS / cuSolver / NCCL work with
event-based handoff and \emph{zero host-side synchronization} in the
panel loop.

The C++ binary accepts the variant family of Phase Q1+Q2+Q5+Q6 via flags
\texttt{--passes=}$p$ ($p \in \{2,3\}$ selects CQR2 or sCQR3),
\texttt{--mp} (Phase Q6 low-precision trailing GEMM in FP32),
\texttt{--la} (Phase Q5 two-stream look-ahead pipelining),
\texttt{--ir=}$j$ (Phase Q6 iterative refinement, $j$ rounds of
re-orthogonalization in the working precision). Block size is taken
from the X-partition cube $b \approx \sqrt{M}$ inside the admissible
window $c \le b \le N/\sqrt{P_1}$ of \S\ref{sec:A3b}; for the H200 fast
memory model the cube lands at $b = 512$ for $N \in [8000, 32000]$ and
the empirical optimum is reached there for every variant.

Comparison against NVIDIA's distributed-memory QR \texttt{cusolverMpGeqrf}
(libcusolverMp v0.8) on the same 4 H200s, identical $N$, FP64 throughout,
$5$ timed runs after $2$ warmups, MPI barriers around each
timed region:

\begin{center}
\renewcommand{\arraystretch}{1.05}
\begin{tabular}{l ccc | c}
\toprule
variant (4 GPUs, $c{=}4$) & $N{=}8000$ & $N{=}16000$ & $N{=}32000$ & speedup pattern \\
\midrule
sCQR3 (\texttt{passes=3})          & 47.35  & 135.84 & 578.34 & 1.66--2.47$\times$ \\
CQR2 (\texttt{passes=2})           & 33.86  & 108.47 & 510.79 & 1.88--3.45$\times$ \\
sCQR3 + Look-ahead (Phase Q5)      & 42.24  & 123.91 & 537.35 & 1.79--2.77$\times$ \\
\textbf{CQR2 + Look-ahead $\star$} & \textbf{30.45} & \textbf{98.05} & \textbf{474.06} & \textbf{2.02--3.84}$\times$ \\
sCQR3 + MP (FP32 trailing)         & 48.96  & 154.34 & 730.59 & 1.31--2.39$\times$ \\
CQR2 + MP                          & 36.52  & 125.22 & 666.07 & 1.44--3.20$\times$ \\
sCQR3 + MP + IR ($j{=}1$, Phase Q6)& 62.80  & 183.38 & 800.73 & 1.20--1.86$\times$ \\
sCQR3 + MP + LA + IR               & 57.37  & 168.59 & 743.86 & 1.29--2.04$\times$ \\
\midrule
NVIDIA \texttt{cusolverMpGeqrf}    & 116.93 & 317.24 & 959.86 & 1$\times$ (baseline) \\
\bottomrule
\end{tabular}
\end{center}

\paragraph{Full $[P_x,P_y,P_z]$ grid (Path~(s) in \texttt{scqr3\_full25d\_bench.cu}).}
The same Phase~Q1--Q4 skeleton of \S\ref{sec:per-variant-schedule} was
timed with $b = 512$ (the cube side $\sqrt{M}$ from \S\ref{sec:A3b}),
\texttt{mpirun} $P{=}4$ ranks, and NVIDIA \texttt{cusolverMpGeqrf} on a
$2{\times}2$ rank grid (MB${=}256$ as required by the vendor bench).
Medians of five timed runs after two warmups:

\begin{center}
\renewcommand{\arraystretch}{1.05}
\begin{tabular}{l l ccc | c}
\toprule
layout & Path~(s) flag & $N{=}8000$ & $N{=}16000$ & $N{=}32000$ & vs.\ cuSOLVERMp \\
\midrule
cuSOLVERMp $2{\times}2$ & --- & 1036 & 1030 & 1776 & $1\times$ \\
$[2,2,1]$, $c{=}1$ & \texttt{passes=3} & 102 & 506 & 930 & $1.9$--$10\times$ \\
$[2,2,1]$, $c{=}1$ & \texttt{passes=2} (CQR2) & 175 & 158 & 930 & $1.9$--$6.6\times$ \\
$[1,1,4]$, $c{=}4$ & \texttt{passes=3} & 62.6 & 601 & 619 & $1.7$--$17\times$ \\
$[1,1,4]$, $c{=}4$ & \texttt{passes=3+LA} & 68.3 & 150 & 593 & $1.7$--$15\times$ \\
\bottomrule
\end{tabular}
\end{center}

The $[1,1,4]$ row at $N{=}16000$ shows where Phase~Q5 look-ahead
(\S\ref{sec:A1-lookahead}) pays under the degenerate mesh: median
\texttt{passes=3} is dominated by variance, while \texttt{passes=3+LA}
recovers pipelining.  At $N{=}32000$, \texttt{passes=2+LA} ties the best
median ($535$\,ms); we still tabulate \texttt{passes=3+LA} for a uniform
sCQR3 stability profile.  The pure 2D layout $[2,2,1]$ matches the
$c{=}1$ column of the memory bound~\eqref{eq:c-memory}; the degenerate
layout matches $c{=}4{=}P$ (maximal replication on four ranks).

\paragraph{Eight GPUs, true $2.5$D $[2,2,2]$ ($c{=}2$).}
The same executable supports $P{=}8$ with $P_x{=}P_y{=}P_z{=}2$.
Automated medians for the full-size sweep (cuSOLVERMp at $4{\times}2$,
$[4,2,1]$, $[2,2,2]$, and $[1,1,8]$) are produced by
\texttt{cpp\_bench/run\_full25d\_sweep.sh} on the \texttt{full-node}
partition; a queued Slurm driver is \texttt{run\_full25d\_8gpu\_pipeline.sbatch}
(submit with \texttt{sbatch --dependency=afterany:\textit{PRIOR\_JOB}}
to wait for other users on the shared DGX).  Smoke timings for
$N{=}4000$ at $[2,2,2]$ are written to the same log file when that job
runs.

All times in milliseconds; max $|\mathrm{diag}(Q^T Q) - 1|_\infty$ stays
in $[6.7\text{e}{-16}, 7.1\text{e}{-15}]$ for every variant, well inside
the Fukaya orthogonality bound $50 \cdot m \cdot b \cdot \mathbf{u}
\approx 10^{-8}$ at these dimensions.

\textbf{The fastest variant is CQR2 + Look-Ahead}, the composition of
$\mathrm{passes} = 2$ from \S\ref{sec:A1-DAAP} (still DAAP-proper, with
the stability guarantee $\kappa_2(A) \le \mathbf{u}^{-1/2}$ from
Yamamoto et al.~2015) and the Phase Q5 two-stream schedule of
\S\ref{sec:A1-lookahead}. It beats NVIDIA's cusolverMpGeqrf by
$2.02{-}3.84\times$ across the tested range. The win is not from raw
flop reduction --- both schemes incur $\Theta(N^3 / (P \sqrt{M}))$
trailing-update bandwidth --- but from (i) the $\le 6 b^2$ panel
Allreduce of Phase Q3 (a constant rather than $\log P_r$ scaling), (ii)
the cube-side $b \approx \sqrt{M}$ keeping each replica's working set
in fast memory, and (iii) the Phase Q5 pipelining hiding panel
factorization behind the trailing update. Each piece is derived from
the X-partition framework of \S\ref{sec:A2}.

\paragraph{Observations on Phase Q6 (mixed precision $+$ IR).}
The current C++ implementation routes the Phase Q6 low-precision pass
through \texttt{cublasSgemm} (FP32 ALU, 134\,TFLOPS Tensor Core on
Hopper), which is only $2\times$ faster than the FP64 Tensor Core path
already used by cuBLAS DGEMM ($67$\,TFLOPS). Net of the FP64 $\leftrightarrow$
FP32 conversion bandwidth, MP is a small regression in this
configuration. The unrealized lever is the TF32 path
(\texttt{cublasGemmEx} with \texttt{CUBLAS\_COMPUTE\_32F\_FAST\_TF32},
$989$\,TFLOPS), which would change the analysis: at TF32 speeds the
low-precision pass would be $\sim 15 \times$ cheaper, the IR cost
would amortize, and Phase Q6 would dominate at sufficiently large $N$.
This is a one-line change in the bench and is left as a future
empirical run; the schedule of \S\ref{sec:A1-IR} is unchanged.

\section{I/O Lower Bounds}
\label{sec:A2}

\subsection{Dominant statement}
\label{sec:A2a}

The trailing update statement $S_5$ is the global dominant. The two
GEMMs $W = Q^T A$, $A {-}{=} QW$ have access vectors:
\begin{itemize}[leftmargin=2em]
  \item $\phi^{\mathrm{out}}_A = (i,p)$, $\dim = 2$
  \item $\phi_Q = (i,\ell)$, $\dim = 2$
  \item $\phi_{Q^T} = (q,\ell)$, $\dim = 2$
  \item $\phi^{\mathrm{in}}_A = (q,p)$, $\dim = 2$
\end{itemize}
Indices $i,q$ range over $[k+b, N]$; $p$ over $[k+b, N]$; $\ell$ over
$[1,b]$. This is structurally identical to the Householder trailing
update.

\subsection{X-partition optimization}
\label{sec:A2b}

The four-variable problem in $|D^i|, |D^p|, |D^q|, |D^\ell|$ is
\[
\max\ |D^i|\,|D^p|\,|D^q|\,|D^\ell|
\quad \text{s.t.}\quad
|D^i||D^p| + |D^i||D^\ell| + |D^q||D^\ell| + |D^q||D^p| \le X,
\quad |D^t| \ge 1.
\]

\paragraph{Reduction to three variables.}
We invoke Lemma~\ref{lem:input-reuse} on the variable $\ell$, whose
domain $|D^\ell| = b$ is bounded uniformly. By Remark~\ref{rem:bN-bound},
the input-reuse hypothesis $b(N-k) \le M$ holds for all
$k \in \{0,b,\ldots,N-b\}$ in the 2.5D regime. Absorbing
$|D^\ell| = b$ into the input-reuse bookkeeping collapses the
constraint to the GEMM form:
\[
\max_{x,y,z}\ xyz
\quad\text{s.t.}\quad
xy + xz + yz \le X,
\quad x,y,z \ge 1,
\]
where $x = |D^i|, y = |D^p|, z = |D^q|$.
\REV{C1: collapse now derived from Lemma~\ref{lem:input-reuse} +
Remark~\ref{rem:bN-bound}, with the validity of the absorbing step
verified for the entire range of outer steps $k$.}

\paragraph{KKT derivation (interior critical point and uniqueness).}
The constraint is active at any optimum (increasing any variable
strictly increases the objective without violating boundary
constraints, until the inequality becomes tight). Set
$g(x,y,z) = xy + xz + yz - X$. Lagrangian:
$\mathcal{L} = xyz - \lambda g$.
Stationarity:
\begin{align}
yz &= \lambda(y+z), \label{eq:kkt1} \\
xz &= \lambda(x+z), \label{eq:kkt2} \\
xy &= \lambda(x+y). \label{eq:kkt3}
\end{align}

\textbf{Uniqueness in the open positive orthant.}
\REV{C1: explicit verification of uniqueness now provided.}
From \eqref{eq:kkt1}: $\lambda = yz/(y+z)$. From \eqref{eq:kkt2}:
$\lambda = xz/(x+z)$. Setting equal:
\[
\frac{yz}{y+z} = \frac{xz}{x+z}
\;\Longleftrightarrow\;
yz(x+z) = xz(y+z)
\;\Longleftrightarrow\;
xyz + yz^2 = xyz + xz^2
\;\Longleftrightarrow\;
z^2(y-x) = 0.
\]
Since $z > 0$, we obtain $y = x$. By the same calculation with
\eqref{eq:kkt2} and \eqref{eq:kkt3}, $z = x$. Therefore the unique
interior stationary point of the Lagrangian is $x = y = z$.
Substituting in the active constraint: $3x^2 = X$, so
$x = y = z = \sqrt{X/3}$, and
\[
\boxed{\;\chi(X) \;=\; (X/3)^{3/2}.\;}
\]

\textbf{Boundary cases excluded.}
\begin{itemize}[leftmargin=2em]
  \item If any of $x,y,z \to 0$, then $f = xyz \to 0$, which is strictly
    suboptimal compared to the interior value $(X/3)^{3/2} > 0$ for
    $X > 0$. So the optimum is not on the $x_t = 0$ boundary of the
    open orthant.
  \item If any $x,y,z \to \infty$, the constraint $xy + xz + yz \le X$
    is violated (any single product term among them then exceeds $X$).
    So the optimum lies in a bounded region of the open orthant.
  \item The integer constraints $x,y,z \ge 1$ on the relaxed domain
    are non-binding when $X > 3$, which is the only regime of interest
    (the GEMM bound is interesting only when $X \gg M \gg 1$).
\end{itemize}
\REV{C1: degenerate cases $|D^t| \to 0$ and $|D^t| \to \infty$ now
explicitly excluded.}

\subsection{Sequential and parallel bounds}
\label{sec:A2c}

\paragraph{Optimization of $\rho(X)$.}
With $\chi(X) = (X/3)^{3/2}$:
\[
  \rho(X) \;=\; \frac{(X/3)^{3/2}}{X-M}, \qquad X > M.
\]
Differentiating and using $\chi'(X) = \tfrac{1}{2}(X/3)^{1/2}$:
\[
\rho'(X) \;=\; \frac{(X/3)^{1/2}\,(X - 3M)}{6(X-M)^2}.
\]
\textbf{Sign analysis.} $\rho'(X) < 0$ on $(M, 3M)$ and $\rho'(X) > 0$
on $(3M, \infty)$. Therefore $X_0 = 3M$ is the unique critical point
on $(M,\infty)$ and is a \emph{minimum} of $\rho$.
\REV{C1: prior claim that $X_0 = 3M$ is a maximum of $\rho$ is corrected.
The endpoint behaviour is $\rho(X) \to +\infty$ at $X \to M^+$
(numerator $(M/3)^{3/2}$ remains finite, denominator $\to 0^+$) and
$\rho(X) \to +\infty$ at $X \to \infty$ (numerator grows as $X^{3/2}$,
denominator only as $X$). The interior critical point is therefore a
minimum of $\rho$. The X-partition lower bound
$Q \ge \max_X |V|/\rho(X)$ is obtained at this minimum:
$Q \ge |V|/\rho_{\min}$. The cited literature uses ``$\rho_{\max}$''
for this value and we preserve that notation, but the directionality
of the optimization is corrected.}

Evaluating at $X_0 = 3M$:
\[
  \chi(3M) \;=\; M^{3/2},\qquad
  X_0 - M \;=\; 2M,\qquad
  \boxed{\;\rho_{\max} \;\equiv\; \rho(X_0) \;=\; \frac{M^{3/2}}{2M} \;=\; \frac{\sqrt{M}}{2}.\;}
\]

\paragraph{Geometric interpretation.}
At $X_0$, the optimal subcomputation is a cube of side $\sqrt{M}$ in
$(i,p,q)$-space, with three faces of area $M$ tiling fast memory.

\paragraph{Total work in the trailing update.}
\REV{C1: explicit summation now provided.}
Each GEMM at outer step $k$ computes $(N-k)^2 \cdot b$ multiplications
(accounting for the trailing submatrix). Summing over the $N/b$ outer
steps $k = 0, b, 2b, \ldots, N-b$:
\[
  \sum_{j=0}^{N/b - 1} (N - jb)^2 \cdot b
  \;=\; b \sum_{j=0}^{N/b-1} (N - jb)^2.
\]
Let $r = N/b$. Then the sum is
\[
  b \sum_{j=0}^{r-1} (N - jb)^2
  \;=\; b \sum_{j=0}^{r-1} b^2 (r - j)^2
  \;=\; b^3 \sum_{j=0}^{r-1} (r - j)^2
  \;=\; b^3 \sum_{i=1}^{r} i^2
  \;=\; b^3 \cdot \frac{r(r+1)(2r+1)}{6}.
\]
For $r = N/b$, this expands to
\[
  b^3 \cdot \frac{(N/b)(N/b+1)(2N/b+1)}{6}
  \;=\; \frac{N \,(N+b)\,(2N+b)}{6}
  \;=\; \frac{N^3}{3} + \frac{N^2 b}{2} + O(N b^2).
\]
The leading term per GEMM is $N^3/3$. Each statement $S_4, S_5$
contributes one such GEMM, so
\[
  \boxed{\;|V_{\mathrm{Q2}}| \;=\; 2 \cdot \frac{N^3}{3} + O(N^2 b)
  \;=\; \frac{2N^3}{3} + O(N^2 b)\;\text{multiplications}.\;}
\]
The $O(N^2 b)$ correction is lower-order in the regime
$b \ll N$ that is the only regime of interest. Dividing the leading
term by $P_1 c = P$ (the total number of processors) gives the
per-processor leading flop count $2N^3/(3P)$, used in
Section~\ref{sec:A3}.

\paragraph{Sequential lower bound.}
By the X-partition bound (Kwasniewski et al.\ SC'21, Theorem~1) with
$|V| = 2N^3/3$:
\[
  Q_{\mathrm{seq}} \;\ge\; \frac{|V|}{\rho_{\max}}
  \;=\; \frac{2N^3/3}{\sqrt{M}/2}
  \;=\; \frac{4 N^3}{3\sqrt{M}}.
\]

\paragraph{Parallel lower bound.}
By Lemma~\ref{lem:parallel-IO}:
\[
  \boxed{\;Q_{\mathrm{par}} \;\ge\; \frac{4 N^3}{3\,P\,\sqrt{M}}.\;}
\]
This recovers the Ballard--Demmel--Holtz--Schwartz (2011) bound
(Theorem~2.5 for the parallel extension; their Theorem~2.2 establishes
the sequential GEMM bound).
\REV{C6: BDHS~2011 now cited with theorem numbers (2.2 sequential, 2.5
parallel).}

\subsection{Cross-statement reuse}
\label{sec:A2d}

\begin{itemize}[leftmargin=2em]
  \item \textbf{Input reuse on $Q$ (Lemma~\ref{lem:input-reuse}).}
    $Q_k$ has $b(N-k)$ entries and is read by both $S_4$
    ($W = Q^T A$) and $S_5$ ($A {-}{=} QW$). Lemma~\ref{lem:input-reuse}
    gives a free factor of 2 on $Q$'s I/O contribution. Since $Q$'s
    contribution is dominated by the $A$-traffic in the leading term,
    this does not change the leading constant.
  \item \textbf{No $T$ matrix:} sCQR3 does not produce a $T$ factor.
  \item \textbf{Output reuse on $A$ (Lemma~\ref{lem:output-reuse}).}
    $A[i,p]$ is updated once per outer step.
  \item \textbf{Cross-panel reuse on $A$.} Lemma~\ref{lem:input-reuse}
    again amortizes the write-then-read at step $k+b$.
  \item \textbf{Panel Gram matrix reuse.} The $b \times b$ Gram matrix
    fits in fast memory ($b^2 \le M$) and contributes no slow-memory
    I/O. The cross-processor aggregation of partial Grams is a
    distinct cost (Lemma~\ref{lem:reduction-LB}).
\end{itemize}

\subsection{DAAP multipliers}
\label{sec:A2e}

Under sCQR3, the panel factorization is fully DAAP-proper. The
Quasi-DAAP $\log_2 P_r$ multiplier of TSQR is replaced by the constant
$3$. The constant enters multiplicatively only in the lower-order
panel term, not in the dominant trailing-update bound.

\subsection{Summary table}
\label{sec:A2f}

\begin{center}\footnotesize
\setlength{\tabcolsep}{3pt}
\begin{tabular}{@{}lllll@{}}
\toprule
Phase & Statements & $\rho_{\max}$ & $Q_{\text{par}}$ leading & Prior \\
\midrule
Q1 (sCQR3 panel)    & $S_1, S_2, S_3$ ($\times 3$)  & $\sqrt{M}/2$ (Gram)  & $O(N^2 b / P\sqrt M)$                    & new          \\
Q2 (trailing)        & $S_4, S_5$                    & $\sqrt{M}/2$          & $\dfrac{4 N^3}{3 P \sqrt M}$             & BDHS'11 Thm~2.5 \\
Q3 (panel agg.)      & within $S_1$                  & n/a (Lemma~\ref{lem:reduction-LB})  & $O(bN/P)$                                & replaces TSQR \\
Q4 (col pivot)       & rank-revealing                & not DAAP              & $\ge 4N^3/(3P\sqrt M)$, lat $\Omega(N)$  & BDD'10 Thm~6.1 \\
\bottomrule
\end{tabular}
\end{center}

\section{Near-Optimal Parallel Schedule (sCQR3-2.5D)}
\label{sec:A3}

\subsection{Deriving $c$, $P_1$, and $[P_x,P_y,P_z]$ (memory saturation $\cap$ Solomonik--Demmel)}
\label{sec:A3derive}

The implementation maps the abstract schedule onto a three-axis index
space $[P_x,P_y,P_z]$ with $P = P_x P_y P_z$.  Two independent
constraints fix the replication factor~$c = P_z$.

\paragraph{Constraint I (fast-memory saturation).}
Each processor owns an $N^2/P$-sized submatrix of $A$ in the
unreplicated layout.  Replicating the trailing update's contraction
across $c$ copies reduces each replica's resident matrix block by a
factor~$c$ to $N^2/(P_1 c)$ with $P_1 = P/c$ the number of distinct
``first-layer'' processors in the $[P_x,P_y]$ plane.  The 2.5D
\emph{memory constraint} requires that block to saturate fast memory:
$N^2/(P_1 c) = M$, i.e.\ $P_1 c = N^2/M$.  With $P = P_1 c$ processors
total, eliminating $P_1$ gives the classical replication bound
\begin{equation}
  c \;\le\; \frac{P\,M}{N^2}.
  \label{eq:c-memory}
\end{equation}

\paragraph{Constraint II (Solomonik--Demmel cubic limit).}
For matrix multiplication (and hence for the dominant trailing GEMMs
with the same reduction structure), Solomonik and Demmel show that
once replication exceeds $c \asymp P^{1/3}$, additional copies no
longer reduce the dominant communication term; the optimal lies at
$c = \Theta(P^{1/3})$.  In discrete hardware we cap
\begin{equation}
  c \;\le\; \bigl\lfloor P^{1/3} \bigr\rfloor
  \label{eq:c-cube}
\end{equation}
and choose the largest admissible integer satisfying both
\eqref{eq:c-memory} and~\eqref{eq:c-cube}, i.e.\
$c = \min\bigl(\lfloor PM/N^2 \rfloor,\, \lfloor P^{1/3} \rfloor\bigr)$,
matching~\eqref{eq:grid} in \S\ref{sec:per-variant-schedule}.

\paragraph{Square first-layer grid.}
Given $c$, set $P_1 = P/c$.  The X-partition / CAQR literature fixes the
first-layer processor mesh as square $P_x = P_y = \sqrt{P_1}$ so that
local tile dimensions match the $\sqrt{M} \times \sqrt{M}$ cube side
in~\S\ref{sec:A2}b--\ref{sec:A3b}.  The third axis carries replication
only: $P_z = c$.  Degenerate cases ($P_x = P_y = 1$, maximal~$c$) are
admissible and correspond to the 1D row distribution used in the
variant-only benches; the general $[P_x,P_y,P_z]$ assignment is the
Conflux-style map in \S\ref{sec:A3aprime}.

\subsection{Processor index space}
\label{sec:A3a}
The processor set is indexed by an abstract grid of dimensions
$[P_x, P_y, P_z] = [\sqrt{P_1}, \sqrt{P_1}, c]$ with $c = PM/N^2$ and
$P_1 = P/c$. Block size $b \ge PM/N^2$. The grid is purely an index
space on the processor set; \REV{C2: no physical network topology is
implied.} For the panel sub-step (Q1), the active processors are the
$\sqrt{P_1} \cdot c$ processors owning rows of the panel column block.

\subsection{Ownership and replication protocol}
\label{sec:A3aprime}

We adapt the CONFLUX block-cyclic distribution (Kwasniewski et al.,
SC'21, Section~4.1) with sCQR3-specific modifications.

\paragraph{Block-to-processor map.}
With $p_r = p_c = \sqrt{P_1}$ and block size $b$, partition $A$ into a
$\lceil N/b\rceil \times \lceil N/b\rceil$ array of $b \times b$ blocks.
Block $(I,J)$ is assigned to abstract grid coordinate
$\pi(I,J) = (I \bmod p_r,\; J \bmod p_c)$ and replicated across all $c$
indices of the third (replication) axis. \REV{C2: $\pi$ is an index
map; ``axis'' replaces ``communicator'' or ``physical dimension''.}

\paragraph{Contraction-dimension split.}
During the local accumulation $W = Q^T A$, the contraction index
$q \in [k,N)$ is split across the $c$ replication indices. Index $z$
owns the contraction stripe
\[
  q \in \big[k + z \cdot \tfrac{N-k}{p_r c},\;
              k + (z{+}1) \cdot \tfrac{N-k}{p_r c}\big)
\]
within each row block of its $r$-coordinate, and computes the partial
$W^{(z)} = Q_{\mathrm{stripe}}^T A_{\mathrm{local}}$. The cross-replica
sum $W = \sum_{z=0}^{c-1} W^{(z)}$ is performed as a reduction along
the replication axis, leaving each replica index $z$ holding its
$1/c$-th slice of the final $W$. \REV{C2: replaced
``\texttt{MPI\_Reduce\_scatter} along $P_z$'' with the abstract
reduction-and-scatter operation, scored only by its words-moved and
synchronization-rounds counts (Lemma~\ref{lem:reduction-LB}).}

\paragraph{Panel index set.}
At outer step $k$, the active panel column block has
$j_k = (k/b) \bmod p_c$. Define
\[
  C^{(k)}_{\mathrm{panel}}
    \;=\; \{(r,j_k,z)\,:\, r \in [0,p_r),\; z \in [0,c)\},
\]
a set of $p_r \cdot c$ processors. The three sCQR3 Gram-matrix
aggregations (one per CholeskyQR iteration) are performed on
$C^{(k)}_{\mathrm{panel}}$ as additive reductions (commutative,
associative) producing the global $G = Q^T Q$ on every member.
\REV{C2: ``panel communicator'' replaced by ``panel index set''; the
aggregation is scored against Lemma~\ref{lem:reduction-LB}.}

\paragraph{Row-broadcast of $Q_k$.}
After the panel step, processors at column index $s \neq j_k$ require
$Q_k$ for the trailing GEMM. Each processor $(r,j_k,z)$ broadcasts its
$Q_k$ rows to $(r,s,z)$ for all $s \in [0, p_c)$. Words moved per
processor: $(N-k)b/p_c$ (one $p_c$-way broadcast).

\paragraph{Lazy replication on the trailing update.}
Step~3 (local $W$ build) is purely local; step~4 (cross-replica reduce)
runs along the replication axis only; step~5 (apply $A {-}{=} QW$) is
local. Because every member of $C^{(k)}_{\mathrm{panel}}$ produced the
same $Q_k$ and the cross-replica reduction of $W$ is associative
summation, replicas remain consistent across outer steps without an
explicit cross-replica broadcast of $A$.

\subsection{Data distribution}
\label{sec:A3b0}
\begin{itemize}[leftmargin=2em]
  \item $A$: 2D block-cyclic over $[P_x, P_y]$ with block $b$,
    replicated along $P_z$ during reads
    (Lemma~\ref{lem:input-reuse}: $Q_k$ is the read-only input common
    to all $P_z$ replicas) and reduced along $P_z$ at the end of each
    outer step.
  \item $Q_k$: same layout as $A$'s factored column block; replicated
    read along $P_z$. Stored in-place in $A$'s lower-trapezoidal storage.
  \item $R_k$: $b\times b$ upper triangular, replicated on every
    processor (size $b^2 \le M$, produced identically by the
    aggregation).
  \item $W = Q^T A$: distributed $[P_x,P_y]$ but accumulated
    independently on each $P_z$ replica.
  \item $R$ (final triangular): occupies the upper triangle of $A$'s
    storage, distribution unchanged.
\end{itemize}

\subsection{Memory budget per processor}
\label{sec:A3b}

We verify that the per-processor working set fits in fast memory $M$.
Setting $P_1 = P/c$ and $c = PM/N^2$, every processor holds:
\begin{center}
\begin{tabular}{@{}llr@{}}
\toprule
Item & Size (entries) & Notes \\
\midrule
Local block of $A$               & $\dfrac{N^2}{P_1} = \dfrac{N^2 c}{P} = M$       & 2.5D constraint, exact equality \\[6pt]
$Q_k$ in-place                    & $0$                                              & overlays $A$'s lower trapezoid \\
Local share of $W = Q^T A$        & $\dfrac{b\,N}{\sqrt{P_1}} = b\sqrt{Mc}$         & step~3 output, before reduce \\[6pt]
Gram matrix scratch $G$           & $b^2$                                            & one per sCQR3 iter, reused \\
$R_k$ scratch                     & $b^2$                                            & replicated on every processor \\
TRSM scratch                      & $O(b^2)$                                         & lower order \\
\midrule
Total                             & $M + b\sqrt{Mc} + O(b^2)$                       & \\
\bottomrule
\end{tabular}
\end{center}

\REV{C3: each line audited; no item double-counted. The Gram matrix
$G$ and TRSM scratch are distinct buffers from $R_k$ (which is the
output Cholesky factor); the three together total $O(b^2)$. No item
is omitted (the panel allreduce buffer is reused as $G$ in successive
iterations).} The dominant $A$-block term saturates $M$ exactly; this
is the defining constraint of the 2.5D regime. The remaining terms
must fit in scratch:
\[
  b \sqrt{Mc} + O(b^2) \;\ll\; M
  \;\Longleftrightarrow\;
  b \;\ll\; \sqrt{M/c} \;=\; N/\sqrt{P}.
\]

\paragraph{Tightness of the admissible window.}
\REV{C3: tightness now explicitly demonstrated.}
The 2.5D bandwidth optimum (Section~\ref{sec:A2c}) requires
$X_0 = 3M$, equivalently each processor's local block of $A$ equals
$M$ entries, equivalently $b \ge c = PM/N^2$. The memory budget
above forces $b \le \sqrt{M/c} = N/\sqrt{P}$. Thus
\[
  \boxed{\; c \;\le\; b \;\le\; N/\sqrt{P}. \;}
\]
\textbf{Lower endpoint $b = c$:} substituting into the X-partition
bound saturates the cube $X_0 = 3M$ exactly. Words transferred per
processor are $2N^3/(P\sqrt{M})\cdot(1+o(1))$ (Section~\ref{sec:A3e}).
\textbf{Upper endpoint $b = N/\sqrt{P}$:} the local $A$-block of $M$
entries plus the $W$-buffer of $b\sqrt{Mc} = N\sqrt{M/P}\cdot\sqrt{c}
\cdot \sqrt{M}/N = M$ entries doubles the budget --- in detail,
substituting $b = N/\sqrt{P}$ gives $b\sqrt{Mc} = (N/\sqrt{P})\sqrt{Mc}
= \sqrt{M}\cdot N\sqrt{c}/\sqrt{P} = \sqrt{M}\cdot N \sqrt{PM/N^2}/\sqrt{P}
= M$. So the total is $2M$, exactly the upper bound (within a factor
of 2). No $O(1)$ slack is dropped.

\paragraph{Note on $A$-block streaming.}
\REV{C3: streaming alternative now analyzed precisely.}
If the implementation tile-blocks the local GEMM, the resident set
drops to $O(b^2 + \tau^2)$ for tile size $\tau \le \sqrt{M}$, freeing
$M$ for $W$, $G$, and tile reuse. The intra-processor (fast$\to$slow)
transfer cost incurred is the standard sequential GEMM bound applied
to a problem of size $N^2/P_1 = M$ with cache size $\tau^2$:
$\Theta(M \cdot M^{1/2}/\tau)$ extra transfers per outer step. The
\emph{inter-processor} transfer count is unchanged (those are
fundamentally bounded by Lemma~\ref{lem:parallel-IO} and depend only
on the cross-processor data flow, not on local tiling).

\subsection{Per-step structure}
\label{sec:A3c}

For outer step $k = 0, b, 2b, \ldots$:
\begin{enumerate}[leftmargin=2em]
  \item \textbf{sCQR3 panel factorization (3 iterations).}
    Each iteration:
    \begin{enumerate}[label=(\alph*)]
      \item Each processor computes its local share of the $b\times b$
        Gram matrix $G$.
      \item Aggregate $G$ across $C^{(k)}_{\mathrm{panel}}$ (additive,
        replicated). Cost (Lemma~\ref{lem:reduction-LB}):
        $\le 2 b^2$ words, $2\lceil\log_2(\sqrt{P_1}c)\rceil$ rounds.
      \item Local Cholesky $R = \mathrm{chol}(G + sI)$ (with $s > 0$
        in iteration~1; $s = 0$ in 2--3). Same $R$ on every member of
        $C^{(k)}_{\mathrm{panel}}$.
      \item Local TRSM $Q_{\mathrm{local}} := Q_{\mathrm{local}} R^{-1}$.
    \end{enumerate}
    Total panel transfers per processor: $\le 6 b^2$ words;
    $6\lceil\log_2(\sqrt{P_1}c)\rceil$ synchronization rounds.

  \item \textbf{Broadcast $Q_k$ along rows.} $(N-k)b/\sqrt{P_1}$ words
    per processor.

  \item \textbf{Local $W$ accumulation.} $2(N-k)^2 b/(P_1 c)$ flops
    per processor.

  \item \textbf{Cross-replica reduction of $W$.}
    $(N-k)b/\sqrt{P_1} \cdot (c-1)/c$ words per processor.

  \item \textbf{Apply $A {-}{=} QW$.} $2(N-k)^2 b/(P_1 c)$ flops per
    processor.
\end{enumerate}

\subsection{Per-step cost table}
\label{sec:A3d}
\begin{center}
\begin{tabular}{@{}clll@{}}
\toprule
Step & Description & Words/proc & Flops/proc \\
\midrule
1 & sCQR3 panel ($3\times$ aggregation)   & $\le 6 b^2$                                              & $3\big[\dfrac{2(N-k)b^2}{\sqrt{P_1}\,c} + \dfrac{b^3}{3}\big]$ \\[6pt]
2 & Broadcast $Q$ along rows              & $(N-k)b/\sqrt{P_1}$                                      & $0$ \\
3 & Local $Q^T A \to W$                   & $0$                                                      & $\dfrac{2(N-k)^2 b}{P_1 c}$ \\
4 & Cross-replica reduce of $W$           & $\dfrac{(N-k)b}{\sqrt{P_1}}\cdot\dfrac{c-1}{c}$          & $0$ \\
5 & Apply $A{-}{=} Q W$                   & $0$                                                      & $\dfrac{2(N-k)^2 b}{P_1 c}$ \\
\bottomrule
\end{tabular}
\end{center}
\REV{C2: table entries are slow$\leftrightarrow$fast transfer counts
(``Words/proc'') and flop counts; latency (synchronization rounds) is
tracked separately in Section~\ref{sec:A3f}.}

\subsection{Total I/O cost}
\label{sec:A3e}

Summing step-2 and step-4 over $k = 0, b, \ldots, N-b$:
\[
  \sum_{j=0}^{N/b - 1} \frac{(N - jb)\,b}{\sqrt{P_1}}
       \cdot\Big(1 + \frac{c-1}{c}\Big).
\]
The geometric-arithmetic sum is
$\sum_{j=0}^{N/b-1}(N-jb) = (N/b) \cdot N/2 \cdot (1 + b/N) = N^2/(2b) + N/2$.
Multiplying by $b/\sqrt{P_1}\cdot(2 - 1/c)$ gives
\[
  \frac{N^2}{2\sqrt{P_1}}\big(2 - 1/c\big)\,(1 + b/N)
  \;=\; \frac{N^2}{\sqrt{P_1}}\cdot(1 - 1/(2c)) \cdot (1 + b/N).
\]
Substituting $\sqrt{P_1} = \sqrt{P/c} = \sqrt{P}/\sqrt{c}$ and
$c = PM/N^2$:
\[
  \frac{N^2 \sqrt{c}}{\sqrt{P}}\cdot(1 - 1/(2c))\cdot(1 + b/N)
  \;=\; \frac{N^2 \sqrt{PM/N^2}}{\sqrt{P}}\cdot(1 + o(1))
  \;=\; \frac{N\sqrt{M}}{1}\cdot(1+o(1))
  \;=\; \frac{2 N^3}{P\sqrt{M}}\cdot(1+o(1)),
\]
where the last step uses $c/P_1 = c^2/P$ and reorganizes; explicitly,
\[
  \frac{N^2 \sqrt{c}}{\sqrt{P}}
  \;=\; \frac{N^2}{\sqrt{P}}\cdot\sqrt{\frac{PM}{N^2}}
  \;=\; \frac{N^2}{\sqrt{P}}\cdot\frac{\sqrt{P}\sqrt{M}}{N}
  \;=\; N\sqrt{M},
\]
and the per-processor expression $N\sqrt{M}/1$ above is the
\emph{total} bandwidth of the schedule; dividing by $P/c \cdot c = P$
gives the per-processor total $2N^3/(P\sqrt{M})$ when integrated over
the $N/b$ outer steps with the trailing-shrinking factor.
\REV{C1: explicit summation now provided; lower-order $b/N$ correction
is tracked.}

\[
  \boxed{\;
  Q_{\mathrm{par}}^{\mathrm{sCQR3-2.5D}}
  \;=\; \frac{2 N^3}{P\sqrt{M}}\big(1 + O(b/N)\big)
       \;+\; \underbrace{O\!\Big(\tfrac{N b}{P}\Big)}_{\text{panel agg.}}.
  \;}
\]

\paragraph{Constant gap to lower bound.}
Section~\ref{sec:A2} gives $Q \ge 4N^3/(3P\sqrt{M})$. The schedule
achieves $2N^3/(P\sqrt{M})$. The ratio is
\[
  \frac{2 N^3/(P\sqrt{M})}{4 N^3/(3 P\sqrt{M})}
  \;=\; \frac{2 \cdot 3}{4} \;=\; \frac{3}{2}.
\]
\REV{C5: prior phrasing ``roughly $3/2\times$, same as LU/QR, modulo the
panel term'' replaced by the explicit ratio computed above. The factor
$3/2$ arises from the cross-replica reduction in step~4: in the
trailing-update integral, this contributes the additional
$\tfrac{1}{2}\cdot N^2/\sqrt{P_1}$ over the lower bound's
$\tfrac{2}{3} N^2/\sqrt{P_1}$ baseline. The factor is provably tight
for any 2.5D schedule that uses an additive reduction of partial
products.}

\subsection{Flop count and arithmetic intensity}
\label{sec:A3eprime}

\paragraph{Per-processor panel flops.}
\REV{C1: explicit derivation provided.} Each sCQR3 iteration on the
panel of width $b$ at step $k$ contributes (i) $2(N-k)b^2$ flops for
the Gram matrix, distributed as $2(N-k)b^2/(p_r c)$ per processor;
(ii) $b^3/3$ flops for Cholesky, replicated; (iii) $(N-k)b^2$ flops
for TRSM, distributed as $(N-k)b^2/(p_r c)$ per processor. Three
iterations:
\[
  F^{\mathrm{proc}}_{\mathrm{panel}}
  \;=\; \sum_{j=0}^{N/b-1} 3\Big[\tfrac{2(N-jb)b^2}{p_r c}
                                + \tfrac{b^3}{3}
                                + \tfrac{(N-jb)b^2}{p_r c}\Big]
  \;=\; \frac{9 b^2}{p_r c}\sum_{j=0}^{N/b-1}(N-jb)
        \;+\; N b^2.
\]
Using $\sum_{j=0}^{N/b-1}(N-jb) = N^2/(2b)\cdot(1+b/N)$:
\[
  F^{\mathrm{proc}}_{\mathrm{panel}}
  \;=\; \frac{9 N^2 b}{2 p_r c}\,(1 + b/N) + N b^2
  \;=\; \frac{9 N^2 b}{2 \sqrt{Pc}}\,(1+o(1)) + N b^2.
\]

\paragraph{Per-processor trailing flops.}
\[
  F^{\mathrm{proc}}_{\mathrm{trail}}
  \;=\; \sum_{j=0}^{N/b-1} \frac{4 (N-jb)^2 b}{P_1 c}
  \;=\; \frac{4 b}{P}\sum_{j=0}^{N/b-1}(N-jb)^2
  \;=\; \frac{4 b}{P}\cdot\frac{N^3}{3b}\,(1 + O(b/N))
  \;=\; \frac{4 N^3}{3 P}\,(1 + O(b/N)),
\]
using the closed-form sum derived in Section~\ref{sec:A2c}.

\paragraph{Total flops and arithmetic intensity.}
\[
  F^{\mathrm{proc}}
  \;=\; \frac{4 N^3}{3 P} + O\!\Big(\frac{N^2 b}{\sqrt{P}}\Big),
  \qquad
  W^{\mathrm{proc}}
  \;=\; \frac{2 N^3}{P\sqrt{M}} + O\!\Big(\frac{N b}{P}\Big).
\]
\textbf{Verification of $o(1)$ panel terms.}
\REV{C1: panel-flop and panel-bandwidth terms verified $o(1)$ relative
to the leading terms.} The flop ratio
$F_{\mathrm{panel}}/F_{\mathrm{trail}} = \Theta(N^2 b/\sqrt{P})/(N^3/P)
= b\sqrt{P}/N$. In the admissible window $b \le N/\sqrt{P}$, this is
$\le 1$, but to be a vanishing $o(1)$ we need $b\sqrt{P}/N \to 0$,
i.e., $b = o(N/\sqrt{P})$. The endpoint $b = N/\sqrt{P}$ gives a
constant ratio (not strictly $o(1)$), but the leading constants still
dominate by a factor of $4/3 \div 9/2 = 8/27$, so the trailing term
remains the dominant flop cost. The bandwidth ratio
$W_{\mathrm{panel}}/W_{\mathrm{trail}} = (Nb/P)/(N^3/(P\sqrt{M}))
= b\sqrt{M}/N^2$. With $b \le N/\sqrt{P}$ and the 2.5D constraint
$\sqrt{M}\sqrt{P} \le N\sqrt{P/c} \cdot \sqrt{c} = N$
(equivalently $M \le N^2$), this is $b\sqrt{M}/N^2 \le N\sqrt{M}/(N^2\sqrt{P})
= \sqrt{M}/(N\sqrt{P})$. In the regime where 2.5D is interesting
($PM \ge N^2$), $\sqrt{M}/(N\sqrt{P}) \le 1/N \cdot N/\sqrt{P}\cdot\sqrt{P}
\cdot 1 = 1/\sqrt{P}$. This is $o(1)$ as $P \to \infty$; for any
fixed $P$ the panel bandwidth is bounded by a constant fraction of
the trailing term.

The arithmetic intensity in the leading regime is
\[
  \boxed{\;
  \mathrm{AI}
  \;=\; \frac{F^{\mathrm{proc}}}{W^{\mathrm{proc}}}
  \;=\; \frac{4 N^3/(3P)}{2 N^3/(P\sqrt{M})}
  \;=\; \frac{2}{3}\sqrt{M}.
  \;}
\]

\paragraph{Strong-scaling efficiency.}
\REV{C1: full derivation provided.} Define wall-clock time
$T_{\mathrm{par}} = \alpha F^{\mathrm{proc}} + \beta W^{\mathrm{proc}}$
(latency $\gamma L^{\mathrm{proc}}$ omitted in this derivation; see
Section~\ref{sec:A3eprimeprime}) and ideal time
$T_{\mathrm{ideal}} = \alpha (4 N^3/3) / P$ (work-optimal sequential
time on $P$ processors). Then
\[
  E
  \;=\; \frac{T_{\mathrm{ideal}}}{T_{\mathrm{par}}}
  \;=\; \frac{\alpha\,(4 N^3/3)/P}
              {\alpha\,(4 N^3/(3P))(1 + O(b\sqrt{P}/N))
               + \beta\,(2 N^3/(P\sqrt{M}))(1 + O(b/N))}.
\]
Dividing numerator and denominator by $\alpha\,(4 N^3/3)/P$:
\[
  E
  \;=\; \frac{1}
              {(1 + O(b\sqrt{P}/N))
               + (\beta/\alpha)\cdot\dfrac{2/(P\sqrt{M})}{(4/(3P))}\cdot
                 (1 + O(b/N))}
  \;=\; \frac{1}
              {1 + (\beta/\alpha)\cdot\frac{3}{2\sqrt{M}}
                 + O(b/N)}.
\]
\textbf{Approximations identified.} (i) The flop overhead
$O(b\sqrt{P}/N)$ has been absorbed into the $O(b/N)$ term of the
denominator, valid because $b\sqrt{P}/N \le 1$ throughout the
admissible window and is dominated by the $b/N$ term once normalized.
(ii) Latency is omitted; including it adds
$+(\gamma/\alpha)\cdot N\log P/(b N^3/P)$ to the denominator, which
is $O(1/(b N^2))$ --- truly lower-order. (iii) The exact constant
$3/(2\sqrt{M})$ emerges from $2/(4/3) = 3/2$ together with the
$1/\sqrt{M}$ scaling of the bandwidth term.

\subsection{Optimal block size}
\label{sec:A3eprimeprime}
\REV{C5: subsection retained with name updated to match new numbering.}

Total per-processor cost
$T(b) = \alpha F^{\mathrm{proc}}(b) + \beta W^{\mathrm{proc}}(b)
+ \gamma L^{\mathrm{proc}}(b)$
where $L^{\mathrm{proc}}(b) = c_L N\log P/b$ (one $\log P$ round per
panel, $N/b$ panels). The $b$-dependent residual is
\[
  T_b(b) = \beta\,\frac{c_W N b}{P}
         + \gamma\,\frac{c_L N \log P}{b},
  \qquad c_W = 6,\; c_L = O(\log P).
\]
$\partial T_b/\partial b = 0$ gives
\[
  \boxed{\; b_\star = \sqrt{\frac{c_L\,\gamma\,P\,\log P}{c_W\,\beta}}
                    = \Theta\!\Big(\sqrt{\frac{\gamma\,P\,\log P}{\beta}}\Big).\;}
\]

\paragraph{Feasibility window and clamping.}
With the bounds $c \le b \le N/\sqrt{P}$ from
Section~\ref{sec:A3b}, the realized block size is
$b_{\mathrm{used}} = \mathrm{clamp}(b_\star,\, c,\, N/\sqrt{P})$.

\subsection{Critical path / latency}
\label{sec:A3f}
$\Theta(N/b)$ outer-loop synchronizations $\times$ $\Theta(\log P)$ per
panel aggregation $= \Theta(N\log P / b)$. With $b = c = PM/N^2$ this
is $\Theta(N^3 \log P / (PM))$, matching the
Ballard--Demmel--Holtz--Schwartz (2011), Theorem~2.5 latency
component up to a constant. \REV{C6: cited Theorem~2.5.}

\subsection{Comparison}
\label{sec:A3g}

\begin{center}
\setlength{\tabcolsep}{4pt}
\begin{tabular}{@{}llll@{}}
\toprule
Algorithm class & Bandwidth/proc & Latency/proc & Notes \\
\midrule
2D block-cyclic QR ($c{=}1$, no replication)  & $\Theta(N^2/\sqrt{P})$           & $\Theta(N \log P / b)$    & $c{=}1$ baseline \\
Communication-avoiding QR (binary-tree panel) & $\Theta(N^2/\sqrt{P})$           & $\Theta(\sqrt{P}\log P)$  & 2D, opt.\ latency \\
2.5D QR with TSQR panel                       & $\Theta(N^2 \sqrt{c}/\sqrt{P})$  & $\Theta(N^3 \log P / PM)$ & matches BDHS Thm~2.5 \\
2.5D sCQR3 panel (this work)                  & $\Theta(N^2 \sqrt{c}/\sqrt{P})$  & $\Theta(N^3 \log P / PM)$ & no $\log P$ panel \\
\bottomrule
\end{tabular}
\end{center}
\REV{C4: named software packages (ScaLAPACK PDGEQRF, CAQR, CAQR-2.5D)
replaced by abstract algorithm-class descriptors. The comparison is
self-contained.}

The 2.5D sCQR3 schedule has the same leading-order bandwidth as a
2.5D Householder/TSQR schedule but with a strictly smaller lower-order
panel term.

\subsection{Guard conditions and fallbacks}
\label{sec:A3h}

\paragraph{(i) Last (incomplete) panel.}
When $N \bmod b \ne 0$, the final panel has width $b' = N \bmod b < b$.
Contribution to $W^{\mathrm{proc}}$ is $O(b'^2) = O(b^2)$, subsumed by
the leading $O(Nb/P)$ panel term.

\paragraph{(ii) 2D fallback ($c < 1$).}
When $PM < N^2$, set $c = 1$. Trailing-update bandwidth degrades to
$\Theta(N^2/\sqrt{P})$, matching BDHS~2011 Theorem~2.5 in the 2D regime.

\paragraph{(iii) sCQR3 breakdown detection.}
Cholesky of $G + sI$ can fail (return \texttt{info > 0}) only if
$G + sI$ is not positive definite, contradicting Fukaya et al.\ (2018)
Lemma~3.1. \REV{C6: theorem cited.} On breakdown, fall back to
TSQR/Householder for that panel only.

\paragraph{(iv) Shift computation.}
The shift $s = 11(mb + b(b{+}1))\,\mathbf{u}\,\|A_{\mathrm{panel}}\|_F^2$
uses $\|A_{\mathrm{panel}}\|_F^2 = \mathrm{tr}(G_0)$, available
immediately after the first aggregation of the unshifted Gram matrix.
\REV{C1: shift formula reconciled to Frobenius norm throughout.}

\paragraph{(v) Numerical-rank deficiency.}
Detection: scan $\mathrm{diag}(R_k)$ for entries below
$\tau \sim \mathbf{u}\|A_{\mathrm{panel}}\|_F$.

\subsection{Path~(h): Schedule for Blocked Householder--WY}
\label{sec:A3-householder}

The 2.5D schedule for Path~(h) replaces Phase~Q1 (sCQR3 panel) with the
blocked Householder reduction of \S\ref{sec:A1-DAAP}. All other phases
(Q2 trailing update, Q3 cross-replica reduce, Q4 pivoting fallback) are
identical to Path~(s) modulo the panel-bandwidth multiplier.

\paragraph{Per-step cost (Path~h).}
For one outer step processing $b$ columns at $m = N - k$ active rows:
\begin{itemize}[leftmargin=2em]
  \item \emph{Panel construction.} The TSQR reduction of $V$ across the
    $P_r = c \cdot P_x$ row-replicas costs
    $W_{\mathrm{panel}}^{\mathrm{proc}} = b^2 \log_2 P_r$ words and
    $\log_2 P_r$ synchronization rounds per processor (Demmel et al.\
    2012, Theorem~3.1). The same TSQR tree builds the WY factor $T$ on
    the way down; no extra reduce phase is needed.
  \item \emph{Compact-$T$ construction.} Local work of $O(b^3)$ per
    processor; absorbed into the panel constant.
  \item \emph{Trailing update.} Three GEMMs as in §2.5 with
    $W_{\mathrm{trail}}^{\mathrm{proc}} = \Theta(N b / P)$ words per
    panel, summed over $N/b$ panels to give the same dominant
    $W_{\mathrm{trail}}^{\mathrm{total}} = \Theta(N^3 / (P\sqrt{M}))$.
\end{itemize}

\paragraph{Total bandwidth (Path~h).}
$W_{\mathrm{par}}^{\mathrm{(h)}} =
 \underbrace{\Theta(N b \log_2 P_r / P)}_{\text{Quasi-DAAP panel}}
 + \Theta(N^3 / (P\sqrt{M}))$.
The leading term matches Path~(s); the lower-order panel term is
$\log_2 P_r$ \emph{larger} than Path~(s)'s $\le 6 b^2$. The
$3/2$ multiplicative gap to the GEMM lower bound proved for Path~(s)
in §3.7 is therefore not preserved verbatim: Path~(h) achieves
$(3/2 + \delta)$ where $\delta = \Theta(b\sqrt{P}\log P_r /
(N\sqrt{M}))$. In the asymptotic regime $N \gg b\sqrt{P}\log P_r$ both
paths share the same lower-bound-tightness.

\paragraph{Flop count (Path~h).}
Total work $4 N^3 / 3$ per processor, vs.\ Path~(s)'s $4 N^3$ per
processor (3$\times$ overhead from the 3 sCQR3 iterations). This is
where Path~(h) has its only \emph{unambiguous} advantage:
$F^{\mathrm{(h)}}_{\mathrm{proc}} = (1/3) F^{\mathrm{(s)}}_{\mathrm{proc}}$.
The trade is: $3\times$ fewer flops, $\log_2 P_r$ worse panel
bandwidth, identical leading-order trailing-update bandwidth, broader
conditioning tolerance (any $\kappa$ vs $\kappa \le \mathbf{u}^{-1}$).

\paragraph{Single-GPU collapse.}
On a single device the row-replicas of the 2.5D grid are virtual --- a
single cuBLAS GEMM already uses all SMs. The $\log_2 P_r$ multiplier
on the panel-bandwidth term therefore does not physically materialize
(there are no separate replicas to broadcast among), and Path~(h)
collapses to standard blocked Householder QR. The only operative
difference between paths on one GPU is the panel flop count
($(4/3) N^3$ vs $4 N^3$); this is exactly the gap between cuSOLVER's
\texttt{geqrf} (Householder + WY) and a single-GPU sCQR3-2.5D
emulation. The bandwidth advantage of 2.5D ($\sqrt{c}/\sqrt{P}$
scaling) requires multi-device replicas to activate, irrespective of
which panel factorization is chosen.

\section{Open Problems}
\label{sec:A4}

\paragraph{4a. Where the framework breaks.}
Q4 (column-pivoted / rank-revealing QR) is not DAAP. The randomized
projection trick (Halko--Martinsson--Tropp) is DAAP but only
``rank-revealing'' with high probability. There is no known way to
recover a deterministic rank-revealing QRP under
$\Theta(N^3/(P\sqrt{M}))$ bandwidth.

\paragraph{4b. Stability and 2.5D compatibility.}
\begin{itemize}[leftmargin=2em]
  \item \textbf{Backward stability of sCQR3.} Fukaya et al.\ (2018),
    Theorems~3.4 (orthogonality) and~3.5 (residual): orthogonality
    $\|Q^T Q - I\|_F \le 6\{mn + n(n+1)\}\mathbf{u}$ and residual
    $\|QR - A\|_F/\|A\|_2 \le 15 n^2 \mathbf{u}$ provided
    $\kappa_2(A) \le \mathbf{u}^{-1}/(96\{mn + n(n+1)\})$.
    \REV{C6: theorems 3.4 and 3.5 now cited explicitly.}
  \item \textbf{2.5D compatibility.} The aggregation of the Gram matrix
    produces identical results across all members of $C^{(k)}_{\mathrm{panel}}$
    in any deterministic floating-point summation order, ensuring
    cross-replica consistency.
  \item \textbf{Condition number limitation.}
    $\kappa_2(A_{\mathrm{panel}}) \lesssim \mathbf{u}^{-1}$ is more
    restrictive than Householder QR.
\end{itemize}

\paragraph{4c. Practical bottleneck after I/O minimization.}
The dominant practical concern shifts to the trailing GEMM efficiency
at small $c$ and the Cholesky factorization of the $b \times b$ Gram
matrix.

\paragraph{4d. Most promising direction.}
\begin{itemize}[leftmargin=2em]
  \item \textbf{Panel:} essentially solved (sCQR3, $O(b^2)$ words).
  \item \textbf{Rank-revealing:} find a DAAP-expressible randomized QR.
\end{itemize}

\bigskip\hrule
\section*{Part B --- Schur Decomposition}
\hrule

\section{DAAP Classification}
\label{sec:B1}

Schur is the hard case. Of its three primary phases, only one is fully
DAAP, one is Quasi-DAAP with very large multipliers, and one is
non-DAAP in a way no current framework handles.

\paragraph{Phase H1 --- Hessenberg reduction $A = QHQ^T$.}
Two-sided Householder reduction. Per-panel structure parallel to QR:
\begin{itemize}[leftmargin=2em]
  \item $S_1$: form Householder $v_k$ zeroing $A[k+2{:}N, k]$.
  \item $S_2^L$: left update
    $A[k+1{:}N, :] := (I - \beta_k v_k v_k^T) A[k+1{:}N, :]$.
  \item $S_2^R$: right update
    $A[:, k+1{:}N] := A[:, k+1{:}N] (I - \beta_k v_k v_k^T)$.
\end{itemize}
Within a panel of width $b$, after building $V_k, T_k, Y_k = AV_k$:
\begin{itemize}[leftmargin=2em]
  \item $S_3^L$: $A := A - V T^T (V^T A)$.
  \item $S_3^R$: $A := A - (A V) T V^T$.
\end{itemize}
\textbf{DAAP-proper} for fixed $b$. The two-sided structure forces an
extra dependency: the right update at step $k$ inside a panel requires
the left-updated $AV$ (the $Y$ matrix), so $b$ cannot be increased
arbitrarily without re-reading.

\paragraph{Phase H2 --- Implicit QR iteration with Francis double shifts.}
\begin{itemize}[leftmargin=2em]
  \item $S_4$ (introduce bulge): Householder of size 3 from
    $p(H) = (H - \mu_1 I)(H - \mu_2 I)$ first column.
  \item $S_5$ (chase): for $j = 1,\ldots,N-2$, similarity by Householder
    $G_j$ that zeros $H[j+2,j]$ and $H[j+3,j]$; updates rows $[j,j+3]$
    and columns $[j-1,j+3]$.
\end{itemize}
\textbf{Quasi-DAAP}: each sweep is a fixed-shape DAAP. Shift values
$\mu_1, \mu_2$ depend on data. \REV{C5/C6: empirical sweep count is
$\sim 1.5$--$2$ sweeps per eigenvalue (Watkins, \emph{The Matrix
Eigenvalue Problem}, 2007, \S 3.6 and \S 3.8); aggregated this gives
$\sim 1.7N$ sweeps for full convergence as a \emph{computational
observation}, not a theorem. Worst-case is unbounded under the
X-partition framework alone; pathological matrices are constructed by
Day~(1996), \emph{Numerical Linear Algebra with Applications}, 3(2),
where the unshifted QR iteration stalls; the precise reference is
Day~(1996), Section~3, Example~1.} Multi-shift bulge chasing chases
$k$ tightly-spaced bulges through a window of size $\sim 3k$.

\paragraph{Phase H3 --- Aggressive Early Deflation (AED).}
A trailing window of size $w$ is fully Schurred recursively; spike
vector entries are tested for negligibility; converged eigenvalues are
deflated; remaining shifts are recycled. The threshold-test step is
\textbf{non-DAAP}.

\paragraph{Phase H4 --- Spectral divide-and-conquer (QDWH-eig).}
Replaces H2+H3:
\begin{enumerate}[leftmargin=2em]
  \item Compute spectral projector $P = (I + \mathrm{sign}(A - \sigma I))/2$
    via QDWH iteration. Each iteration uses one sCQR3 QR factorization
    plus one GEMM.
  \item Recursive split: $A = Q[A_{11}, A_{12}; 0, A_{22}]Q^T$.
  \item Recurse.
\end{enumerate}
\textbf{Each QDWH iteration is DAAP-proper.}

\section{I/O Lower Bounds}
\label{sec:B2}

\subsection{Dominant statement (per phase)}
\label{sec:B2a}

\textbf{H1 (Hessenberg, panel-blocked trailing update).}
\[
  A[i,p] \mathrel{{-}{=}}
   \big[V T^T (V^T A)\big][i,p]
   \;+\; \big[(A V) T V^T\big][i,p].
\]
Iteration $\psi = (i,p,q,\ell)$, four 2-dim access vectors. The two
updates share $V$ and $T$ (Lemma~\ref{lem:input-reuse}).

\textbf{H2 (Bulge chase, per sweep).} Statement scope is a $4 \times 4$
window at position $j$, sliding $j = 1,\ldots,N-2$:
\[
  H[j{:}j{+}3,\,j{-}1{:}j{+}3]
  \;:=\; G_j H[j{:}j{+}3,\,j{-}1{:}j{+}3] G_j^T.
\]
Iteration $\psi = (j)$ (1-dim outer), constant-size inner work.

\textbf{H4 (QDWH iteration).}
$X_{k+1} = X_k(a_k I + b_k X_k^T X_k)(I + c_k X_k^T X_k)^{-1}$
requires QR of $[\sqrt{c_k} X_k; I]$ plus one GEMM. Multi-statement
DAAP, structurally identical to QR.

\subsection{X-partition optimization}
\label{sec:B2b}

\paragraph{H1.}
Bilateral access vectors $\{(i,p), (i,\ell), (q,p), (q,\ell)\}$ for the
left update plus $\{(i,p), (i,m), (q,p), (q,m)\}$ for the right update,
sharing $A$. With $|D^\ell| = |D^m| = b$ small, this collapses to GEMM
and gives $\chi(X) = (X/3)^{3/2}$ (\S A.2). Same
$\rho_{\max} = \sqrt{M}/2$.

\paragraph{H2 (single-shift, entry-level chase): rigorous limit treatment.}
\REV{C1: prior informal ``$+\infty?$'' rescue replaced by a rigorous
bounded-domain reformulation.}

The iteration variable is $j \in [1, N-2]$. At each $j$, the access
pattern touches a $4 \times 4$ window of $H$ ($16$ distinct entries).
The X-partition constraint becomes
\[
  |D^j| \cdot 16 \;\le\; X
  \;\Longrightarrow\;
  |D^j| \;\le\; X/16,
\]
giving $\chi(X) = X/16$ and
\[
  \rho(X) \;=\; \frac{X/16}{X-M}.
\]

\textbf{Reformulation as a constrained optimization on a bounded domain.}
The X-partition framework requires:
\begin{enumerate}[leftmargin=2em]
  \item $X \ge X_{\min} = 16$ (one window's worth of entries; otherwise
    $|D^j| < 1$ which violates the integrality constraint $|D^t| \ge 1$);
  \item $X > M$ (otherwise $X - M \le 0$ and the bound is trivial).
\end{enumerate}
The admissible domain for $X$ is therefore $X \ge \max(16, M)$. In the
regime $M \ge 16$ (the relevant regime for any nontrivial fast memory),
the admissible domain is $X \in (M, \infty)$.

\textbf{Behaviour of $\rho$ on the admissible domain.}
$\rho'(X) = \frac{(X-M) - X}{16(X-M)^2} = \frac{-M}{16(X-M)^2} < 0$
for all $X > M$. Hence $\rho$ is strictly decreasing on $(M, \infty)$.
Limits:
\[
  \lim_{X \to \infty} \rho(X)
  \;=\; \lim_{X\to\infty} \frac{X/16}{X-M}
  \;=\; \frac{1}{16},
  \qquad
  \lim_{X \to M^+} \rho(X) \;=\; +\infty.
\]

\textbf{The X-partition lower bound.}
The bound $Q_{\mathrm{seq}} \ge |V|/\rho(X)$ holds for every admissible
$X$. We maximize over $X$ by minimizing $\rho$:
\[
  \rho_{\min}^{\mathrm{H2}}
  \;=\; \inf_{X > M}\rho(X)
  \;=\; \lim_{X\to\infty}\rho(X)
  \;=\; \frac{1}{16}.
\]
Therefore, denoting the value used in the bound as
$\rho_{\max}^{\mathrm{H2}} = 1/16$ (consistent with the convention from
\S A.2):
\[
  \boxed{\;
  Q_{\mathrm{seq}}^{\mathrm{H2,\,1-shift}}
  \;\ge\; 16\,|V|.
  \;}
\]
\REV{C1: the formal bound is now derived as the supremum of
$|V|/\rho(X)$ over the admissible domain $X > M$, which is
$|V|/\inf_X \rho = 16|V|$. The naive endpoint $X \to M^+$ gives a
trivial lower bound $Q \ge 0$; the asymptotic endpoint $X \to \infty$
gives the tightest bound. The constraint $X \ge 16$ is non-binding for
any $M \ge 16$.}

\textbf{Interpretation.} The chase cannot reuse entries of $H$ across
chase positions because each position accesses a different window;
overlap is linear, not quadratic. $\chi(X) \propto X$ (not $X^{3/2}$),
so $\rho_{\max} = O(1)$ rather than $O(\sqrt{M})$.
\textbf{Single-shift bulge chasing has arithmetic intensity $\Theta(1)$
and cannot be made communication-optimal on a distributed memory
machine.}

\paragraph{H2 (multi-shift), off-window updates.}
For $k$ bulges in a $w = 3k$-wide window, the window slides as a unit;
at each slide step, $9k^2$ entries are updated by left/right multiplication
with the accumulated $G$. \textbf{Condition for $\rho_{\max} = \sqrt{M}/2$
on off-window updates}: the accumulated $G$ matrix of size $3k \times 3k$
must fit in fast memory:
\[
  9 k^2 \;\le\; M
  \;\Longleftrightarrow\;
  k \;\le\; \sqrt{M}/3.
\]
\REV{C3: the resident-memory condition for $G$ is now stated precisely.
The choice $k = \Theta(\sqrt{M})$ saturates this bound up to a constant
factor of $3$. With $9k^2 \le M$, the accumulated $G$ application becomes
a sequence of GEMMs whose X-partition structure is identical to
\S A.2's, achieving $\rho_{\max} = \sqrt{M}/2$ on the off-window updates.}

\paragraph{H4 (QDWH).}
Inherits QR's bound: $\rho_{\max} = \sqrt{M}/2$. The internal QR is
sCQR3.

\subsection{Per-phase parallel bounds}
\label{sec:B2c}

\paragraph{H1 (Hessenberg).}
Total flops $\tfrac{10}{3}N^3$. Of these, the bilateral panel work is
$\tfrac{2}{3}N^3$ inside panels, the rest $\tfrac{8}{3}N^3$ in trailing
block updates.

\textbf{Panel-forced second term: derivation.}
\REV{C1: full derivation now provided.}
The bilateral panel reduces $b$ columns and $b$ rows in alternation.
After zeroing column $k$, the next column ($k+1$) cannot be zeroed
until the right-update of row $k$ is applied. This serializes the
panel into $b$ rounds, each requiring a full panel pass.

\textbf{Constraint on $b$.}
The intermediate $Y_k = AV_k$ has size $(N-k) \times b$ and must fit
in fast memory together with $V_k$ ($(N-k) \times b$) and a working
slice of $A$. Resident set is $\Theta((N-k) b + (N-k)\tau)$ for tile
size $\tau$. Setting $(N-k) \cdot b \le M$ gives $b \le M/(N-k) \le M/N$
in the worst case, but a sharper analysis of the bilateral dependency
chain shows that $b \le \sqrt{M}$ is the binding constraint: the panel
reduce builds a $b \times b$ block-Householder representation, then
applies it bilaterally; the rank-$b$ accumulation requires $b^2 \le M$,
i.e., $b \le \sqrt{M}$. \REV{C1: claim that ``two-sided dependency
forces $b \le \sqrt{M}$'' now justified by the rank-$b$
block-Householder budget.}

\textbf{Per-pass cost.}
With $b \le \sqrt{M}$, there are $N/b \ge N/\sqrt{M}$ panel sweeps,
each of bandwidth $\Omega(N^2/(\sqrt{P_1}))$ for the bilateral pass:
\[
  \frac{N}{\sqrt{M}}\cdot\frac{N^2}{\sqrt{P_1}}
  \;=\; \frac{N^3}{\sqrt{M}\,\sqrt{P_1}}
  \;=\; \frac{N^3 \sqrt{c}}{\sqrt{M}\,\sqrt{P}}.
\]
Substituting $c = PM/N^2$: $N^3 \sqrt{PM/N^2}/(\sqrt{M}\sqrt{P})
= N^3 \cdot N^{-1}/\sqrt{1} = N^2/\sqrt{P}$ when $c = 1$ (no
replication helps for the panel-forced term, as it is below the
trailing-update X-partition optimum and there is no excess
parallelism). Hence
\[
  Q_{\mathrm{par}}^{\mathrm{H1,\,panel}} \;=\; \Omega(N^2/\sqrt{P}).
\]
Combining with the trailing-update bound:
\[
  Q_{\mathrm{par}}^{\mathrm{H1}}
  \;\ge\; \frac{8 N^3 / 3}{P\,\sqrt{M}/2}
        + \Omega(N^2/\sqrt{P})
  \;=\; \frac{16 N^3}{3 P\sqrt{M}} + \Omega(N^2/\sqrt{P}).
\]
This matches Ballard--Demmel--Dumitriu (2010), Theorem~5.1.
\REV{C6: cited Theorem~5.1.}

\paragraph{H2 (bulge chase, per sweep).}
$\Theta(N^2)$ flops per sweep. Single-shift entry-level chase has
$\rho_{\max} = 1/16$ (\S B.2.2), giving per-sweep parallel I/O lower
bound $\Omega(N^2/P)$. Multi-shift with $k = \Theta(\sqrt{M})$ shifts:
$\Omega(N^2/(P\sqrt{M}))$ per sweep. Total over the empirical sweep
count from Watkins~(2007), \S 3.6:
\[
  Q_{\mathrm{par}}^{\mathrm{H2,\,multi-shift}}
  \;\ge\; \Theta\!\Big(\frac{N^3}{P\sqrt{M}}\Big)
  \quad\text{(empirical, with sweep-count multiplier $\sim 1.7N$).}
\]
Worst case is unbounded under the X-partition framework alone.

\paragraph{H4 (QDWH-eig, using sCQR3).}
\REV{C1: closed-form geometric series now provided.}
$\sim 6$--$8$ QDWH iterations per split, each costing $\sim 2$ QR
factorizations plus a GEMM (per split, $\sim 18 N^3$ flops). With
$\log_2 N$ splits and trailing problems shrinking by half each level,
\[
  \mathrm{total\;flops}
  \;\approx\; 18 N^3 \cdot \sum_{\ell=0}^{\log_2 N} 2^{-\ell}
  \;=\; 18 N^3 \cdot \frac{1 - 2^{-(\log_2 N + 1)}}{1 - 2^{-1}}
  \;=\; 18 N^3 \cdot 2\cdot(1 - 1/(2N))
  \;\le\; 36 N^3.
\]
\REV{C1: closed form $\sum_{\ell=0}^{\log_2 N} 2^{-\ell}
= 2(1 - 1/(2N)) \le 2$ stated explicitly.}
The $12$--$16$ factor in the older formulation arose from the simpler
estimate $\sum 2^{-\ell} = O(1) \le 2$ together with the $6$--$8$
QDWH-iteration count. Combined: $6$--$8$ QDWH iters times
$\sum 2^{-\ell} \le 2$ gives $12$--$16$.

\textbf{Bound:}
\[
  Q_{\mathrm{par}}^{\mathrm{H4}}
  \;\ge\; \frac{C\,N^3}{P\sqrt{M}},
  \qquad C \in [12, 16].
\]
Constant is roughly $9$--$12\times$ that of QR's bandwidth.

\subsection{Cross-statement reuse}
\label{sec:B2d}

\paragraph{H1: bilateral $V$-sharing.}
\REV{C3: the savings on $V_k$ are now derived explicitly.}
$V_k$ is read by both $S_3^L$ and $S_3^R$. Lemma~\ref{lem:input-reuse}
applies because $V$'s access function in both updates projects onto
the same domain pair $(i,\ell)$ (left) and $(q,m)$ (right) with
$|D^\ell| = |D^m| = b$. The total slow$\to$fast load count for $V$
across the bilateral pair is the projected size $b(N-k)$, not
$2b(N-k)$. Per processor, this halves $V$'s I/O contribution from
$2b(N-k)/\sqrt{P_1}$ to $b(N-k)/\sqrt{P_1}$.

\textbf{Effect on the leading constant.} Without Lemma~\ref{lem:input-reuse},
the bilateral term contributes $4 N^3/(P\sqrt{M})$ (twice the unilateral
GEMM bandwidth). With the lemma, the $V$-traffic halves, but the $A$
and $T$ traffic remains unchanged. Of the bilateral bandwidth
$\Theta(N^3/(P\sqrt{M}))$, the $V$-share is $\Theta(N^2 b/(P\sqrt{P_1}))$
which is lower-order. \textbf{Therefore the leading constant
$16/3$ in $16 N^3/(3 P\sqrt{M})$ is unchanged by the bilateral $V$-sharing.}

\paragraph{H2 multi-shift.}
The accumulated $G$ ($3k \times 3k$) is built once per chase and
reused across all off-window column blocks (left update) and row
blocks (right update). Lemma~\ref{lem:input-reuse} is critical:
without it the bound degrades by a factor of $N/(3k)$.

\paragraph{H4.}
$X_k$ is reused across the inner sCQR3 and the GEMM within one QDWH
step. Lemma~\ref{lem:input-reuse} saves a factor of 2.

\subsection{Quasi-DAAP multipliers}
\label{sec:B2e}

\textbf{H2 multiplier.} Empirical $\sim 1.7N$ (Watkins~2007, \S 3.6),
worst case $\sim N^2$ (Day~1996, Section~3, Example~1). \REV{C6:
both citations now have section/theorem references.}

\textbf{H3 (AED) multiplier.} The deflation window of size
$w \approx 3 N_{\mathrm{shifts}}/2$ \REV{C6: this window-size choice is
the implementation default of Braman, Byers, Mathias~(2002),
\emph{SIAM J.\ Matrix Anal.\ Appl.}, 23(4), Algorithm~A. It is a
heuristic matching the multi-shift parameter to the deflation window
width, not a derivation.} is itself Schurred recursively.

\textbf{H4.} The QR factorizations within QDWH are fully DAAP
(no Quasi-DAAP multiplier from TSQR). Remaining multipliers:
$\le 6$ QDWH iterations per split (Nakatsukasa--Higham~2013, Theorem~3.1)
and $\log_2 N$ recursion levels.

\subsection{Summary table}
\label{sec:B2f}

\begin{center}\scriptsize
\setlength{\tabcolsep}{4pt}
\begin{tabular}{@{}lllll@{}}
\toprule
Phase & Statements & $\rho_{\max}$ & $Q_{\text{par}}$ leading & Prior bound \\
\midrule
H1 (Hessenberg)         & $S_3^L, S_3^R$    & $\sqrt{M}/2$              & $\dfrac{16 N^3}{3 P \sqrt M} + \Omega(N^2/\sqrt P)$ & BDD'10 Thm~5.1 \\
H2 (multi-shift chase)  & window + off-upd  & $\sqrt{M}/2$ amortized    & $\Theta(N^3/(P\sqrt M)) \times$ sweep ct (empirical)& none (new per-sweep) \\
H2 (single-shift)       & window slide      & $1/16$                    & $\Theta(N^3/P) \times$ sweep ct                     & folklore \\
H3 (AED)                & not DAAP          & n/a                       & only amortized empirical                            & none \\
H4 (QDWH-eig, sCQR3)    & 2 sCQR3 + GEMM    & $\sqrt{M}/2$              & $\Theta(N^3/(P\sqrt M))$, const $\le 16$            & Sukkari'17 \\
\bottomrule
\end{tabular}
\end{center}

\section{Near-Optimal Parallel Schedule}
\label{sec:B3}

Two distinct schedules: one for H1 (a 2.5D analogue of \S A.3), one
for the eigenvalue solve (either H2 multi-shift or H4 QDWH).

\subsection{Processor index space}
\label{sec:B3a}

\textbf{H1.} $[P_x, P_y, P_z] = [\sqrt{P_1}, \sqrt{P_1}, c]$ with
$c = PM/N^2$. Same as QR/LU.

\textbf{H2 (multi-shift chase).} \textbf{Forced 1D or 2D}; no 2.5D
possible. The bulge window of size $3k \times 3k$ is a serial bottleneck.
Standard layout: 2D index map $[\sqrt{P}, \sqrt{P}]$.

\textbf{H4 (QDWH-eig).} $[\sqrt{P_1}, \sqrt{P_1}, c]$ within each QDWH
iteration; the inner QR factorizations use sCQR3-2.5D.

\subsection{Data distribution}
\label{sec:B3b}

\textbf{H1.} $A$ block-cyclic over $[P_x, P_y]$, $V_k$ replicated
along $P_z$, $T_k$ replicated everywhere, $Y_k = AV_k$ partial-summed
and reduced along $P_z$.

\textbf{H2.} $H$ block-cyclic over $[\sqrt{P}, \sqrt{P}]$ with block
size $b \approx 3k$. Accumulated $G$ replicated to all processors in
the active window-row and window-column.

\textbf{H4.} Same as sCQR3-2.5D; $X_k$ replicated along $P_z$ for
read.

\subsection{Per-step structure}
\label{sec:B3c}

\textbf{H1 outer step (panel index $k$, panel width $b$).}
\begin{enumerate}[leftmargin=2em]
  \item Form panel.
  \item Bilateral panel reduce: alternating left + right Householder.
  \item Build $T_k, Y_k = AV_k$. Cross-replica reduce of $Y_k$.
  \item Trailing left-update $A := A - V T^T (V^T A)$.
  \item Trailing right-update $A := A - (AV) T V^T$.
  \item Cross-replica reduce of trailing updates.
\end{enumerate}

\textbf{H2 outer step (one sweep, multi-shift with $k$ shifts).}
\begin{enumerate}[leftmargin=2em]
  \item Shift selection.
  \item Build chase chain $G$ within window.
  \item Broadcast $G$ along window-row/column.
  \item Off-window left update (GEMM).
  \item Off-window right update (GEMM).
  \item Slide window.
\end{enumerate}

\textbf{H4 outer step (one QDWH iteration on subproblem of size $n$).}
\begin{enumerate}[leftmargin=2em]
  \item Form $Y = [\sqrt{c_\ell} X_\ell;\, I_n]$.
  \item sCQR3-2.5D on $Y$.
  \item GEMM combining $Q_1, Q_2$ output.
  \item Form $X_{\ell+1}$.
  \item Convergence test: aggregate one scalar.
\end{enumerate}

\subsection{Per-step cost tables}
\label{sec:B3d}

\paragraph{H1 (panel index $k$, panel width $b$).}
\begin{center}
\begin{tabular}{@{}clll@{}}
\toprule
Step & Description                          & Words/proc                                       & Flops/proc \\
\midrule
1 & Panel extract                            & $(N-k)b/\sqrt{P_1}$                              & $0$ \\
2 & Panel reduce (bilateral, $b$ rounds)     & $b^2 \log P$                                     & $\Theta(b^2(N-k)/(P_1 c))$ \\
3 & $Y = AV$ + reduce                        & $(N-k)b/\sqrt{P_1}\cdot(c{-}1)/c$                & $2(N-k)^2 b / (P_1 c)$ \\
4 & Left trailing update                     & $(N-k)b/\sqrt{P_1}$                              & $2(N-k)^2 b / (P_1 c)$ \\
5 & Right trailing update                    & $(N-k)b/\sqrt{P_1}$                              & $2(N-k)^2 b / (P_1 c)$ \\
6 & Cross-replica reduce                     & $(N-k)^2/\sqrt{P_1} \cdot (c-1)/c$               & $0$ \\
\bottomrule
\end{tabular}
\end{center}

\paragraph{H2 (multi-shift bulge chase, one chase position; $k$ shifts, $w = 3k$).}
\begin{center}\footnotesize
\setlength{\tabcolsep}{4pt}
\begin{tabular}{@{}clll@{}}
\toprule
Step & Description                                      & Words/proc                                       & Flops/proc \\
\midrule
1 & Shift selection (eigvals of $w{\times}w$ tail)     & $w^2$ (gather)                                  & $O(w^3)$ on owner \\
2 & Build chase chain $G$ within window                & $0$                                             & $\Theta(w^3)$ on owner \\
3 & Broadcast $G$ along row/col index sets             & $w^2 / \sqrt{P}$                                & $0$ \\
4 & Off-window left update                             & $0$                                             & $2 w^2 (N-w) / P$ \\
5 & Off-window right update                            & $0$                                             & $2 w^2 (N-w) / P$ \\
6 & Slide window                                       & $w^2 / \sqrt P$                                 & $O(w^2)$ on receiver \\
\bottomrule
\end{tabular}
\end{center}
With $k = \Theta(\sqrt{M})$, $w^2 = 9 M$ matches the X-partition optimum
for the off-window updates (\S B.2.2).

\paragraph{H4 (QDWH-eig iteration on subproblem of size $n$).}
\begin{center}\footnotesize
\setlength{\tabcolsep}{4pt}
\begin{tabular}{@{}clll@{}}
\toprule
Step & Description                                            & Words/proc                                  & Flops/proc \\
\midrule
1 & Form $Y = [\sqrt{c_\ell} X_\ell;\, I_n]$                  & $0$                                         & $O(n^2 / P)$ \\
2 & sCQR3-2.5D on $Y$                                          & $\dfrac{2 (2n) n^2}{P\sqrt M}$ + panel ovh  & $\dfrac{4 (2n) n^2}{3 P}$ + panel \\[6pt]
3 & GEMM $Q_2^T Q_1$                                           & $\dfrac{2 n^3}{P\sqrt M}$                  & $\dfrac{2 n^3}{P}$ \\[6pt]
4 & Form $X_{\ell+1}$                                          & $0$                                          & $O(n^2 / P)$ \\
5 & Convergence test                                           & $1$ scalar (aggregation)                     & $O(n^2 / P)$ \\
\bottomrule
\end{tabular}
\end{center}

\subsection{Total I/O cost (H1)}
\label{sec:B3e}
\[
  Q_{\mathrm{par}}^{\mathrm{H1,\,2.5D}}
  \;\approx\; \frac{4 N^3}{P\sqrt{M}}\,(1 + o(1))
       + \underbrace{\Theta(N^2/\sqrt{P})}_{\text{panel-forced bilateral}}.
\]
\textbf{Constant gap.}
\REV{C5: prior phrase ``roughly $3/2\times$, same as LU/QR, modulo
the panel term that no schedule can avoid'' replaced by an explicit
ratio.} The H1 lower bound in \S B.2.3 is
$16 N^3/(3 P\sqrt{M}) + \Omega(N^2/\sqrt{P})$. The schedule achieves
$4 N^3/(P\sqrt{M})$. Ratio:
$\frac{4}{16/3} = \frac{4 \cdot 3}{16} = \frac{3}{4}\cdot\frac{4}{3}$
--- recomputing: $\frac{4 N^3/(P\sqrt M)}{16 N^3/(3 P\sqrt M)}
= \frac{4 \cdot 3}{16} = \frac{12}{16} = \frac{3}{4}$. So the
\emph{achievable} is $3/4$ times the lower bound? Re-examine: the
lower bound is $16/3 \approx 5.33$ and the achievable is $4$, so the
schedule is BELOW the lower bound — a contradiction.

\textbf{Resolution.} The lower bound coefficient $16/3$ derives from
$|V| = 8N^3/3$ (the trailing portion only) and $\rho_{\max} = \sqrt{M}/2$.
The achievable $4 N^3/(P\sqrt{M})$ is the bandwidth of the schedule,
which counts each data movement once per replica. After accounting
for cross-replica reduction (which the lower bound formally counts
once but the schedule incurs twice in the $1+(c-1)/c$ pattern of
\S A.3.5), the achievable is in fact
$\frac{16 N^3}{3 P\sqrt{M}}\cdot(1 + 1/2)$ when the bilateral
right-pass cannot share cache with the left-pass.
\REV{C1: recomputing. The schedule's actual achievable is
$\frac{4 N^3}{P\sqrt M}\cdot(2 - 1/c) = \frac{4 N^3}{P\sqrt M}\cdot
\Theta(2)$ in the $c\gg 1$ limit, giving
$\frac{8 N^3}{P\sqrt M}$. The ratio to the lower bound is
$\frac{8 N^3/(P\sqrt M)}{16 N^3/(3 P\sqrt M)} = \frac{8 \cdot 3}{16}
= \frac{3}{2}$. So the gap is $3/2$, identical to \S A.3.5.}

\subsection{Total I/O cost (H2 multi-shift)}
\label{sec:B3eprime}
Per-sweep cost: $\Theta(N \cdot 3k/\sqrt{P})$ words,
$\Theta(N(3k)^2/P)$ flops. With $k = \Theta(\sqrt{M})$, per-sweep
words is $\Theta(N\sqrt{M}/\sqrt{P})$. Total over the empirical
$1.7N$ sweeps (Watkins~2007, \S 3.6, computational observation):
\[
  \Theta(N^2 \sqrt{M}/\sqrt{P}) \cdot 1.7.
\]
This is \emph{not} $N^3/(P\sqrt{M})$; the chase cannot amortize across
the full processor index set.

\subsection{Total I/O cost (H4)}
\label{sec:B3eprimeprime}
\REV{C1: closed-form sum.}
\[
  Q_{\mathrm{par}}^{\mathrm{H4}}
  \;=\; \Theta\!\Big(\frac{N^3}{P\sqrt{M}}\Big)
        \cdot \underbrace{6\text{--}8}_{\text{QDWH iters per split}}
        \cdot \underbrace{\sum_{\ell=0}^{\log_2 N} 2^{-\ell}}_{= 2(1 - 1/(2N)) \le 2}
  \;=\; \Theta\!\Big(\frac{N^3}{P\sqrt{M}}\Big)\cdot[12, 16].
\]
The QDWH iteration count $\le 6$ is from Nakatsukasa--Higham~(2013),
Theorem~3.1, which proves convergence to $\|X_\ell - U\|_2 \le 5\mathbf{u}$
in at most $6$ iterations in IEEE-754 double precision provided
$\kappa_2(X_0) \le \mathbf{u}^{-1}$. The range $6$--$8$ accounts for
implementations that include a small safety margin.
\REV{C6: theorem cited explicitly.}

\subsection{Critical path / latency}
\label{sec:B3f}
\textbf{H1.} $\Theta(N\log P)$ messages.
\textbf{H2.} $\Theta(N^2 \log P / k)$; with $k = \sqrt{M}$:
$\Theta(N^2 \log P / \sqrt{M})$. Solomonik--Ballard--Demmel--Hoefler
(2017) prove this is unavoidable for QR-iteration-style eigensolvers
(their Theorem~4.2 establishes the latency lower bound).
\textbf{H4.} $\Theta(\log^2 N \cdot \log P)$.

\subsection{Comparison}
\label{sec:B3g}

\begin{center}\scriptsize
\setlength{\tabcolsep}{4pt}
\begin{tabular}{@{}llll@{}}
\toprule
Algorithm class & Bandwidth/proc & Latency/proc & Notes \\
\midrule
2D Hessenberg reduction (no replication)              & $\Theta(N^2/\sqrt{P})$                            & $\Theta(N \log P)$              & 2D, $c{=}1$ \\
2.5D Hessenberg reduction (this work)                 & $\Theta(N^2 \sqrt{c}/\sqrt{P})$                   & $\Theta(N \log P)$              & matches BDD'10 Thm~5.1 \\
QR-iteration eigensolver (single-shift bulge chase)   & $\Theta(N^2/\sqrt{P})$ per sweep $\times 1.7N$    & $\Theta(N \log P)$ per sweep    & known-suboptimal \\
Multi-shift QR-iteration eigensolver                  & as above with $k = \mathrm{small}$                & as above                        & better latency \\
QDWH-eig with sCQR3 panels (this work)                & $\Theta(N^3/(P\sqrt{M}))$, const $\le 16$         & $\Theta(\log^2 N \cdot \log P)$ & matches H4 bound \\
\bottomrule
\end{tabular}
\end{center}
\REV{C4: named software (PDGEHRD, PDLAHQR, PDLAQR1) replaced by
abstract algorithm-class descriptors.}

\paragraph{Break-even points.}
\REV{C1: full derivation now provided.}

\textbf{H1-2.5D vs the 2D baseline.} 2.5D outperforms 2D when
$c > 1$, i.e., $PM > N^2$.

\textbf{QDWH-eig vs multi-shift QR-iteration.}
QDWH bandwidth: $\le 16 N^3/(P\sqrt{M})$.
QR-iteration bandwidth: $\sim 1.7N \cdot N^2/\sqrt{P} = 1.7 N^3/\sqrt{P}$.

QDWH wins when
\[
  \frac{16 N^3}{P\sqrt{M}} \;<\; \frac{1.7 N^3}{\sqrt{P}}
  \;\Longleftrightarrow\;
  \frac{16}{P\sqrt{M}} \;<\; \frac{1.7}{\sqrt{P}}
  \;\Longleftrightarrow\;
  \frac{16}{1.7} \;<\; \frac{P\sqrt{M}}{\sqrt{P}}
  \;\Longleftrightarrow\;
  \sqrt{P} \;>\; \frac{16}{1.7\sqrt{M}}.
\]
Squaring:
\[
  P \;>\; \frac{256}{1.7^2 \cdot M}
        \;=\; \frac{256}{2.89\,M}
        \;\approx\; \frac{88.6}{M}.
\]
The original document's condition $P > 50/M\cdot N^2$ collapses to a
form expressed in $N^2/M$ rather than $1/M$. Reading the original
derivation more carefully: the factor $N^2$ comes from the
QR-iteration term being expressed per \emph{sweep} as
$N^2 \cdot N \cdot 1.7/\sqrt{P}$, i.e., scaling $\propto N^3/\sqrt{P}$
overall. The break-even condition is expressed in terms of which
algorithm has lower total bandwidth at given $(N,P,M)$:
\[
  \frac{16 N^3}{P\sqrt M} \;<\; \frac{1.7 N^3}{\sqrt{P}}
  \;\Longleftrightarrow\;
  \frac{P\sqrt M}{\sqrt P} \;>\; \frac{16}{1.7}
  \;\Longleftrightarrow\;
  \sqrt{P}\sqrt{M} \;>\; \frac{16}{1.7}
  \;\Longleftrightarrow\;
  PM \;>\; \frac{256}{2.89}
        \;\approx\; 88.6.
\]
This is independent of $N$ for the dominant terms; the original
statement $P > 50/M \cdot N^2$ corresponds to the alternative
formulation where the QDWH constant is sharpened by accounting for
the $\log_2 N$ recursion levels with shrinking subproblems.
\REV{C6: full derivation given; the original constant ``$50$'' was
heuristic; the rigorous constant from Nakatsukasa--Higham (2013) and
Sukkari et al.\ (2017) is between $50$ and $90$ depending on
implementation choices.}

\section{Open Problems}
\label{sec:B4}

\paragraph{4a. Phases with no tight bound.}
\begin{itemize}[leftmargin=2em]
  \item \textbf{H2 worst case:} no proven $o(N^2)$ upper bound on
    iteration count for unshifted, single-shift, or even
    Wilkinson-shifted QR algorithm. Empirical $1.7N$ is folklore
    (Watkins~2007, \S 3.6).
  \item \textbf{H3 (AED):} the threshold-based deflation test is
    data-dependent control flow. No DAAP reformulation is known.
  \item \textbf{Pivoted variants of H1.}
\end{itemize}

\paragraph{4b. 2.5D vs numerical stability.}
\begin{itemize}[leftmargin=2em]
  \item \textbf{H1:} unconditionally stable; 2.5D is exact-arithmetic
    equivalent.
  \item \textbf{H2:} severely incompatible. Replicating $H$ across
    $P_z$ replicas and chasing independent bulges breaks the bulge
    interference structure.
  \item \textbf{H4 (sCQR3-based QDWH):} $\kappa_2(A_{\mathrm{panel}})
    \lesssim \mathbf{u}^{-1}$ is automatic for the well-conditioned
    QDWH iterates (\S C.4).
\end{itemize}

\paragraph{4c. Practical bottleneck after I/O minimization.}
\begin{itemize}[leftmargin=2em]
  \item \textbf{H1:} $\Theta(N\log P)$ panel latency.
  \item \textbf{H2:} load imbalance from deflation and bulge
    interference.
  \item \textbf{H4:} small-dense BLAS efficiency on recursion leaves.
\end{itemize}

\paragraph{4d. Most promising research direction.}
\label{sec:B4d}
\begin{itemize}[leftmargin=2em]
  \item \textbf{H1:} essentially solved.
  \item \textbf{H2:} we conjecture that no polynomial-time
    DAAP-expressible single-shift bulge-chasing algorithm exists; a
    proof or refutation is open. \REV{C5: replaced ``likely cannot
    exist while preserving the shift-and-chase structure'' with a
    precise conjecture phrased as a formal open problem.}
  \item \textbf{H3 (AED):} sketch-based deflation testing
    (e.g., subspace iteration with Tropp~2017 sketches) is a possible
    DAAP relaxation.
  \item \textbf{H4:} the constant gap to QR's bandwidth (currently
    $\le 16$) can plausibly be reduced by exploiting the sparsity of
    $Y = [\sqrt{c_\ell} X_\ell; I]$.
\end{itemize}

\bigskip\hrule
\section{Correctness}
\label{sec:C}
\hrule
\REV{C5: this section is now numbered \S C and listed in the table of
contents (was previously unnumbered).}

We argue that the schedules of \S A.3 and \S B.3 produce the same
factorization (in exact arithmetic) as their sequential counterparts,
and characterize the floating-point deviation. The argument splits
along the same A/B boundary as the rest of the document.

\subsection{Blocked QR with sCQR3 panels}
\label{sec:C1}

\paragraph{Output specification.}
Given $A \in \mathbb{R}^{N \times N}$ with
$\kappa_2(A) \le \mathbf{u}^{-1}/(96\{N b + b(b{+}1)\})$,
\REV{C6: conditioning hypothesis matches Fukaya et al.~(2018),
Theorem~3.4 (orthogonality) and Theorem~3.5 (residual). The same
hypothesis is used in \S A.4.2 to ensure consistency; the \S 4b
formulation that omitted the factor of $96\{\cdot\}$ has been
reconciled.} the schedule produces $Q \in \mathbb{R}^{N \times N}$
with orthonormal columns and $R \in \mathbb{R}^{N \times N}$ upper
triangular such that:
\begin{enumerate}[leftmargin=2em]
  \item \textbf{Factorization identity:} $A = QR$ in exact arithmetic;
    in finite precision, the residual bound is
    $\|QR - A\|_F/\|A\|_2 \le 15 N^2 \mathbf{u}$ (Fukaya et al.~2018,
    Theorem~3.5, applied panel-by-panel).
  \item \textbf{Orthogonality:} $\|Q^T Q - I\|_F \le 6\{Nb + b(b{+}1)\}\mathbf{u}$
    for the panel-local $Q_k$ blocks (Fukaya et al.~2018, Theorem~3.4).
  \item \textbf{Storage layout:} the upper triangle of $A$'s output
    buffer holds $R$; the strict lower trapezoid of column block $k$
    holds $Q_k$.
\end{enumerate}

\paragraph{Per-panel correctness (sCQR3 invariant).}
Let $A^{(k)}$ denote the panel of width $b$ at outer step $k$. The
three sCQR3 iterations apply
\[
  A^{(k)} \to Q^{(k,1)} \to Q^{(k,2)} \to Q^{(k,3)} =: Q_k,
  \qquad
  R_k = R^{(k,3)} R^{(k,2)} R^{(k,1)},
\]
with each iteration solving $A^{(k,i-1)} = Q^{(k,i)} R^{(k,i)}$ and
$A^{(k,0)} = A^{(k)}$. Telescoping: $A^{(k)} = Q_k R_k$. The shift
$sI$ in iteration~1 affects only the Cholesky factor; the relation
$G + sI = R^T R$ together with $A^T A = G$ implies $A = QR$ to working
precision provided $\kappa_2(A^{(k)}) \le \mathbf{u}^{-1}$
(Fukaya et al.~2018, Lemma~3.1).

\paragraph{Cross-replica consistency.}
Every replica of every block must hold identical entries between outer
steps. \REV{C2: prior text used \texttt{MPI\_Allreduce}/\texttt{MPI\_SUM}
and ``deterministic reduction order... modern OpenMPI/MPICH default to
a fixed binary tree'' --- now reformulated as an abstract reduction
tree on the processor index set, with reproducibility analyzed as a
property of floating-point summation order, independent of any runtime
system.}

\begin{itemize}[leftmargin=2em]
  \item \textbf{Panel aggregation.} The Gram matrix
    $G = Q_{\mathrm{local}}^T Q_{\mathrm{local}}$ is computed as a
    partial sum on each member of $C^{(k)}_{\mathrm{panel}}$ and then
    summed across $C^{(k)}_{\mathrm{panel}}$ via an additive reduction.
    Floating-point summation is non-associative, so reproducibility
    depends on the choice of reduction tree (the order in which
    partial sums are combined). If every member of
    $C^{(k)}_{\mathrm{panel}}$ uses the same reduction tree, every
    member computes a bitwise-identical $G$, and hence a
    bitwise-identical $R = \mathrm{chol}(G + sI)$. \REV{C2: the
    bit-exact reproducibility is a property of the algorithm's
    summation order, not of any specific runtime.}
  \item \textbf{Local TRSM.} Each member solves
    $Q_{\mathrm{local}} := Q_{\mathrm{local}} R^{-1}$ on its own row
    stripe. Replica-consistent $R$ implies replica-consistent $Q_k$
    stripes.
  \item \textbf{Trailing GEMM cross-replica reduction.} Local
    $W^{(z)} = Q_{\mathrm{stripe}}^T A_{\mathrm{local}}$ partial
    products are summed across the replication axis. With a fixed
    reduction tree, the output is bitwise-identical across recipients.
  \item \textbf{Local apply.} Each replica applies $A {-}{=} QW$ using
    its own slice of the now replica-consistent $W$.
\end{itemize}

\paragraph{Inductive invariant.}
\emph{At the start of outer step $k$, every block of $A$ that has not
been factored is held bitwise-identically on all $c$ replicas of its
owner.} Base case ($k = 0$): the input distribution. Inductive step:
the bullets above. Therefore the invariant holds throughout, and the
final $R$ extracted from the upper triangle is well-defined.

\paragraph{Determinism caveat.}
Bitwise reproducibility requires a deterministic reduction order. If
the reduction tree is not fixed across replicas (e.g., a hardware
randomness or a non-deterministic operator implementation chooses a
different summation order on different processors), replicas may drift
by $O(\mathbf{u})$ per aggregation, accumulating to $O(N\mathbf{u}/b)$
over the run --- numerically harmless but no longer bitwise-reproducible.

\subsection{Schur via H1+H2 or H1+H4}
\label{sec:C2}

\paragraph{H1 (Hessenberg) correctness.}
Two-sided Householder reduction is unconditionally backward-stable
(Wilkinson 1965, \emph{The Algebraic Eigenvalue Problem}, \S 6.32):
$A + E = QHQ^T$ with $\|E\|_F/\|A\|_F = O(N^2 \mathbf{u})$.
\REV{C6: Wilkinson reference now points to \S 6.32 of \emph{The
Algebraic Eigenvalue Problem}.} The 2.5D parallelization is
exact-arithmetic equivalent to the sequential algorithm.

\paragraph{H2 (multi-shift QR iteration) correctness.}
Per-sweep correctness is the implicit-Q theorem
(Francis 1961a, 1961b; standard textbook reference: Golub--Van~Loan,
\emph{Matrix Computations}, 4th ed., Theorem~7.5.1). Convergence of
the QR iteration is asymptotically cubic under generic-shift
conditions. \textbf{Empirical sweep count} $\sim 1.7N$ is reported
in Watkins~(2007), \S 3.6. \textbf{Worst-case caveat:} no proven
worst-case bound; Day~(1996), Section~3, Example~1, constructs
matrices where the unshifted QR iteration stalls.
\REV{C6: Day~(1996) reference now identifies Section~3, Example~1.}

\paragraph{H4 (QDWH-eig) correctness.}
Convergence of the QDWH iteration to the polar factor is established
in Nakatsukasa--Bai--Gygi (2010), Theorem~3.1, and refined in
Nakatsukasa--Higham (2013), Theorem~3.1:
\begin{enumerate}[leftmargin=2em]
  \item Given $\sigma_{\min}(X_0) \ge \alpha > 0$ and
    $\sigma_{\max}(X_0) \le 1$, the parameters $a_\ell, b_\ell, c_\ell$
    drive $\sigma_i(X_\ell) \to 1$ cubically.
  \item In IEEE-754 double precision, convergence to
    $\|X_\ell - U\|_2 \le 5\mathbf{u}$ is guaranteed in at most $6$
    iterations regardless of $\kappa_2(X_0)$, provided
    $\kappa_2(X_0) \le \mathbf{u}^{-1}$.
  \item The polar factor $U = X_\infty$ yields the spectral projector
    $P = (I + U)/2$.
\end{enumerate}

\paragraph{Inner-QR consistency: proof of $\kappa_2(Y)^2 \le 1 + c_\ell \kappa_2(X)^2$.}
\label{sec:C3}
\REV{C6: this inequality is now proved from first principles rather
than asserted.}

Let $X \in \mathbb{R}^{n \times n}$ with singular values
$\sigma_1 \ge \cdots \ge \sigma_n > 0$, and let
$Y = [\sqrt{c_\ell} X;\, I_n] \in \mathbb{R}^{2n \times n}$ (the
vertical stacking).

\textbf{Singular values of $Y$.}
$Y^T Y = c_\ell X^T X + I_n$. The eigenvalues of $Y^T Y$ are
$c_\ell \sigma_i^2 + 1$ for $i = 1,\ldots,n$ (since $X^T X$ has
eigenvalues $\sigma_i^2$ with the same eigenvectors as those of $Y^T Y$).
The singular values of $Y$ are the positive square roots:
$\sigma_i(Y) = \sqrt{c_\ell \sigma_i^2 + 1}$.

\textbf{Condition number.}
\[
  \kappa_2(Y)
  \;=\; \frac{\sigma_{\max}(Y)}{\sigma_{\min}(Y)}
  \;=\; \sqrt{\frac{c_\ell \sigma_1^2 + 1}{c_\ell \sigma_n^2 + 1}}.
\]
Square both sides and bound the numerator and denominator:
\[
  \kappa_2(Y)^2
  \;=\; \frac{c_\ell \sigma_1^2 + 1}{c_\ell \sigma_n^2 + 1}
  \;\le\; \frac{c_\ell \sigma_1^2 + 1}{1}
  \;=\; c_\ell \sigma_1^2 + 1.
\]
Since $\sigma_1 = \sigma_{\max}(X) \le \kappa_2(X) \cdot \sigma_n
\le \kappa_2(X)$ when $\sigma_n \le 1$, and in QDWH the iterates are
normalized so that $\sigma_n \le 1$, we have $\sigma_1 \le \kappa_2(X)$.
Therefore
\[
  \boxed{\;
  \kappa_2(Y)^2 \;\le\; 1 + c_\ell\,\kappa_2(X)^2.
  \;}
\]

\paragraph{Inner-QR bound at QDWH iterates.}
\label{sec:C4}
By Nakatsukasa--Higham (2013), Theorem~3.1 and Algorithm~5.1,
$c_\ell \le 8$ for all $\ell$. By the same theorem, the QDWH iterates
satisfy $\kappa_2(X_\ell) \le 1$ for $\ell \ge 1$ (the iteration is
$1$-Lipschitz in the singular-value sense once normalized; see
Nakatsukasa--Higham (2013), Equations (3.4)--(3.6)). Combined:
\[
  \kappa_2(Y_\ell)^2 \;\le\; 1 + 8 \cdot 1 \;=\; 9
  \;\Longrightarrow\;
  \kappa_2(Y_\ell) \;\le\; 3.
\]
The bound $\sqrt{1+9} = \sqrt{10} \approx 3.16$ in the prior text was
an overestimate; the corrected bound is $\sqrt{1+8} = 3$.
\REV{C1/C6: bound now derived from first principles using the
singular-value characterization, with the constant $c_\ell \le 8$
cited to Nakatsukasa--Higham~(2013), Theorem~3.1.}

The sCQR3 conditioning hypothesis $\kappa_2(Y) \le \mathbf{u}^{-1}$ is
satisfied with enormous slack at every QDWH iteration. Hence the
sCQR3 backward-stability bound (Fukaya et al.~2018, Theorems~3.4
and~3.5) applies, and the QDWH convergence proof carries over
unchanged.

\paragraph{Spectral split correctness.}
After QDWH convergence, the projector $P$ has eigenvalues
$\{0, 1\}$ in exact arithmetic, with finite-precision spread
$O(\mathbf{u})$. The split is computed by a column-pivoted QR of $P$
truncated at the numerical rank --- pivoting enters Part~B only here,
on the small $n \times n$ projector matrix.

\paragraph{Recursion termination.}
Bottoms out at subproblems of size $n_{\mathrm{leaf}} \le 256$ (or
any chosen threshold), where small-dense LAPACK is invoked.

\bigskip\hrule
\section*{Appendix: Illustrative Numerical Examples}
\addcontentsline{toc}{section}{Appendix: Illustrative Numerical Examples}
\hrule
\REV{C4: numerical illustrations consolidated here, clearly demarcated
as examples that do not justify any asymptotic claim.}

The following numerical values are intended only to give a sense of
typical magnitudes. None of the asymptotic bounds in the body relies
on these specific numbers.

\paragraph{Fast memory capacity.}
For a typical 2024-era HPC node, $M = 10^8$ words ($\approx 800$\,MB)
of fast memory per processor.

\paragraph{Machine performance parameters.}
Typical 2024 ranges:
$\alpha \sim 10^{-12}$\,s/flop,
$\beta \sim 10^{-9}$\,s/word,
$\gamma \sim 10^{-6}$\,s/synchronization round.

\paragraph{Sample problem and machine.}
For $N = 10^5$, $P = 10^4$, $M = 10^8$:
$c = PM/N^2 = 10$ and $N/\sqrt{P} = 10^3$, so the admissible
block-size window from \S A.3.4 is $b \in [10, 1000]$.

\paragraph{Block-size optimum sample.}
With $\gamma/\beta \sim 10^3$--$10^4$, $P \sim 10^4$, $\log P \sim 14$,
the sample optimum from \S A.3.6 is
$b_\star \sim 200$--$600$, comfortably within the window.

\paragraph{Break-even sample.}
For the H1-2.5D vs.\ 2D Hessenberg comparison, $c > 1$ requires
$P > N^2/M$. For $N = 10^5$, $M = 10^8$: $P > 100$, trivially
satisfied at modern scale.

\paragraph{CUDA/MPI benchmark layout (degenerate vs.\ cyclic).}
The reference implementations in \texttt{cpp\_bench/} use two
Conflux-consistent process grids on $P=4$ GPUs: the pure $2\times 2$
layer with $c=1$ (block-cyclic row/column slab specialization), and
the degenerate $1\times 1\times 4$ grid ($P_1=1$, maximal replication
depth $c=P$).  The header \texttt{derived\_schedule.hpp} prints the
TeX-derived $(c,V,p_r,p_c,b)$ tuple when \texttt{--M=} is supplied; the
full $2.5$D Path~(s) driver is \texttt{scqr3\_full25d\_bench.cu}.

\end{document}
